% Lectura y comentario : contaminación y medidas de preservación ; creación de carteles de sensibilización — Espagnol Tle A — Cours signé OnBuch+
% Programme officiel MINESEC. Compiler avec : tectonic E09-contaminacion-carteles-sensibilizacion.tex
\newcommand{\DOCMATIERE}{Espagnol}
\newcommand{\DOCNIVEAU}{Tle A}
\newcommand{\DOCMODULE}{Módulo 2 : Environnement, santé et bien-être}
\newcommand{\DOCLECON}{E09}
\newcommand{\DOCTITRE}{Lectura y comentario : contaminación y medidas de preservación ; creación de carteles de sensibilización}
% OnBuch+ — préambule « cours signés » ENRICHI (compilé avec tectonic / XeLaTeX)
% Système visuel « Pop » d'OnBuch+ : fond crème (jamais blanc), bordures encre
% épaisses, ombres « dures » décalées, titres Archivo Black en capitales.
% Version MATHÉMATIQUES : graphes (pgfplots/tikz), boîtes pédagogiques
% (définition, propriété, méthode, attention, le savais-tu, exercice résolu,
% exercice, TP, bilan), page de garde et signature OnBuch+ ancrée en bas.
%
% Macros attendues dans main.tex AVANT \input{preamble} :
%   \DOCMATIERE  \DOCNIVEAU  \DOCMODULE  \DOCLECON (numéro)  \DOCTITRE
\documentclass[11pt]{article}

\usepackage{fontspec}
\usepackage[spanish,french]{babel} % français = langue principale ; l'espagnol via \esp{..} / espagnol
\usepackage[a4paper,margin=2cm,top=3.4cm,bottom=2.8cm,headheight=1.4cm,footskip=1.3cm]{geometry}
\usepackage{amsmath,amssymb,amsfonts}
\usepackage{mathtools} % \xmapsto, \coloneqq... souvent utilisés par les modèles
\usepackage{cancel} % simplifications barrées (bilans de forces)
% \abs{z} (module, valeur absolue) et \norm{u} (norme), utilisés par les modèles
\providecommand{\abs}{}\let\abs\relax\DeclarePairedDelimiter{\abs}{\lvert}{\rvert}
\providecommand{\norm}{}\let\norm\relax\DeclarePairedDelimiter{\norm}{\lVert}{\rVert}
\usepackage[table]{xcolor} % [table] : fournit aussi \rowcolors (alternance de lignes)
\usepackage{pagecolor}
\usepackage{tikz}
\usetikzlibrary{arrows.meta,positioning,calc,shapes.geometric,shapes.misc,shapes.arrows,decorations.pathmorphing,decorations.pathreplacing,decorations.markings,patterns,shadows,mindmap,trees,fit,backgrounds,matrix,intersections,angles,quotes}
\usepackage{pgfplots}
\usepackage[version=4]{mhchem} % équations nucléaires \ce{^{235}_{92}U}
\usepackage[european,siunitx]{circuitikz} % schémas électriques
\pgfplotsset{compat=1.18}
\usepgfplotslibrary{fillbetween,groupplots} % aires entre courbes (calcul intégral)
\usepackage{enumitem}
\usepackage{fancyhdr}
% [most] charge toutes les bibliothèques tcolorbox (listings, minted, raster,
% documentation...) et gonfle inutilement la mémoire TeX sur un document long
% (~300 boîtes) : on ne charge que ce qui est réellement utilisé.
\usepackage[skins,breakable]{tcolorbox}
\usepackage{graphicx}
\usepackage{colortbl}
\usepackage{booktabs}
\usepackage{tabularx}
\usepackage{diagbox} % case d'en-tête diagonale des tableaux à double entrée
% Colonnes extensibles centrée / gauche / droite (C, L, R), souvent utilisées par les modèles
\newcolumntype{C}{>{\centering\arraybackslash}X}
\newcolumntype{L}{>{\raggedright\arraybackslash}X}
\newcolumntype{R}{>{\raggedleft\arraybackslash}X}
\usepackage{array}
\usepackage{multicol}
\usepackage{siunitx}
% parse-numbers=false : accepte aussi \SI{2,5 \times 10^{-3}}{...}
\sisetup{locale=FR,output-decimal-marker={,},group-minimum-digits=4,parse-numbers=false}
\DeclareSIUnit\ohm{\ensuremath{\symup{\Omega}}} % U+2126 absent de la police
\DeclareSIUnit\division{div} % réglages d'oscilloscope (ms/div, V/div)
\DeclareSIUnit\tr{tr} % tours (vitesse de rotation, ex. \SI{1500}{\tr\per\minute})
\DeclareSIUnit\molar{M} % concentration molaire (ex. \SI{3}{\molar}, \SI{1}{\micro\molar})
\DeclareSIUnit\an{an} % année (le modèle confond parfois avec \year, un compteur TeX) — ex. \SI{5}{\kilo\meter\per\an}
\DeclareSIUnit\FCFA{FCFA} % franc CFA (ex. \SI{500}{\FCFA\per\kilo\gram})
\DeclareSIUnit\degreeBrix{\textdegree Brix} % teneur en sucre d'un sirop/jus (réfractométrie)
\DeclareSIUnit\month{mois} % le modèle utilise parfois \month, un compteur TeX, comme unité de durée
\usepackage{qrcode}
% \degree : commande fréquemment utilisée par les modèles pour un angle ou une
% température en dehors de \SI{}{} (non fournie par les paquets chargés).
\providecommand{\degree}{^{\circ}}
% \qed (fin de démonstration), utilisé par les modèles sans amsthm
\providecommand{\qed}{\hfill$\square$}
% \barre : double barre des valeurs interdites dans un tableau de variations (tkz-tab)
\providecommand{\barre}{\ensuremath{\|}}
% environnement proof (amsthm non chargé) utilisé par les modèles
\ifdefined\proof\else\newenvironment{proof}[1][Démonstration]{\par\noindent\textit{#1.}\ }{\qed\par}\fi
\usepackage{lastpage}
\usepackage{hyperref}
\hypersetup{hidelinks}

% ============================================================
% Palette « Pop » OnBuch+ — mêmes hex que la classe P (plus_ui.dart)
% ============================================================
\definecolor{poporange}{HTML}{F59321}
\definecolor{popdark}{HTML}{E07A0C}
\definecolor{popcream}{HTML}{FFF6E8}
\definecolor{popink}{HTML}{1C1714}
\definecolor{poppurple}{HTML}{7A5AE0}
\definecolor{popgreen}{HTML}{2E9E6A}
\definecolor{poppink}{HTML}{E0457B}
\definecolor{popblue}{HTML}{2F7DE1}
\definecolor{popgold}{HTML}{C98A12}
\definecolor{popmuted}{HTML}{8A7F76}
\definecolor{popline}{HTML}{EADFD2}
\definecolor{popgray}{HTML}{8A7F76}
\colorlet{popgrey}{popgray}
% Alias tolérants (le modèle nomme parfois les couleurs différemment)
\colorlet{popwhite}{white}
\colorlet{popblack}{popink}
\colorlet{popred}{poppink}
\colorlet{popyellow}{popgold}
% Teintes claires (fonds de tableaux/encadrés) — le modèle nomme parfois
% les couleurs avec un suffixe L (ex. popblueL) sans les avoir définies.
\colorlet{poporangeL}{poporange!20}
\colorlet{popdarkL}{popdark!20}
\colorlet{poppurpleL}{poppurple!20}
\colorlet{popgreenL}{popgreen!20}
\colorlet{poppinkL}{poppink!20}
\colorlet{popblueL}{popblue!20}
\colorlet{popgoldL}{popgold!20}
\colorlet{popredL}{popred!20}
\colorlet{popgrayL}{popgray!20}
% Teintes claires pour fonds de tableaux / zones
\colorlet{popblueL}{popblue!10!white}
\colorlet{poporangeL}{poporange!14!white}
\colorlet{popgreenL}{popgreen!12!white}
\colorlet{poppurpleL}{poppurple!12!white}
\colorlet{poppinkL}{poppink!12!white}
\colorlet{popgoldL}{popgold!14!white}

\pagecolor{popcream}

% ============================================================
% Typographie
% ============================================================
\defaultfontfeatures{Ligatures=TeX}
\setmainfont{PlusJakartaSans-Medium.ttf}[Path=../../../fonts/,BoldFont=PlusJakartaSans-Bold.ttf]
% Maths en sans-serif (Fira Math) avec chiffres et lettres droites de Plus
% Jakarta, pour que les formules se fondent dans le texte.
\usepackage[math-style=ISO,bold-style=ISO]{unicode-math}
\setmathfont{FiraMath-Regular.otf}[Path=../../../fonts/]
\setmathfont{PlusJakartaSans-Medium.ttf}[Path=../../../fonts/,range={up/num,up/{Latin,latin},\mathup}]
\usepackage{icomma} % virgule décimale sans espace en mode math
% Caractères mathématiques Unicode absents des polices de texte (Plus Jakarta,
% Archivo Black) : rendus par leur équivalent mathématique, en texte comme en
% formule — sinon ils s'affichent en carré vide (ex. « Arithmétique dans ℤ »).
\usepackage{newunicodechar}
\newunicodechar{ℝ}{\ensuremath{\mathbb{R}}}
\newunicodechar{ℤ}{\ensuremath{\mathbb{Z}}}
\newunicodechar{ℂ}{\ensuremath{\mathbb{C}}}
\newunicodechar{ℕ}{\ensuremath{\mathbb{N}}}
\newunicodechar{ℚ}{\ensuremath{\mathbb{Q}}}
\newunicodechar{𝔻}{\ensuremath{\mathbb{D}}}
\newunicodechar{α}{\ensuremath{\alpha}}
\newunicodechar{β}{\ensuremath{\beta}}
\newunicodechar{γ}{\ensuremath{\gamma}}
\newunicodechar{δ}{\ensuremath{\delta}}
\newunicodechar{ε}{\ensuremath{\varepsilon}}
\newunicodechar{θ}{\ensuremath{\theta}}
\newunicodechar{λ}{\ensuremath{\lambda}}
\newunicodechar{μ}{\ensuremath{\mu}}
\newunicodechar{σ}{\ensuremath{\sigma}}
\newunicodechar{τ}{\ensuremath{\tau}}
\newunicodechar{φ}{\ensuremath{\varphi}}
\newunicodechar{ψ}{\ensuremath{\psi}}
\newunicodechar{ω}{\ensuremath{\omega}}
\newunicodechar{Σ}{\ensuremath{\Sigma}}
% majuscules produites par \MakeUppercase dans les titres (α-aminés -> Α)
\newunicodechar{Α}{\ensuremath{\alpha}}
\newunicodechar{Β}{\ensuremath{\beta}}
\newunicodechar{Γ}{\ensuremath{\Gamma}}
\newunicodechar{Δ}{\ensuremath{\Delta}}
\newunicodechar{Ε}{\ensuremath{\varepsilon}}
\newunicodechar{Θ}{\ensuremath{\Theta}}
\newunicodechar{Λ}{\ensuremath{\Lambda}}
\newunicodechar{Μ}{\ensuremath{\mu}}
\newunicodechar{Τ}{\ensuremath{\tau}}
\newunicodechar{Φ}{\ensuremath{\Phi}}
\newunicodechar{Ψ}{\ensuremath{\Psi}}
\newunicodechar{Ω}{\ensuremath{\Omega}}
\newunicodechar{↦}{\ensuremath{\mapsto}}
\newunicodechar{∈}{\ensuremath{\in}}
\newunicodechar{∉}{\ensuremath{\notin}}
\newunicodechar{∧}{\ensuremath{\wedge}}
\newunicodechar{⇔}{\ensuremath{\Leftrightarrow}}
\newunicodechar{⟺}{\ensuremath{\Longleftrightarrow}}
\newunicodechar{⇒}{\ensuremath{\Rightarrow}}
\newunicodechar{⟹}{\ensuremath{\Longrightarrow}}
\newunicodechar{∘}{\ensuremath{\circ}}
\newunicodechar{∪}{\ensuremath{\cup}}
\newunicodechar{∩}{\ensuremath{\cap}}
\newunicodechar{≡}{\ensuremath{\equiv}}
\newunicodechar{⊂}{\ensuremath{\subset}}
\newunicodechar{⊥}{\ensuremath{\perp}}
\newunicodechar{∃}{\ensuremath{\exists}}
\newunicodechar{∀}{\ensuremath{\forall}}
\newunicodechar{∛}{\ensuremath{\sqrt[3]{\,}}}
\newunicodechar{∜}{\ensuremath{\sqrt[4]{\,}}}
\newunicodechar{⃗}{\ensuremath{{}^{\to}}}
\newunicodechar{ⁿ}{\ifmmode^{n}\else\textsuperscript{\ensuremath{n}}\fi}
\newunicodechar{ˣ}{\ifmmode^{x}\else\textsuperscript{\ensuremath{x}}\fi}
\newunicodechar{ᵇ}{\ifmmode^{b}\else\textsuperscript{\ensuremath{b}}\fi}
\newunicodechar{ʳ}{\ifmmode^{r}\else\textsuperscript{\ensuremath{r}}\fi}
\newunicodechar{ᵘ}{\ifmmode^{u}\else\textsuperscript{\ensuremath{u}}\fi}
\newunicodechar{ʸ}{\ifmmode^{y}\else\textsuperscript{\ensuremath{y}}\fi}
\newunicodechar{ᵃ}{\ifmmode^{a}\else\textsuperscript{\ensuremath{a}}\fi}
\newunicodechar{ᵉ}{\ifmmode^{e}\else\textsuperscript{\ensuremath{e}}\fi}
\newunicodechar{ᵏ}{\ifmmode^{k}\else\textsuperscript{\ensuremath{k}}\fi}
\newunicodechar{ᵖ}{\ifmmode^{p}\else\textsuperscript{\ensuremath{p}}\fi}
\newunicodechar{ᴺ}{\ifmmode^{N}\else\textsuperscript{\ensuremath{N}}\fi}
\newunicodechar{ᶜ}{\ifmmode^{c}\else\textsuperscript{\ensuremath{c}}\fi}
\newunicodechar{ʰ}{\ifmmode^{h}\else\textsuperscript{\ensuremath{h}}\fi}
\newunicodechar{ᵠ}{\ifmmode^{q}\else\textsuperscript{\ensuremath{q}}\fi}
\newunicodechar{ₙ}{\ifmmode_{n}\else\textsubscript{\ensuremath{n}}\fi}
\newunicodechar{ₐ}{\ifmmode_{a}\else\textsubscript{\ensuremath{a}}\fi}
\newunicodechar{ᵢ}{\ifmmode_{i}\else\textsubscript{\ensuremath{i}}\fi}
\newunicodechar{ᵣ}{\ifmmode_{r}\else\textsubscript{\ensuremath{r}}\fi}
\newunicodechar{ₖ}{\ifmmode_{k}\else\textsubscript{\ensuremath{k}}\fi}
\newunicodechar{ᵦ}{\ifmmode_{\beta}\else\textsubscript{\ensuremath{\beta}}\fi}
\newunicodechar{ₓ}{\ifmmode_{x}\else\textsubscript{\ensuremath{x}}\fi}
\newfontfamily\popdisplay{ArchivoBlack-Regular.ttf}[Path=../../../fonts/]
\newfontfamily\poplogo{SpaceGrotesk-Bold.ttf}[Path=../../../fonts/]
\newfontfamily\popsemi{PlusJakartaSans-SemiBold.ttf}[Path=../../../fonts/]
\newfontfamily\popxbold{PlusJakartaSans-ExtraBold.ttf}[Path=../../../fonts/]
\newfontfamily\popxboldls{PlusJakartaSans-ExtraBold.ttf}[Path=../../../fonts/,LetterSpace=3.0]

\setlist[enumerate]{leftmargin=1.7em,itemsep=5pt,topsep=4pt}
\setlist[itemize]{leftmargin=1.7em,itemsep=4pt,topsep=4pt,label={\color{poporange}\textbullet}}
\setlength{\parindent}{0pt}
\setlength{\parskip}{5pt}
\renewcommand{\arraystretch}{1.25}

% pgfplots : style Pop par défaut
\pgfplotsset{
  every axis/.append style={axis line style={popink,line width=1pt},tick style={popink},
    grid style={popline,line width=0.5pt},label style={font=\small},tick label style={font=\footnotesize}},
  popcurve/.style={poporange,line width=1.8pt,no markers,smooth},
}
% Même style disponible dans un \draw TikZ simple (hors pgfplots)
\tikzset{popcurve/.style={poporange,line width=1.8pt}}
\tikzset{popgrid/.style={popline,very thin}}
\colorlet{popcurve}{poporange} % le nom du style est parfois utilisé comme couleur

% ============================================================
% Pill / squiggle
% ============================================================
\newcommand{\poppill}[3]{%
  \tikz[baseline=-0.55ex]\node[draw=popink,line width=1.2pt,rounded corners=2.6pt,%
    fill=#2,inner xsep=6.5pt,inner ysep=3pt]{%
    \popxboldls\fontsize{7}{7}\selectfont\color{#3}\MakeUppercase{#1}};%
}
\newcommand{\popsquiggle}[1]{%
  \tikz[baseline=-0.4ex]{
    \draw[#1,line width=2.6pt,line cap=round]
      plot [smooth] coordinates {(0,0.09) (0.34,0.01) (0.68,0.21) (1.02,0.01) (1.36,0.10)};
  }%
}
% Mot-clé surligné (marqueur orange sous le texte)
\newcommand{\cle}[1]{{\popxbold\color{popdark}#1}}

% ============================================================
% En-tête / pied
% ============================================================
\pagestyle{fancy}
\fancyhf{}
\renewcommand{\headrulewidth}{0pt}
\renewcommand{\footrulewidth}{0pt}
\lhead{\raisebox{-0.15ex}{\poplogo\fontsize{13}{13}\selectfont\color{popink}OnBuch\color{poporange}+}}
\rhead{\poppill{\DOCMATIERE\ \textbullet\ \DOCNIVEAU\ \textbullet\ Leçon \DOCLECON}{white}{popink}}
\lfoot{\popxbold\fontsize{7.6}{7.6}\selectfont\color{popmuted}\MakeUppercase{Cours signé OnBuch+}\ \textbullet\ {\popsemi\DOCTITRE}}
\rfoot{\tikz[baseline=-0.6ex]\node[rounded corners=8.5pt,fill=popink,minimum height=17pt,minimum width=17pt,inner xsep=5pt,inner sep=0]{\popxbold\fontsize{8}{8}\selectfont\color{popcream}\thepage\,/\,\pageref*{LastPage}};}

% ============================================================
% Titres
% ============================================================
\newcommand{\chaptitre}[2]{%
  \begin{tcolorbox}[enhanced,colback=poporange,colframe=popink,boxrule=2.6pt,arc=11pt,
      drop shadow=popink,left=15pt,right=15pt,top=12pt,bottom=11pt,before skip=3pt,after skip=18pt]
    \poppill{#2}{popink}{popcream}\par\vspace{9pt}
    {\raggedright\hyphenpenalty=10000\exhyphenpenalty=10000\popdisplay\fontsize{22}{24}\selectfont\color{popcream}\MakeUppercase{#1}\par}
  \end{tcolorbox}
}
% Section numérotée : pastille orange avec le numéro + titre Archivo
\newcounter{popsec}
\newcommand{\coursec}[1]{%
  \stepcounter{popsec}\setcounter{popsub}{0}%
  \par\vspace{16pt}\noindent
  \begin{minipage}{\linewidth}
  \tikz[baseline=-0.7ex]\node[draw=popink,line width=1.6pt,fill=poporange,rounded corners=4pt,minimum size=22pt,inner sep=0]{\popdisplay\fontsize{12}{12}\selectfont\color{popcream}\thepopsec};%
  \hspace{8pt}{\popdisplay\fontsize{15}{17}\selectfont\color{popink}\MakeUppercase{#1}}\par
  \vspace{4pt}\noindent{\color{poporange}\rule{36pt}{3.6pt}}
  \end{minipage}\par\nopagebreak\vspace{8pt}
}
\newcounter{popsub}
\newcommand{\courssub}[1]{%
  \stepcounter{popsub}%
  \par\vspace{9pt}\noindent{\popxbold\fontsize{11.5}{13}\selectfont\color{popink}{\color{poporange}\thepopsec.\thepopsub}\hspace{6pt}#1}\par\nopagebreak\vspace{4pt}
}

% ============================================================
% Boîtes pédagogiques — même recette : fond blanc, bordure épaisse colorée,
% ombre dure encre, badge plein. Titre optionnel [#1].
% ============================================================
\tcbset{popbase/.style={breakable,enhanced,colback=white,boxrule=2.3pt,arc=9pt,
  shadow={3pt}{-3pt}{0pt}{popink},left=14pt,right=14pt,top=11pt,bottom=11pt,before skip=11pt,after skip=15pt,
  fontupper=\popsemi\fontsize{10.6}{14.5}\selectfont\color{popink},
  fontlower=\popsemi\fontsize{10.6}{14.5}\selectfont\color{popink}}}
% Coche dessinée (le glyphe ✓ n'existe pas dans les polices du document)
\newcommand{\popcheck}[1]{\tikz[baseline=-0.5ex]\draw[#1,line width=1.6pt,line cap=round,line join=round](-0.13,0.0)--(-0.04,-0.09)--(0.13,0.1);}
\providecommand{\coche}{\popcheck{popgreen}}
\providecommand{\resultat}[1]{\par\smallskip\noindent\fcolorbox{popgreen}{popgreenL}{\parbox{\dimexpr\linewidth-2\fboxsep-2\fboxrule\relax}{\textbf{Résultat.} #1}}\par\smallskip}
\newcommand{\popboxhead}[3]{\poppill{#1}{#2}{white}\ifx\relax#3\relax\else\hspace{6pt}{\popxbold\fontsize{10.6}{12}\selectfont\color{popink}#3}\fi\par\vspace{7pt}}

\newtcolorbox{aretenir}[1][]{popbase,colframe=popgreen,before upper={\popboxhead{À retenir}{popgreen}{#1}}}
% Texte en espagnol (document de lecture, dialogue, poème, extrait) : cadre
% d'étude. Un titre optionnel (source, auteur) s'affiche dans l'en-tête.
\newtcolorbox{textebox}[1][]{popbase,colframe=popblue,before upper={\popboxhead{Texto}{popblue}{#1}\begin{otherlanguage}{spanish}},after upper={\end{otherlanguage}}}
% Marques ✓ / ✗ (forme correcte / fautive) souvent utilisées par les modèles
\providecommand{\cross}{\textcolor{poppink}{\ensuremath{\times}}}
\providecommand{\xmark}{\cross}\providecommand{\crossmark}{\cross}\providecommand{\cmark}{\ensuremath{\checkmark}}
% Ligne de réponse / justification à compléter (inventée par les modèles)
\providecommand{\justification}[1]{\par\noindent\textit{Justification~:} \dotfill\par}
% \textipa (package tipa non chargé) : la police ne couvre pas l'API -> simple passage
\providecommand{\textipa}[1]{#1}
% Mot ou phrase en espagnol à l'intérieur d'un paragraphe français
\newcommand{\esp}[1]{\begingroup\selectlanguage{spanish}\itshape #1\endgroup}
\newtcolorbox{exemplebox}[1][]{popbase,colframe=poporange,before upper={\popboxhead{Exemple}{poporange}{#1}}}
\newtcolorbox{definition}[1][]{popbase,colframe=popblue,before upper={\popboxhead{Définition}{popblue}{#1}}}
\newtcolorbox{propriete}[1][]{popbase,colframe=poppurple,before upper={\popboxhead{Propriété}{poppurple}{#1}}}
\newtcolorbox{methode}[1][]{popbase,colframe=popgold,before upper={\popboxhead{Méthode}{popgold}{#1}}}
\newtcolorbox{attention}[1][]{popbase,colframe=poppink,before upper={\popboxhead{Attention}{poppink}{#1}}}
\newtcolorbox{experience}[1][]{popbase,colframe=popblue,colback=popblueL,before upper={\popboxhead{Expérience}{popblue}{#1}}}
% « Le savais-tu ? » : carte violette pleine (contexte, culture, Cameroun)
\newtcolorbox{savaistu}[1][]{popbase,colback=poppurple,colframe=popink,
  fontupper=\popsemi\fontsize{10.4}{14.2}\selectfont\color{white},
  before upper={\poppill{Le savais-tu\,?}{white}{poppurple}\ifx\relax#1\relax\else\hspace{6pt}{\popxbold\color{white}#1}\fi\par\vspace{7pt}}}
% Exercice résolu : énoncé en haut, \tcblower puis la solution sur fond crème
\newcounter{popexo}\newcounter{popexr}
\newtcolorbox{exoresolu}[1][]{popbase,colframe=poporange,colbacklower=popcream,
  segmentation style={popink,line width=1.2pt,dashed},
  before upper={\stepcounter{popexr}\popboxhead{Exercice résolu \thepopexr}{poporange}{#1}},
  before lower={\poppill{Solution}{popink}{popcream}\par\vspace{6pt}}}
% Exercice (non résolu en place) — la correction est dans la partie Corrigés
\newtcolorbox{exercice}[1][]{popbase,colframe=popink,
  before upper={\stepcounter{popexo}\popboxhead{Exercice \thepopexo}{popink}{#1}}}
% Corrigé
\newtcolorbox{corrige}[1][]{popbase,colframe=popgreen,colback=popgreenL,
  before upper={\popboxhead{Corrigé}{popgreen}{#1}}}
% Objectifs / prérequis / bilan
\newtcolorbox{objectifs}{popbase,colframe=popink,colback=poporangeL,
  before upper={\popboxhead{Objectifs de la leçon}{popink}{}}}
\newtcolorbox{prerequis}{popbase,colframe=popink,colback=popblueL,
  before upper={\popboxhead{Prérequis}{popblue}{}}}
\newtcolorbox{bilan}{popbase,colframe=popink,colback=white,boxrule=2.8pt,
  before upper={\popboxhead{Fiche bilan}{poporange}{L'essentiel à connaître pour l'examen}}}
% Figure encadrée avec légende
\newtcolorbox{popfigure}[1][]{enhanced,breakable=false,colback=white,colframe=popline,boxrule=1.4pt,arc=8pt,
  left=10pt,right=10pt,top=10pt,bottom=8pt,before skip=10pt,after skip=12pt,halign=center,
  lower separated=false,halign lower=center,
  fontlower=\popsemi\fontsize{9}{11.5}\selectfont\color{popmuted},
  #1}
\newcommand{\legende}[1]{\par\vspace{6pt}{\popsemi\fontsize{9}{11.5}\selectfont\color{popmuted}#1}}

% ============================================================
% Signature OnBuch+
% ============================================================
\newcommand{\poplogoinline}[1]{{\poplogo\fontsize{#1}{#1}\selectfont\color{popink}OnBuch\color{poporange}+}}
\newcommand{\signatureonbuch}{%
  \begin{tcolorbox}[enhanced,colback=popink,colframe=popink,boxrule=2.2pt,arc=10pt,
      drop shadow=poporange,left=14pt,right=14pt,top=10pt,bottom=10pt,before skip=0pt,after skip=0pt]
    \tikz[baseline=-0.7ex]\node[circle,fill=popgreen,draw=popcream,line width=1.3pt,minimum size=22pt,inner sep=0]{\popcheck{white}};%
    \hspace{9pt}\begin{minipage}[c]{0.62\linewidth}
      {\popxbold\fontsize{12}{13}\selectfont\color{popcream}Signé\ \textbullet\ {\poplogo OnBuch\color{poporange}+}}\par\vspace{2pt}
      {\raggedright\popsemi\fontsize{8}{10.5}\selectfont\color{popline}\MakeUppercase{Équipe pédagogique OnBuch+}\\\MakeUppercase{Conforme au programme officiel MINESEC}\par}
    \end{minipage}\hfill
    {\poplogo\fontsize{22}{22}\selectfont\color{popcream}OnBuch\color{poporange}+}
  \end{tcolorbox}%
}

% ============================================================
% Page de garde
% #1 titre · #2 matière · #3 niveau · #4 module · #5 n° leçon · #6 durée ·
% #7 description / objectifs (texte ou itemize)
% ============================================================
\newcommand{\pagegarde}[7]{%
  \thispagestyle{empty}
  \newgeometry{a4paper,margin=1.7cm,top=1.6cm,bottom=1.4cm}%
  \begin{tikzpicture}[remember picture,overlay]
    \fill[poporange!18] ([xshift=-3.2cm,yshift=-3.6cm]current page.north east) circle (4.6cm);
    \fill[poppurple!14] ([xshift=2.4cm,yshift=7.6cm]current page.south west) circle (3.8cm);
  \end{tikzpicture}%
  {\poplogo\fontsize{30}{30}\selectfont\color{popink}OnBuch\color{poporange}+}\hfill
  \raisebox{0.6ex}{\poppill{Cours signé \textbullet\ Premium}{popink}{popcream}}\par
  \vspace{1pt}\popsquiggle{poppurple}\par
  \vspace{8pt}
  \begin{tcolorbox}[enhanced,colback=poporange,colframe=popink,boxrule=2.8pt,arc=13pt,
      drop shadow=popink,left=18pt,right=18pt,top=12pt,bottom=13pt,after skip=10pt]
    \poppill{#2 \textbullet\ #3}{popink}{popcream}\hspace{5pt}\poppill{#4}{white}{popink}\par\vspace{8pt}
    {\popdisplay\fontsize{14}{14}\selectfont\color{popink}LEÇON #5}\par\vspace{4pt}
    {\raggedright\hyphenpenalty=10000\exhyphenpenalty=10000\popdisplay\fontsize{25}{27}\selectfont\color{popcream}\MakeUppercase{#1}\par}
    \vspace{8pt}
    {\popxbold\fontsize{9.5}{9.5}\selectfont\color{popink}\MakeUppercase{Durée conseillée : #6}}
  \end{tcolorbox}
  \begin{tcolorbox}[enhanced,colback=white,colframe=popink,boxrule=2.2pt,arc=10pt,
      drop shadow=popink,left=16pt,right=16pt,top=10pt,bottom=10pt,after skip=8pt]
    \poppill{Dans cette leçon}{poppurple}{white}\par\vspace{5pt}
    \popsemi\fontsize{9.3}{12.6}\selectfont\color{popink}#7
  \end{tcolorbox}
  \noindent\begin{minipage}[t]{0.487\linewidth}
    \begin{tcolorbox}[enhanced,colback=popgreen,colframe=popink,boxrule=2.2pt,arc=10pt,
        drop shadow=popink,left=11pt,right=11pt,top=10pt,bottom=10pt,height=3.1cm,valign=center]
      \tcbox[colback=white,colframe=popink,boxrule=1.3pt,arc=5pt,boxsep=0pt,nobeforeafter,
        left=3pt,right=3pt,top=3pt,bottom=3pt,valign=center]{\qrcode[height=1.9cm]{https://play.google.com/store/apps/details?id=com.luvvix.onbuch&hl=fr}}%
      \hspace{8pt}\begin{minipage}[c]{0.5\linewidth}
        {\popdisplay\fontsize{11}{12.5}\selectfont\color{white}\MakeUppercase{Télécharge OnBuch}}\par\vspace{4pt}
        {\raggedright\popsemi\fontsize{8.4}{11}\selectfont\color{white}Gratuit \textbullet\ Google Play\\Tuteur IA, annales, concours, résultats\par}
      \end{minipage}
    \end{tcolorbox}
  \end{minipage}\hfill
  \begin{minipage}[t]{0.487\linewidth}
    \begin{tcolorbox}[enhanced,colback=poppurple,colframe=popink,boxrule=2.2pt,arc=10pt,
        drop shadow=popink,left=11pt,right=11pt,top=10pt,bottom=10pt,height=3.1cm,valign=center]
      \tikz[baseline=-0.7ex]\node[circle,fill=white,draw=popink,line width=1.3pt,minimum size=1.6cm,inner sep=0]{\popdisplay\fontsize{24}{24}\selectfont\color{poppurple}+};%
      \hspace{8pt}\begin{minipage}[c]{0.6\linewidth}
        {\popdisplay\fontsize{11}{12.5}\selectfont\color{white}\MakeUppercase{Passe à OnBuch+}}\par\vspace{4pt}
        {\raggedright\popsemi\fontsize{8.4}{11}\selectfont\color{white}Tous les cours signés, Tuteur IA illimité, fiches PDF — onglet OnBuch+ de l'app\par}
      \end{minipage}
    \end{tcolorbox}
  \end{minipage}
  \vspace{8pt}
  \popcontenu
  \vfill
  \signatureonbuch
  \restoregeometry
  \newpage
}

% Rangée « Ce cours contient » (page de garde)
\newcommand{\poptile}[3]{% #1 couleur #2 gros texte #3 légende
  \begin{tcolorbox}[enhanced,colback=#1,colframe=popink,boxrule=2pt,arc=9pt,shadow={2.5pt}{-2.5pt}{0pt}{popink},
      left=6pt,right=6pt,top=9pt,bottom=9pt,halign=center,valign=center,height=1.75cm,width=\linewidth,nobeforeafter]
    {\popdisplay\fontsize{12.5}{13}\selectfont\color{white}\MakeUppercase{#2}}\par\vspace{3pt}
    {\centering\popsemi\fontsize{7.8}{9.5}\selectfont\color{white}#3\par}
  \end{tcolorbox}}
\newcommand{\popcontenu}{%
  \par\noindent{\popxboldls\fontsize{7.5}{7.5}\selectfont\color{popmuted}CE COURS SIGNÉ CONTIENT}\par\vspace{7pt}
  \noindent\begin{minipage}[t]{0.235\linewidth}\poptile{popblue}{Cours}{complet, illustré, pas à pas}\end{minipage}\hfill
  \begin{minipage}[t]{0.235\linewidth}\poptile{poporange}{Méthodes}{et exercices résolus}\end{minipage}\hfill
  \begin{minipage}[t]{0.235\linewidth}\poptile{poppink}{Exercices}{type examen + corrigés}\end{minipage}\hfill
  \begin{minipage}[t]{0.235\linewidth}\poptile{popgreen}{Fiche bilan}{l'essentiel à retenir}\end{minipage}\par}

% Sommaire
\newtcolorbox{sommaire}{enhanced,colback=white,colframe=popink,boxrule=2.2pt,arc=10pt,
  drop shadow=popink,left=18pt,right=18pt,top=13pt,bottom=13pt,before skip=6pt,after skip=20pt,
  before upper={\poppill{Sommaire}{poppurple}{white}\par\vspace{9pt}\popsemi\fontsize{10.6}{14.5}\selectfont\color{popink}}}

% Signature de fin de cours — poussée tout en bas de la dernière page
\newcommand{\findecours}{%
  \par\vfill\nopagebreak
  \begin{center}\popsquiggle{poporange}\end{center}\vspace{4pt}
  \signatureonbuch
}
\AtBeginDocument{\def\llbracket{\mathopen{[\mkern-3.5mu[}}\def\rrbracket{\mathclose{]\mkern-3.5mu]}}} % intervalles d'entiers (glyphe ⟦ absent de la police)
% Glyphes absents de Fira Math : redéfinis à partir de glyphes disponibles
\AtBeginDocument{%
  \def\setminus{\mathbin{\backslash}}\let\smallsetminus\setminus
  \def\vdots{\mathord{\vbox{\baselineskip3.2pt\lineskiplimit0pt\kern4pt\hbox{.}\hbox{.}\hbox{.}}}}%
  \def\ddots{\mathinner{\mkern1mu\raise6.5pt\hbox{.}\mkern2mu\raise3.6pt\hbox{.}\mkern2mu\raise0.7pt\hbox{.}\mkern1mu}}%
  \def\triangle{\mathord{\scalebox{0.9}{$\symup{\Delta}$}}}%
  \def\nabla{\mathord{\raisebox{\depth}{\scalebox{1}[-1]{$\symup{\Delta}$}}}}%
  \def\checkmark{\popcheck{popdark}}%
  \def\star{\mathbin{\ast}}\def\bigstar{\mathord{\ast}}%
  \def\diamond{\mathbin{\scalebox{0.6}{$\Diamond$}}}%
  \def\bigcup{\mathop{\mathchoice{\raisebox{-0.3em}{\scalebox{1.7}{$\cup$}}}{\scalebox{1.25}{$\cup$}}{\cup}{\cup}}\displaylimits}%
  \def\bigcap{\mathop{\mathchoice{\raisebox{-0.3em}{\scalebox{1.7}{$\cap$}}}{\scalebox{1.25}{$\cap$}}{\cap}{\cap}}\displaylimits}%
  \def\lesseqqgtr{\mathrel{\lessgtr}}\def\bigoplus{\mathop{\oplus}\displaylimits}%
}
\providecommand{\ssi}{si et seulement si} % abréviation fréquente
\AtBeginDocument{\providecommand{\wideparen}{\overparen}} % arc de cercle

\begin{document}

\pagegarde{Lectura y comentario : contaminación y medidas de preservación ; creación de carteles de sensibilización}{Espagnol}{Tle A}{Módulo 2 : Environnement, santé et bien-être}{E09}{2 h}{Cette leçon propose la lecture et le commentaire d'un texte sur la pollution de l'air, de l'eau et des sols ainsi que sur les mesures de préservation de l'environnement. Elle fournit le lexique thématique indispensable et guide l'élève vers la création d'une affiche de sensibilisation en espagnol, compétence directement évaluée au baccalauréat.
\vspace{4pt}
\begin{multicols}{2}\small
\begin{itemize}
\item À la fin de cette leçon, je sais lire et comprendre globalement et finement un texte espagnol sur la pollution et la protection de l'environnement.
\item Je sais identifier les idées principales et les articulations d'un texte argumentatif écologique.
\item Je sais employer correctement le lexique de la pollution (contaminación, calentamiento global, deforestación, sequía, capa de ozono) et des solutions (energías renovables, reciclaje).
\item Je sais rédiger un commentaire structuré d'un texte (introduction, développement, conclusion) en espagnol.
\item Je sais utiliser l'impératif affirmatif et négatif pour formuler des appels à l'action.
\item Je sais construire des slogans percutants avec les verbes d'appel à l'action (proteger, reciclar, ahorrar, cuidar).
\end{itemize}\end{multicols}}

\begin{sommaire}
\begin{multicols}{2}
\begin{enumerate}
\item Texte support : « Un planeta en peligro »
\item Compréhension du texte
\item Lexique thématique : la pollution et la protection
\item Méthode du commentaire de texte
\item L'affiche de sensibilisation : impératif et slogans
\item Production guidée : créer son affiche
\item Activité d'intégration
\item Méthodes et exercices résolus
\item Exercices
\item Corrigés des exercices
\item Fiche bilan
\end{enumerate}
\end{multicols}
\end{sommaire}

\begin{objectifs}
\begin{itemize}
\item À la fin de cette leçon, je sais lire et comprendre globalement et finement un texte espagnol sur la pollution et la protection de l'environnement.
\item Je sais identifier les idées principales et les articulations d'un texte argumentatif écologique.
\item Je sais employer correctement le lexique de la pollution (contaminación, calentamiento global, deforestación, sequía, capa de ozono) et des solutions (energías renovables, reciclaje).
\item Je sais rédiger un commentaire structuré d'un texte (introduction, développement, conclusion) en espagnol.
\item Je sais utiliser l'impératif affirmatif et négatif pour formuler des appels à l'action.
\item Je sais construire des slogans percutants avec les verbes d'appel à l'action (proteger, reciclar, ahorrar, cuidar).
\item Je sais concevoir une affiche de sensibilisation complète en espagnol (titre, slogan, consignes, image mentale).
\item Je sais éviter les faux-amis et les erreurs fréquentes des francophones dans ce champ lexical.
\end{itemize}
\end{objectifs}

\begin{prerequis}
\begin{itemize}
\item L'impératif affirmatif et négatif (Première) : ¡Recicla! / ¡No contamines!
\item Le présent de l'indicatif des verbes réguliers et irréguliers fréquents
\item Le vocabulaire de base de la nature (árbol, río, tierra, aire) vu en Seconde/Première
\item Les connecteurs logiques simples (porque, sin embargo, además, por eso)
\item La structure du commentaire de texte vue en Première
\end{itemize}
\end{prerequis}

\newpage

\coursec{Texte support : « Un planeta en peligro »}

\courssub{Situation d'accroche et problématique}

À Douala, après une forte pluie, les rues se remplissent de déchets plastiques qui bouchent les caniveaux et transforment certains quartiers en marécages. À Yaoundé, la fumée des motos-taxis et des vieux camions rend l'air irrespirable aux heures de pointe, et les passants se couvrent le nez en traversant le rond-point de la poste. De l'autre côté de l'Atlantique, l'Amazonie brûle, Santiago du Chili suffoque sous le smog et les plages du Pacifique reçoivent chaque jour des tonnes de plastique. Face à cela, partout dans le monde hispanophone, des jeunes créent des affiches et des slogans pour défendre la planète : \esp{¡No contamines!}, \esp{¡Recicla, reutiliza, reduce!}, \esp{¡Cuida el agua, cuida la vida!}

Et toi, saurais-tu lire et comprendre un texte espagnol sur la pollution, puis concevoir une affiche capable de mobiliser ton quartier ou ton lycée ? C'est exactement la double compétence que cette leçon va développer : \textbf{comprendre} un article de sensibilisation, puis \textbf{agir} par l'écrit.

\begin{methode}[Première lecture : observer avant de traduire]
Avant de traduire le moindre mot, un bon lecteur suit toujours trois étapes :
\begin{enumerate}
\item \textbf{Repérer le thème} : lire le titre et la première phrase. De quoi parle le texte ? Ici, le titre \esp{Un planeta en peligro} annonce immédiatement le sujet : la planète menacée.
\item \textbf{Identifier le type de texte} : d'où vient le document ? Un article de revue écologique n'a ni le même ton ni le même but qu'un poème ou qu'un règlement administratif. Ici, il s'agit d'un \esp{artículo de sensibilización} : un article qui informe et qui cherche à convaincre.
\item \textbf{Observer l'organisation} : compter les paragraphes et lire la première phrase de chacun. On découvre ainsi le plan de l'auteur avant même de comprendre tous les détails.
\end{enumerate}
Ce n'est qu'après ces trois étapes que l'on passe à la lecture détaillée, dictionnaire ou glossaire à la main.
\end{methode}

\courssub{Lecture silencieuse du texte}

Lis maintenant le texte suivant en silence, une première fois, sans t'arrêter sur les mots inconnus : le glossaire placé après le texte t'aidera ensuite.

\begin{textebox}[« Un planeta en peligro » — Revista \emph{Tierra Viva}, artículo de sensibilización]
El aire de las grandes ciudades está cada vez más contaminado. En Ciudad de México, en Santiago de Chile o en Lima, millones de coches y de fábricas expulsan cada día humos tóxicos que provocan enfermedades respiratorias, sobre todo en los niños y en los ancianos. Este fenómeno, unido al calentamiento global, amenaza la salud de todos los habitantes del planeta.

Pero el aire no es la única víctima. Las aguas de los ríos y de los mares reciben cada año millones de toneladas de residuos plásticos y de productos químicos. Los vertederos ilegales contaminan también los suelos, que se vuelven estériles. En muchas regiones, la sequía castiga a los campesinos y los bosques desaparecen a causa de la deforestación: cada minuto, una superficie de árboles equivalente a varios campos de fútbol es destruida.

¿Cuáles son las causas de esta situación? Los expertos señalan el consumo excesivo de energías sucias como el carbón y el petróleo, la producción masiva de plásticos y la falta de educación ecológica de la población. Además, la destrucción de los bosques reduce la capacidad de la Tierra para absorber el dióxido de carbono y debilita la capa de ozono, que nos protege de los rayos del sol.

Sin embargo, existen soluciones. Las energías renovables, como la energía solar y la energía eólica, no contaminan y son inagotables. El reciclaje de los residuos permite reducir la cantidad de basura que llega a los vertederos. En muchos países, los gobiernos han decidido prohibir las bolsas de plástico y plantar millones de árboles.

Por último, los jóvenes desempeñan un papel fundamental. En las escuelas de España y de América Latina, miles de estudiantes participan en campañas de limpieza y crean carteles con lemas como «¡Recicla, reutiliza, reduce!». Concienciar a la población es la primera batalla: si cada uno de nosotros cambia sus hábitos, el planeta todavía puede salvarse. El medio ambiente es nuestra casa común; protegerlo es proteger la vida.
\end{textebox}

\textbf{Glossaire du texte}

\begin{center}
\begin{tabularx}{\linewidth}{|X|X|X|X|}
\hline
\rowcolor{popblueL} \textbf{Español} & \textbf{Français} & \textbf{Español} & \textbf{Français} \\
\hline
la contaminación & la pollution & los residuos & les déchets \\
\hline
el calentamiento global & le réchauffement climatique & el vertedero & la décharge \\
\hline
la deforestación & la déforestation & el medio ambiente & l'environnement \\
\hline
la sequía & la sécheresse & concienciar & sensibiliser \\
\hline
la capa de ozono & la couche d'ozone & el lema & le slogan \\
\hline
las energías renovables & les énergies renouvelables & el cartel & l'affiche \\
\hline
el humo & la fumée & inagotable & inépuisable \\
\hline
el campesino & le paysan & salvar & sauver \\
\hline
\end{tabularx}
\end{center}

\begin{definition}[Contaminación et medio ambiente]
\esp{La contaminación} désigne la \cle{dégradation du milieu naturel} (air, eau, sols) par des substances nocives : fumées, produits chimiques, plastiques, bruit. \esp{El medio ambiente} désigne \cle{l'environnement}, c'est-à-dire l'ensemble du milieu naturel qui entoure les êtres vivants. Attention à l'ordre des mots : on dit toujours \esp{medio ambiente} (et jamais \esp{ambiente medio}), et l'adjectif correspondant est \esp{medioambiental} : \esp{la protección medioambiental} « la protection de l'environnement ».
\end{definition}

\begin{attention}[Contaminación ou polución ?]
Les deux mots existent en espagnol : \esp{contaminación} et \esp{polución} signifient tous les deux « pollution ». Mais \esp{contaminación} est le terme \textbf{le plus courant et le plus sûr à l'écrit}, surtout dans les examens. Retiens aussi le verbe : \esp{contaminar} « polluer », et l'adjectif : \esp{contaminado/a} « pollué(e) ». Exemple : \esp{El río está muy contaminado.} « La rivière est très polluée. » Ne confonds pas avec le faux ami français : \esp{contaminar} ne signifie jamais « contaminer » au sens médical (maladie) ; pour une maladie, on dit \esp{contagiar}.
\end{attention}

\courssub{Lecture à voix haute et prononciation}

Après la lecture silencieuse, ton professeur lira le texte à voix haute, puis te demandera de lire un paragraphe à ton tour. Quelques pièges de prononciation à surveiller dans ce texte :

\begin{itemize}
\item \esp{calentamiento} se prononce « ka-lén-ta-mién-to » : le \esp{c} devant \esp{a} se prononce « k ».
\item \esp{sequía} se prononce « sé-kî-a » : l'accent sur le \esp{í} brise la diphtongue.
\item \esp{energías} se prononce « é-nèr-hî-as » : le \esp{g} devant \esp{i} se prononce comme un « h » aspiré.
\item \esp{concienciar} se prononce « kon-sién-siar » : le \esp{c} devant \esp{i} se prononce « s » en Amérique latine (et « th » en Espagne).
\item \esp{inagotable} : bien marquer les cinq syllabes « i-na-go-ta-blé ».
\end{itemize}

\courssub{Première compréhension globale}

Réponds maintenant, en français d'abord, aux questions essentielles.

\textbf{De quoi parle le texte ?} Le texte parle de la pollution de la planète : pollution de l'air dans les grandes villes, contamination des eaux et des sols, déforestation, mais aussi des solutions possibles.

\textbf{Quel est le but de l'auteur ?} Ce n'est pas seulement d'informer : le dernier paragraphe appelle clairement à l'action (\esp{si cada uno de nosotros cambia sus hábitos, el planeta todavía puede salvarse}). C'est donc un \textbf{article de sensibilisation} : informer pour convaincre et faire agir.

\textbf{Qui est visé ?} Le texte s'adresse au grand public, et en particulier aux jeunes, puisque le dernier paragraphe parle des écoles et des campagnes d'affichage.

\begin{exoresolu}[Questions de compréhension globale]
Énoncé : réponds en espagnol, par une phrase complète. 1) ¿De qué habla el texto? 2) ¿Cuál es el objetivo del autor? 3) ¿A quiénes se dirige el artículo?
\tcblower
Solution détaillée :
1) \esp{El texto habla de la contaminación del aire, del agua y de los suelos, y de las soluciones para proteger el planeta.} « Le texte parle de la pollution de l'air, de l'eau et des sols, et des solutions pour protéger la planète. »
2) \esp{El objetivo del autor es informar y sensibilizar a los lectores, sobre todo a los jóvenes.} « Le but de l'auteur est d'informer et de sensibiliser les lecteurs, surtout les jeunes. »
3) \esp{El artículo se dirige al público en general y en particular a los estudiantes.} « L'article s'adresse au public en général et en particulier aux élèves. »
Méthode : pour répondre, reprends les mots de la question (\esp{hablar de}, \esp{el objetivo}, \esp{dirigirse a}) et complète-les avec les informations du texte. Ne recopie jamais une phrase entière sans la reformuler.
\end{exoresolu}

\courssub{La structure du texte : problème, causes, solutions}

Un article de sensibilisation suit presque toujours la même logique en trois temps, que l'on retrouve ici très clairement :

\begin{itemize}
\item \textbf{Paragraphes 1 et 2 : el problema.} L'auteur dresse le constat : l'air des villes est contaminé, les eaux reçoivent des déchets, les sols deviennent stériles, les forêts disparaissent.
\item \textbf{Paragraphe 3 : las causas.} L'auteur explique l'origine du mal : énergies sales, production massive de plastique, manque d'éducation écologique.
\item \textbf{Paragraphes 4 et 5 : las soluciones.} L'auteur propose des remèdes : énergies renouvelables, recyclage, interdiction des sacs plastiques, et surtout la mobilisation des jeunes.
\end{itemize}

\begin{popfigure}
\begin{tikzpicture}[
  box/.style={draw=popblue, thick, rounded corners=4pt, fill=popblueL, align=center, minimum width=3.4cm, minimum height=1.5cm, font=\small},
  fleche/.style={-{Stealth[length=3mm]}, very thick, poporange}
]
\node[box, align=center] (p) at (0,0){\textbf{El problema}\\ § 1--2 : aire, agua,\\ suelos contaminados};
\node[box, fill=poporangeL, draw=poporange, align=center] (c) at (5.2,0){\textbf{Las causas}\\ § 3 : energías sucias,\\ plásticos, ignorancia};
\node[box, fill=popgreenL, draw=popgreen, align=center] (s) at (10.4,0){\textbf{Las soluciones}\\ § 4--5 : renovables,\\ reciclaje, jóvenes};
\draw[fleche] (p) -- (c);
\draw[fleche] (c) -- (s);
\end{tikzpicture}
\legende{Organigramme de la structure du texte : constat, explication, remèdes.}
\end{popfigure}

Pour mémoriser le lexique, rien de tel qu'une carte mentale : le mot central \esp{la contaminación} se divise en trois branches correspondant aux trois milieux naturels attaqués.

\begin{popfigure}
\begin{tikzpicture}[
  centro/.style={circle, draw=popdark, very thick, fill=poppurpleL, align=center, minimum size=2.6cm, font=\small\bfseries},
  rama/.style={draw=popblue, thick, rounded corners=4pt, fill=popblueL, align=center, minimum width=3.6cm, font=\small},
  hoja/.style={font=\footnotesize, align=center}
]
\node[centro, align=center] (c) at (0,0){La\\ contaminación};
\node[rama, align=center] (a) at (-4.6,2.2){\textbf{El aire}\\ humos tóxicos, smog,\\ calentamiento global};
\node[rama, fill=popgreenL, draw=popgreen, align=center] (g) at (4.6,2.2){\textbf{El agua}\\ residuos plásticos,\\ productos químicos};
\node[rama, fill=popgoldL, draw=popgold, align=center] (s) at (0,-2.8){\textbf{Los suelos}\\ vertederos ilegales,\\ deforestación, sequía};
\draw[thick, popblue] (c) -- (a);
\draw[thick, popgreen] (c) -- (g);
\draw[thick, popgold] (c) -- (s);
\end{tikzpicture}
\legende{Carte mentale lexicale : les trois milieux touchés par la pollution.}
\end{popfigure}

\begin{aretenir}[Le plan-type d'un article de sensibilisation]
Retiens ce schéma, il te servira pour tous les textes écologiques de l'année : \textbf{constat} (le problème) $\rightarrow$ \textbf{explication} (les causes) $\rightarrow$ \textbf{proposition} (les solutions et l'appel à l'agir). Au baccalauréat, identifier ce plan dans un commentaire de texte est toujours valorisé.
\end{aretenir}

\courssub{Comparer les langues : le lexique de l'environnement}

Beaucoup de mots de ce champ lexical sont des mots transparents, proches du français, ce qui est une chance. Mais quelques pièges subsistent, que ce tableau résume :

\begin{center}
\begin{tabularx}{\linewidth}{|X|X|X|}
\hline
\rowcolor{popblueL} \textbf{Français} & \textbf{Español} & \textbf{Remarque} \\
\hline
la pollution & la contaminación & terme le plus courant ; \esp{polución} existe aussi \\
\hline
le réchauffement climatique & el calentamiento global & du verbe \esp{calentar} « chauffer » \\
\hline
l'environnement & el medio ambiente & ordre fixe des mots ; adjectif \esp{medioambiental} \\
\hline
les déchets & los residuos / la basura & \esp{residuos} = registre soutenu, \esp{basura} = courant \\
\hline
la décharge & el vertedero & du verbe \esp{verter} « verser, déverser » \\
\hline
la sécheresse & la sequía & accent obligatoire sur le \esp{í} \\
\hline
sensibiliser & concienciar & aussi \esp{sensibilizar} ; ne pas dire \esp{conscienciar} \\
\hline
les énergies renouvelables & las energías renovables & \esp{renovable}, pas \esp{renouvelable} \\
\hline
\end{tabularx}
\end{center}

\begin{attention}[Pièges pour les francophones]
\begin{itemize}
\item \esp{La sequía} prend un accent : écris \esp{sequía}, jamais \esp{sequia} (qui se prononcerait « sé-kia »).
\item \esp{Sensibilizar} et \esp{concienciar} signifient « sensibiliser » ; en revanche \esp{sensible} en espagnol veut dire « sensible » au français (émotif), pas « raisonnable » : c'est un faux ami partiel.
\item On dit \esp{proteger el medio ambiente} « protéger l'environnement » : le verbe \esp{proteger} se conjugue avec un \esp{j} devant \esp{o} et \esp{a} : \esp{yo protejo}, \esp{él protege}.
\item \esp{El cartel} signifie « l'affiche », pas « le cartel » au sens de groupe criminel (sens qui existe aussi, selon le contexte !).
\end{itemize}
\end{attention}

\begin{savaistu}[Des villes qui suffoquent dans une cuvette]
Ciudad de México et Santiago du Chili comptent parmi les villes les plus polluées du monde hispanophone. Leur situation géographique explique tout : entourées de montagnes, elles forment une sorte de « cuvette » où les fumées restent prisonnières, surtout quand il n'y a pas de vent. À Mexico, certains jours, les écoles ferment et la moitié des voitures doit rester au garage : c'est le programme \esp{Hoy no circula} « Aujourd'hui, on ne roule pas ». Un exemple concret à citer dans un commentaire de texte ou une expression écrite sur l'environnement !
\end{savaistu}

\begin{exercice}[À toi de jouer : repérage dans le texte]
1) Relève dans le texte trois mots du champ lexical de la pollution et trois mots du champ lexical des solutions.
2) Trouve dans le texte l'équivalent espagnol de : « les décharges illégales », « la couche d'ozone », « sensibiliser la population », « le planète peut encore être sauvée ».
3) Indique à quel paragraphe correspondent ces idées : a) le rôle des jeunes ; b) les origines du problème ; c) l'état des rivières et des mers.
\end{exercice}

\begin{corrige}[Exercice : repérage dans le texte]
1) Pollution : \esp{contaminación}, \esp{residuos}, \esp{deforestación} (aussi \esp{humos tóxicos}, \esp{sequía}). Solutions : \esp{energías renovables}, \esp{reciclaje}, \esp{campañas de limpieza} (aussi \esp{plantar árboles}, \esp{prohibir las bolsas de plástico}).
2) \esp{los vertederos ilegales} ; \esp{la capa de ozono} ; \esp{concienciar a la población} ; \esp{el planeta todavía puede salvarse}.
3) a) paragraphe 5 ; b) paragraphe 3 ; c) paragraphe 2.
\end{corrige}

Tu sais désormais lire globalement ce texte : tu connais son thème, son type, son but et son plan en trois mouvements. La prochaine étape consistera à approfondir la compréhension phrase par phrase, avant d'étudier le lexique en détail et d'apprendre à construire ton propre commentaire.

\coursec{Compréhension du texte}

Dans la section précédente, nous avons lu et observé le texte \esp{« Un planeta en peligro »}, un article de presse qui décrit les grandes menaces qui pèsent sur l'environnement et les solutions possibles. Nous allons maintenant vérifier et approfondir notre compréhension de ce texte, étape par étape, comme on vous le demandera à l'examen : d'abord une compréhension \cle{globale} (de quoi parle le texte ? à qui s'adresse-t-il ?), puis une compréhension \cle{fine}, paragraphe par paragraphe, en citant le texte pour justifier chaque réponse.

\courssub{Compréhension globale : thème, objectif, destinataires}

Avant de répondre aux questions, rappelons le texte support qui sert de base à tout le travail de cette section.

\begin{textebox}[Un planeta en peligro (rappel du texte support)]
El planeta Tierra está en peligro. Cada año, las fábricas y los coches emiten millones de toneladas de CO2 a la atmósfera, lo que provoca el calentamiento global. Según los científicos, la temperatura del planeta ha subido más de un grado en los últimos cien años.

El segundo problema es la contaminación del agua. Los ríos y los mares reciben cada día toneladas de residuos: se calcula que ocho millones de toneladas de plástico llegan al océano cada año. Además, la deforestación avanza rápidamente: cada minuto desaparece una superficie de bosque equivalente a varios campos de fútbol. Sin embargo, todavía hay esperanza.

Existen soluciones sencillas que todos podemos aplicar. Primero, podemos usar el transporte público o la bicicleta en lugar del coche. Además, es necesario reciclar los residuos y reducir el consumo de plástico. Por eso, muchas ciudades han prohibido las bolsas de plástico. También debemos apostar por las energías renovables, como la energía solar y la energía eólica.

En conclusión, proteger el medio ambiente es tarea de todos: los gobiernos deben tomar medidas, pero cada ciudadano puede actuar desde hoy. El futuro del planeta está en nuestras manos.
\end{textebox}

\textbf{Glossaire :} \esp{emitir} = émettre ; \esp{el calentamiento global} = le réchauffement climatique ; \esp{los residuos} = les déchets ; \esp{se calcula que} = on estime que ; \esp{la deforestación} = la déforestation ; \esp{un campo de fútbol} = un terrain de football ; \esp{reciclar} = recycler ; \esp{apostar por} = miser sur ; \esp{la energía eólica} = l'énergie éolienne ; \esp{tomar medidas} = prendre des mesures ; \esp{está en nuestras manos} = est entre nos mains.

Trois questions globales permettent de cerner n'importe quel texte :

\begin{enumerate}
\item \esp{¿De qué trata el texto?} (De quoi traite le texte ?) — Il s'agit du \cle{thème}.
\item \esp{¿Qué quiere el autor?} (Que veut l'auteur ?) — Il s'agit de l'\cle{objectif} : informer, convaincre, dénoncer, sensibiliser.
\item \esp{¿A quién se dirige el texto?} (À qui le texte s'adresse-t-il ?) — Il s'agit des \cle{destinataires}.
\end{enumerate}

\begin{propriete}[Les mots-outils pour répondre]
Pour rédiger une réponse de compréhension en espagnol, utilisez ces formules d'introduction :
\begin{itemize}
\item \esp{Según el texto, ...} = Selon le texte, ...
\item \esp{El autor dice que ... / El autor explica que ...} = L'auteur dit que / explique que ...
\item \esp{Se trata de ...} = Il s'agit de ...
\item \esp{El texto habla de ...} = Le texte parle de ...
\item \esp{El texto se dirige a ...} = Le texte s'adresse à ...
\end{itemize}
Ces formules montrent à l'examinateur que vous vous appuyez sur le document et non sur vos opinions personnelles.
\end{propriete}

\begin{exemplebox}[Réponses globales rédigées (modèles)]
\esp{El texto trata de la contaminación del planeta y de las soluciones para proteger el medio ambiente.} « Le texte traite de la pollution de la planète et des solutions pour protéger l'environnement. »

\esp{El autor quiere informar y sensibilizar a los lectores sobre los problemas ecológicos.} « L'auteur veut informer et sensibiliser les lecteurs aux problèmes écologiques. »

\esp{El texto se dirige a todos los ciudadanos, y en particular a los jóvenes, porque dice: « proteger el medio ambiente es tarea de todos ».} « Le texte s'adresse à tous les citoyens, et en particulier aux jeunes, car il dit : "protéger l'environnement est la tâche de tous". »
\end{exemplebox}

\begin{methode}[Répondre à une question de compréhension en trois étapes]
\begin{enumerate}
\item \textbf{Reformuler} la question pour introduire la réponse : \esp{¿Qué problemas cita el texto?} devient \esp{El texto cita varios problemas...}.
\item \textbf{Citer le texte} entre guillemets pour justifier : \esp{... porque dice: « los ríos y los mares reciben cada día toneladas de residuos ».}
\item \textbf{Traduire la citation} si la consigne le demande : « les fleuves et les mers reçoivent chaque jour des tonnes de déchets ».
\end{enumerate}
Une réponse complète = reformulation + citation + (éventuellement) traduction. Une réponse sans citation n'est jamais pleinement justifiée.
\end{methode}

\begin{attention}[Les deux erreurs qui coûtent des points au Bac]
\begin{itemize}
\item \textbf{Répondre en français} quand la consigne exige l'espagnol. Au Bac, la réponse doit être rédigée \textbf{en espagnol}, avec vos propres mots.
\item \textbf{Recopier de longs passages} du texte sans guillemets ni reformulation. Citez seulement le fragment utile (quelques mots), entre guillemets, et reformulez le reste. Recopier trois lignes entières est considéré comme une absence de compréhension.
\end{itemize}
\end{attention}

\courssub{Compréhension fine, paragraphe par paragraphe}

Passons maintenant à l'analyse détaillée. Pour chaque paragraphe, voici des questions de type \emph{vrai/faux justifié} (\esp{verdadero o falso, justifica tu respuesta}), le format le plus fréquent à l'examen.

\textbf{Paragraphe 1 (le réchauffement).} \esp{Verdadero o falso: « La temperatura del planeta no ha cambiado en los últimos cien años. »}

Réponse attendue : \esp{Falso. Según el texto, « la temperatura del planeta ha subido más de un grado en los últimos cien años ».} (Faux. Selon le texte, « la température de la planète a augmenté de plus d'un degré au cours des cent dernières années ».)

\textbf{Paragraphe 2 (l'eau et les forêts).} \esp{Verdadero o falso: « El plástico no llega nunca al océano. »}

Réponse attendue : \esp{Falso. El autor dice que « ocho millones de toneladas de plástico llegan al océano cada año ».} (Faux. L'auteur dit que « huit millions de tonnes de plastique arrivent dans l'océan chaque année ».)

\textbf{Paragraphe 3 (les solutions).} \esp{¿Qué medios de transporte propone el texto?}

Réponse attendue : \esp{El texto propone « usar el transporte público o la bicicleta en lugar del coche ».} (Le texte propose « d'utiliser les transports en commun ou le vélo à la place de la voiture ».)

\textbf{Paragraphe 4 (la conclusion).} \esp{¿Quién debe proteger el medio ambiente, según el autor?}

Réponse attendue : \esp{Según el autor, todos debemos protegerlo: « los gobiernos deben tomar medidas, pero cada ciudadano puede actuar desde hoy ».} (Selon l'auteur, nous devons tous le protéger : « les gouvernements doivent prendre des mesures, mais chaque citoyen peut agir dès aujourd'hui ».)

\begin{exemplebox}[Exemple de phrase du texte, traduite]
\esp{El aire está contaminado por los coches y las fábricas.} = « L'air est pollué par les voitures et les usines. »

On remarque la tournure passive avec \esp{está contaminado por} : en espagnol comme en français, le complément d'agent s'introduit par \esp{por} (= par). Attention à ne pas traduire \esp{por} par « pour » ici.
\end{exemplebox}

\courssub{Recherche dans le texte : problèmes et solutions}

L'exercice suivant consiste à extraire du texte \textbf{trois problèmes écologiques} cités et \textbf{trois solutions} proposées. C'est un travail de repérage : on cherche les mots-clés (\esp{contaminación}, \esp{deforestación}, \esp{reciclar}...) sans recopier les phrases entières.

\begin{exoresolu}[Problèmes et solutions : relevé guidé]
Énoncé : \esp{Busca en el texto tres problemas ecológicos y tres soluciones. Escríbelos en dos columnas.}
\tcblower
Solution détaillée :

\textbf{Problemas citados en el texto :}
\begin{enumerate}
\item \esp{el calentamiento global} (le réchauffement climatique) — \esp{« las fábricas y los coches emiten millones de toneladas de CO2 »} ;
\item \esp{la contaminación del agua por el plástico} (la pollution de l'eau par le plastique) — \esp{« ocho millones de toneladas de plástico llegan al océano cada año »} ;
\item \esp{la deforestación} (la déforestation) — \esp{« cada minuto desaparece una superficie de bosque »}.
\end{enumerate}

\textbf{Soluciones propuestas :}
\begin{enumerate}
\item \esp{usar el transporte público o la bicicleta} (utiliser les transports en commun ou le vélo) ;
\item \esp{reciclar los residuos y reducir el consumo de plástico} (recycler les déchets et réduire la consommation de plastique) ;
\item \esp{apostar por las energías renovables} (miser sur les énergies renouvelables).
\end{enumerate}
\end{exoresolu}

À votre tour de compléter ce tableau de synthèse, qui servira de base à votre future affiche :

\begin{popfigure}
\begin{tikzpicture}
\node[draw=popblue, fill=popblueL, rounded corners, minimum width=6.2cm, minimum height=0.9cm, align=center, font=\bfseries] (p) at (0,0) {Problemas citados en el texto};
\node[draw=popgreen, fill=popgreenL, rounded corners, minimum width=6.2cm, minimum height=0.9cm, align=center, font=\bfseries] (s) at (7.2,0) {Soluciones propuestas};
\node[draw=popline, fill=white, rounded corners, minimum width=6.2cm, minimum height=0.8cm, align=center] at (0,-1.1) {1. \ldots\ldots\ldots\ldots\ldots\ldots\ldots};
\node[draw=popline, fill=white, rounded corners, minimum width=6.2cm, minimum height=0.8cm, align=center] at (0,-2.1) {2. \ldots\ldots\ldots\ldots\ldots\ldots\ldots};
\node[draw=popline, fill=white, rounded corners, minimum width=6.2cm, minimum height=0.8cm, align=center] at (0,-3.1) {3. \ldots\ldots\ldots\ldots\ldots\ldots\ldots};
\node[draw=popline, fill=white, rounded corners, minimum width=6.2cm, minimum height=0.8cm, align=center] at (7.2,-1.1) {1. \ldots\ldots\ldots\ldots\ldots\ldots\ldots};
\node[draw=popline, fill=white, rounded corners, minimum width=6.2cm, minimum height=0.8cm, align=center] at (7.2,-2.1) {2. \ldots\ldots\ldots\ldots\ldots\ldots\ldots};
\node[draw=popline, fill=white, rounded corners, minimum width=6.2cm, minimum height=0.8cm, align=center] at (7.2,-3.1) {3. \ldots\ldots\ldots\ldots\ldots\ldots\ldots};
\draw[-{Stealth}, thick, poporange] (3.3,-2.1) -- (4.0,-2.1);
\end{tikzpicture}
\legende{Tableau à compléter : à chaque problème correspond une solution. La flèche symbolise le passage du constat à l'action.}
\end{popfigure}

\courssub{Relevé des chiffres et des faits}

Un article de presse utilise des \cle{chiffres} pour rendre l'information crédible. Savoir les repérer et les exploiter est une compétence d'examen. Relevons ceux du texte :

\begin{center}
\begin{tabularx}{\linewidth}{|X|X|X|}
\hline
\rowcolor{popblueL} \textbf{Cifra del texto} & \textbf{Hecho mencionado} & \textbf{Traduction} \\
\hline
\esp{más de un grado} & \esp{la temperatura ha subido} & la température a augmenté de plus d'un degré \\
\hline
\esp{ocho millones de toneladas} & \esp{plástico que llega al océano cada año} & plastique qui arrive dans l'océan chaque année \\
\hline
\esp{cada minuto} & \esp{desaparece una superficie de bosque} & une surface de forêt disparaît chaque minute \\
\hline
\esp{los últimos cien años} & \esp{período del calentamiento} & période du réchauffement constaté \\
\hline
\end{tabularx}
\end{center}

\begin{attention}[Chiffres : les pièges des francophones]
\begin{itemize}
\item \esp{millones de toneladas} = des \textbf{millions} de tonnes (et non « milliards » ; \esp{mil millones} = un milliard).
\item \esp{un grado} = un \textbf{degré} (de température), pas « une grade ».
\item En espagnol, \esp{ocho} est invariable ; \esp{millones} prend toujours la préposition \esp{de} devant le nom : \esp{millones de personas} (des millions de personnes).
\end{itemize}
\end{attention}

\courssub{Relevé des connecteurs logiques et de leur fonction}

Les \cle{connecteurs logiques} (\esp{los conectores}) sont les charnières du texte : ils organisent l'argumentation. Les repérer, c'est comprendre la logique de l'auteur.

\begin{center}
\begin{tabularx}{\linewidth}{|X|X|X|}
\hline
\rowcolor{popblueL} \textbf{Conector} & \textbf{Función} & \textbf{Équivalent français} \\
\hline
\esp{además} & añade una idea & de plus, en outre \\
\hline
\esp{sin embargo} & expresa una oposición & cependant, pourtant \\
\hline
\esp{por eso} & expresa una consecuencia & c'est pourquoi, donc \\
\hline
\esp{también} & añade una idea similar & aussi, également \\
\hline
\esp{en conclusión} & introduce el final & en conclusion \\
\hline
\esp{primero} & ordena las ideas & d'abord, premièrement \\
\hline
\end{tabularx}
\end{center}

Observons leur fonction dans le texte :

\begin{itemize}
\item \esp{« Sin embargo, todavía hay esperanza. »} : le connecteur marque un \textbf{contraste} entre le constat négatif (la planète en danger) et une ouverture positive. « Cependant, il y a encore de l'espoir. »
\item \esp{« Por eso, muchas ciudades han prohibido las bolsas de plástico. »} : le connecteur introduit une \textbf{conséquence} logique. « C'est pourquoi de nombreuses villes ont interdit les sacs en plastique. »
\item \esp{« Además, es necesario reciclar... »} : le connecteur \textbf{ajoute} une solution à la précédente. « De plus, il est nécessaire de recycler... »
\item \esp{« En conclusión, proteger el medio ambiente es tarea de todos. »} : le connecteur annonce la \textbf{conclusion} du raisonnement. « En conclusion, protéger l'environnement est l'affaire de tous. »
\end{itemize}

\begin{attention}[Faux amis et calques à éviter]
\begin{itemize}
\item \esp{sin embargo} ne signifie jamais « sans embargo » : c'est toujours « cependant ».
\item Ne traduisez pas « c'est pourquoi » par \esp{por qué} (= pourquoi, interrogatif). La conséquence se dit \esp{por eso} ou \esp{por lo tanto}.
\item \esp{en conclusión} s'écrit avec l'accent sur le \esp{ó} ; on ne dit pas \esp{en conclusión de que} pour « en concluant que » : dites simplement \esp{en conclusión, ...}.
\end{itemize}
\end{attention}

\courssub{Questions d'expression personnelle}

Après la compréhension, l'examen et la classe demandent souvent une prise de parole personnelle, toujours \textbf{en espagnol}. Deux questions classiques :

\begin{enumerate}
\item \esp{¿Qué problemas ecológicos hay en tu ciudad?} (Quels problèmes écologiques y a-t-il dans ta ville ?)
\item \esp{¿Qué haces tú para proteger el medio ambiente?} (Que fais-tu pour protéger l'environnement ?)
\end{enumerate}

\begin{exemplebox}[Modèles de réponses personnelles (contexte camerounais)]
\esp{En mi ciudad, Yaoundé, hay mucha contaminación del aire por los coches y las motos, y la gente tira plástico en las calles.} « Dans ma ville, Yaoundé, il y a beaucoup de pollution de l'air à cause des voitures et des motos, et les gens jettent du plastique dans les rues. »

\esp{En Douala, los ríos están contaminados y durante la estación seca hay problemas de agua.} « À Douala, les fleuves sont pollués et pendant la saison sèche il y a des problèmes d'eau. »

\esp{Para proteger el medio ambiente, yo reciclo las botellas, apago la luz cuando salgo de una habitación y planto árboles con mi familia.} « Pour protéger l'environnement, je recycle les bouteilles, j'éteins la lumière quand je sors d'une pièce et je plante des arbres avec ma famille. »
\end{exemplebox}

Notez les tournures utiles : \esp{hay mucha contaminación} (il y a beaucoup de pollution), \esp{la gente tira plástico} (les gens jettent du plastique), \esp{apagar la luz} (éteindre la lumière), \esp{plantar árboles} (planter des arbres).

\begin{experience}[Vérification collective et justification par citations]
Travail en deux temps, en classe entière :
\begin{enumerate}
\item Chaque élève rédige individuellement ses réponses aux questions de compréhension fine (vrai/faux des paragraphes 1 à 4) et aux deux questions personnelles.
\item Correction collective au tableau : pour chaque réponse, l'élève doit \textbf{lire sa réponse en espagnol}, puis \textbf{citer le passage du texte} qui la justifie, entre guillemets. La classe vérifie que la citation est exacte et que la réponse n'est pas un simple recopiage.
\end{enumerate}
Critères de réussite affichés au tableau : réponse en espagnol, citation courte entre guillemets, connecteur ou formule d'introduction (\esp{según el texto}, \esp{el autor dice que}).
\end{experience}

\begin{exercice}[À vous de jouer]
\begin{enumerate}
\item \esp{Verdadero o falso (justifica con una cita): « El texto dice que no hay ninguna solución para el planeta. »}
\item \esp{¿Qué dos energías renovables menciona el texto?}
\item \esp{¿Qué conector usa el autor para pasar de los problemas a la esperanza?}
\item Traduis en français : \esp{« Por eso, muchas ciudades han prohibido las bolsas de plástico. »}
\item \esp{¿Qué haces tú para proteger el medio ambiente?} (Réponds en trois phrases en espagnol.)
\end{enumerate}
\end{exercice}

\begin{corrige}[Exercice « À vous de jouer »]
\begin{enumerate}
\item \esp{Falso. El texto dice que « existen soluciones sencillas que todos podemos aplicar » y que « todavía hay esperanza ».} (Faux : le texte affirme qu'il existe des solutions simples que tous peuvent appliquer.)
\item \esp{El texto menciona la energía solar y la energía eólica.} (L'énergie solaire et l'énergie éolienne.)
\item \esp{Usa el conector « sin embargo »: « Sin embargo, todavía hay esperanza. »} (Il utilise le connecteur d'opposition « cependant ».)
\item « C'est pourquoi de nombreuses villes ont interdit les sacs en plastique. »
\item Exemple : \esp{Para proteger el medio ambiente, yo uso la bicicleta para ir al colegio. También reciclo el papel y no tiro plástico en la calle. Además, hablo con mis amigos sobre el calentamiento global.} (Pour protéger l'environnement, j'utilise le vélo pour aller au collège. Je recycle aussi le papier et je ne jette pas de plastique dans la rue. De plus, je parle avec mes amis du réchauffement climatique.)
\end{enumerate}
\end{corrige}

\begin{aretenir}[L'essentiel de la section]
\begin{itemize}
\item Une réponse de compréhension = \textbf{reformulation + citation courte entre guillemets + traduction si demandée}.
\item Les formules-clés : \esp{según el texto}, \esp{el autor dice que}, \esp{se trata de}.
\item Les connecteurs structurent l'argumentation : \esp{además} (addition), \esp{sin embargo} (opposition), \esp{por eso} (conséquence), \esp{en conclusión} (bilan).
\item Au Bac, on répond \textbf{en espagnol}, sans recopier de longs passages.
\item Ce travail de repérage (problèmes, solutions, chiffres, connecteurs) prépare directement la création de votre affiche de sensibilisation, objet de la fin de la leçon.
\end{itemize}
\end{aretenir}

\coursec{Lexique thématique : la pollution et la protection}

Dans la section précédente, nous avons lu et compris le texte \esp{Un planeta en peligro}. Nous avons repéré les idées principales : la planète souffre, l'air et l'eau sont contaminés, mais des solutions existent. Pour aller plus loin dans le commentaire de texte et pour préparer notre future affiche de sensibilisation, il faut maintenant \cle{maîtriser le vocabulaire}. Un bon commentaire, comme une bonne affiche, repose d'abord sur des mots justes. Organisons donc ce lexique en quatre champs : les problèmes, les milieux touchés, les solutions et les acteurs.

\courssub{Observation : le lexique en contexte}

Avant toute liste de mots, lisons ce court texte original. Soulignez mentalement tous les mots qui se rapportent à l'environnement.

\begin{textebox}[El lenguaje de la emergencia ecológica]
Cada mañana, en las grandes ciudades como Douala o Ciudad de México, el aire se llena de humo y de gases tóxicos que salen de los coches y de las fábricas. Esta contaminación del aire agrava el calentamiento global y el efecto invernadero. Al mismo tiempo, la deforestación avanza : cada año desaparecen miles de hectáreas de bosques, y la selva amazónica, el pulmón del planeta, está en peligro. En otras regiones, el problema es contrario : la sequía seca los ríos y los suelos, y los campesinos ya no pueden cultivar.

Sin embargo, hay esperanza. Los ecologistas y las ONG trabajan para concienciar a la población. Proponen soluciones concretas : usar energías renovables como la energía solar o la energía eólica, practicar el reciclaje, tomar el transporte público, plantar árboles y ahorrar agua y electricidad. « Reducir, reutilizar, reciclar » : estas tres palabras resumen la lucha contra la contaminación. Si cada ciudadano cuida el medio ambiente, podemos salvar el planeta y construir un futuro sostenible para nuestros hijos.
\end{textebox}

\textbf{Glossaire :} \esp{el humo} = la fumée ; \esp{los gases tóxicos} = les gaz toxiques ; \esp{agravar} = aggraver ; \esp{el pulmón} = le poumon ; \esp{secar} = assécher ; \esp{los campesinos} = les paysans ; \esp{sin embargo} = cependant ; \esp{resumir} = résumer ; \esp{los ciudadanos} = les citoyens.

\textbf{Questions de repérage :}
\begin{enumerate}
\item Cite trois problèmes écologiques mentionnés dans le texte.
\item Quelles solutions sont proposées dans le deuxième paragraphe ?
\item Qui travaille pour \esp{concienciar a la población} ?
\end{enumerate}

On observe que le vocabulaire se range naturellement en \cle{quatre familles} : les mots qui nomment les \textbf{problèmes}, ceux qui nomment les \textbf{milieux} (air, eau, sols), ceux qui nomment les \textbf{solutions}, et enfin les \textbf{verbes d'action} et les \textbf{acteurs}. C'est exactement la structure de notre carte mentale.

\begin{popfigure}
\begin{tikzpicture}[mindmap, grow cyclic, every node/.style={concept, text=white, font=\small\bfseries},
root concept/.append style={concept color=popdark, font=\bfseries},
level 1/.append style={level distance=3.4cm, sibling angle=90},
level 2/.append style={level distance=2.1cm, sibling angle=45, font=\scriptsize}]
\node[root concept]{Medio ambiente}
  child[concept color=poporange]{ node {Problemas}
      child { node {contaminación} }
      child { node {sequía} }
      child { node {deforestación} }
  }
  child[concept color=popblue]{ node {Medios afectados}
      child { node {el aire} }
      child { node {el agua} }
      child { node {los suelos} }
  }
  child[concept color=popgreen]{ node {Soluciones}
      child { node {reciclaje} }
      child { node {energías renovables} }
      child { node {plantar árboles} }
  }
  child[concept color=poppurple]{ node {Acciones}
      child { node {proteger} }
      child { node {reducir} }
      child { node {concienciar} }
  };
\end{tikzpicture}
\legende{Carte mentale du lexique de l'environnement : quatre champs lexicaux.}
\end{popfigure}

\courssub{Champ 1 — Les problèmes écologiques}

Voici les noms essentiels pour décrire les dangers qui menacent la planète. Remarquez le genre de chaque nom : en espagnol, le genre est souvent imprévisible pour un francophone, il faut l'apprendre avec le mot.

\begin{center}
\begin{tabularx}{\linewidth}{|X|X|X|}\hline
\rowcolor{popblueL} Español & Français & Remarque \\ \hline
\esp{la contaminación} & la pollution & féminin en -ción \\ \hline
\esp{el calentamiento global} & le réchauffement climatique & masculin \\ \hline
\esp{la deforestación} & la déforestation & féminin en -ción \\ \hline
\esp{la sequía} & la sécheresse & accent sur le i \\ \hline
\esp{el agujero en la capa de ozono} & le trou dans la couche d'ozone & \esp{agujero} = trou \\ \hline
\esp{el efecto invernadero} & l'effet de serre & masculin \\ \hline
\esp{los residuos / los desechos} & les déchets & toujours au pluriel \\ \hline
\esp{el humo} & la fumée & masculin \\ \hline
\esp{los gases tóxicos} & les gaz toxiques & pluriel \\ \hline
\end{tabularx}
\end{center}

\begin{exemplebox}[Les problèmes en phrases]
\esp{La deforestación destruye los bosques.} « La déforestation détruit les forêts. »

\esp{El calentamiento global derrite los glaciares.} « Le réchauffement climatique fait fondre les glaciers. »

\esp{La sequía seca los ríos del norte del país.} « La sécheresse assèche les rivières du nord du pays. »

\esp{Las fábricas echan gases tóxicos al aire.} « Les usines rejettent des gaz toxiques dans l'air. »
\end{exemplebox}

\begin{attention}[Ne confonds pas : sequía et inundación]
\esp{La sequía} (la sécheresse) est le manque d'eau ; son contraire est \esp{la inundación} (l'inondation), l'excès d'eau. Les deux mots prennent un accent : \esp{sequía}, \esp{inundación}. Erreur fréquente : écrire \esp{sequia} sans accent, ou confondre les deux termes dans une traduction.
\end{attention}

\courssub{Champ 2 — Les milieux touchés}

\begin{center}
\begin{tabularx}{\linewidth}{|X|X|X|}\hline
\rowcolor{popblueL} Español & Français & Remarque \\ \hline
\esp{el aire} & l'air & masculin \\ \hline
\esp{el agua} & l'eau & féminin ! mais \esp{el agua} au singulier \\ \hline
\esp{los suelos} & les sols & pluriel \\ \hline
\esp{los ríos} & les rivières, les fleuves & masculin pluriel \\ \hline
\esp{los océanos} & les océans & masculin pluriel \\ \hline
\esp{los bosques} & les forêts & masculin pluriel \\ \hline
\esp{la selva amazónica} & la forêt amazonienne & \esp{selva} = forêt tropicale \\ \hline
\end{tabularx}
\end{center}

\begin{attention}[Le cas étrange de « el agua »]
\esp{Agua} est un nom \textbf{féminin}, mais on dit \esp{el agua} (et non \esp{la agua}) au singulier pour éviter la rencontre de deux « a » dont le premier est accentué. Le nom reste féminin : \esp{el agua está contaminada} (et non \esp{contaminado}). Au pluriel, on retrouve la forme normale : \esp{las aguas}. Même phénomène avec \esp{el aula}, \esp{el águila}.
\end{attention}

\begin{attention}[« El medio ambiente » est masculin]
Malgré sa terminaison en \esp{-ente}, on dit \esp{el medio ambiente} (masculin) : \esp{proteger el medio ambiente} = « protéger l'environnement ». Ne dis jamais \esp{la media ambiente}.
\end{attention}

\courssub{Champ 3 — Les solutions}

\begin{definition}[Energías renovables]
Les \esp{energías renovables} sont les énergies qui se renouvellent naturellement : le soleil (\esp{energía solar}), le vent (\esp{energía eólica}), l'eau (\esp{energía hidráulica}). Elles s'opposent aux \esp{combustibles fósiles} (pétrole, charbon), qui s'épuisent et polluent.
\end{definition}

\begin{center}
\begin{tabularx}{\linewidth}{|X|X|}\hline
\rowcolor{popgreenL} Español & Français \\ \hline
\esp{el reciclaje} & le recyclage \\ \hline
\esp{el transporte público} & les transports en commun \\ \hline
\esp{la bicicleta} & le vélo \\ \hline
\esp{plantar árboles} & planter des arbres \\ \hline
\esp{ahorrar agua y electricidad} & économiser l'eau et l'électricité \\ \hline
\esp{las energías renovables} & les énergies renouvelables \\ \hline
\end{tabularx}
\end{center}

\begin{exemplebox}[Les solutions en phrases]
\esp{Debemos ahorrar agua.} « Nous devons économiser l'eau. »

\esp{En mi barrio, reciclamos el plástico y el vidrio.} « Dans mon quartier, nous recyclons le plastique et le verre. »

\esp{El gobierno quiere plantar millones de árboles.} « Le gouvernement veut planter des millions d'arbres. »

\esp{La energía solar es limpia y barata.} « L'énergie solaire est propre et bon marché. »
\end{exemplebox}

\begin{savaistu}[Un modèle hispanophone : le Costa Rica]
Le Costa Rica, petit pays d'Amérique centrale, produit la quasi-totalité de son électricité grâce aux énergies renouvelables (eau, vent, chaleur de la terre). C'est un modèle dans le monde hispanophone : la preuve qu'un pays peut se développer tout en protégeant \esp{el medio ambiente}.
\end{savaistu}

\courssub{Champ 4 — Les acteurs et les verbes d'action}

\begin{center}
\begin{tabularx}{\linewidth}{|X|X|X|}\hline
\rowcolor{poppurpleL} Verbe & Français & Exemple \\ \hline
\esp{proteger} & protéger & \esp{proteger los bosques} \\ \hline
\esp{preservar} & préserver & \esp{preservar el agua} \\ \hline
\esp{cuidar} & prendre soin de & \esp{cuidar el planeta} \\ \hline
\esp{reciclar} & recycler & \esp{reciclar el papel} \\ \hline
\esp{reutilizar} & réutiliser & \esp{reutilizar las bolsas} \\ \hline
\esp{reducir} & réduire & \esp{reducir los residuos} \\ \hline
\esp{contaminar} & polluer & \esp{contaminar el río} \\ \hline
\esp{destruir} & détruire & \esp{destruir la capa de ozono} \\ \hline
\esp{concienciar} & sensibiliser, conscientiser & \esp{concienciar a los jóvenes} \\ \hline
\esp{sensibilizar} & sensibiliser & \esp{sensibilizar a la población} \\ \hline
\end{tabularx}
\end{center}

Les acteurs : \esp{los ecologistas} (les écologistes), \esp{las ONG} (les ONG, organisations non gouvernementales), \esp{los ciudadanos} (les citoyens), \esp{el gobierno} (le gouvernement).

\begin{propriete}[La règle des 3 R]
La campagne écologique internationale se résume en trois verbes à l'infinitif : \esp{Reducir} (réduire sa consommation), \esp{Reutilizar} (réutiliser les objets), \esp{Reciclar} (recycler les matériaux). Ce sont trois verbes du 3\textsuperscript{e} groupe en \esp{-ir}. Attention : \esp{reducir} a une irrégularité à la 1\textsuperscript{re} personne du présent : \esp{yo reduzco}, puis \esp{tú reduces, él reduce, nosotros reducimos, vosotros reducís, ellos reducen}.
\end{propriete}

\begin{popfigure}
\begin{tikzpicture}[>=Stealth, node distance=2.6cm]
\node[draw=popgreen, fill=popgreenL, circle, minimum size=1.7cm, font=\small\bfseries, align=center] (a) {Reducir};
\node[draw=popblue, fill=popblueL, circle, minimum size=1.7cm, font=\small\bfseries, align=center, right=of a] (b) {Reutilizar};
\node[draw=poporange, fill=poporangeL, circle, minimum size=1.7cm, font=\small\bfseries, align=center, below=1.8cm of $(a)!0.5!(b)$] (c) {Reciclar};
\draw[->, very thick, popgreen] (a) to[bend left=20] (b);
\draw[->, very thick, popblue] (b) to[bend left=20] (c);
\draw[->, very thick, poporange] (c) to[bend left=20] (a);
\end{tikzpicture}
\legende{Le cycle des 3 R : Reducir, Reutilizar, Reciclar.}
\end{popfigure}

\courssub{Expressions utiles pour le commentaire et l'affiche}

Ces expressions toutes faites te serviront dans le commentaire de texte comme dans la rédaction de slogans :

\begin{itemize}
\item \esp{cuidar el medio ambiente} = « prendre soin de l'environnement » ;
\item \esp{luchar contra la contaminación} = « lutter contre la pollution » ;
\item \esp{salvar el planeta} = « sauver la planète » ;
\item \esp{un futuro sostenible} = « un avenir durable » ;
\item \esp{proteger la naturaleza} = « protéger la nature » ;
\item \esp{estar en peligro (de extinción)} = « être en danger (d'extinction) ».
\end{itemize}

\begin{attention}[Faux-amis et pièges lexicaux]
\begin{itemize}
\item \esp{Una fábrica} signifie « une usine », jamais « une fabrique » au sens artisanal : \esp{Las fábricas contaminan el aire} = « Les usines polluent l'air ».
\item « Sensible » au sens de « sensibilisé » ne se traduit pas par \esp{sensible} (qui veut dire « sensible, fragile »). On dit : \esp{concienciar a la gente} ou \esp{sensibilizar a la población} = « sensibiliser les gens / la population ».
\item \esp{Los desechos} désigne les déchets matériels (ordures, résidus) ; ne l'emploie pas au sens figuré.
\item \esp{La sequía} = la sécheresse ; ne pas confondre avec \esp{la inundación} (l'inondation).
\end{itemize}
\end{attention}

\begin{exoresolu}[Complète avec le mot qui convient]
Énoncé. Complète chaque phrase avec un mot du lexique étudié : \esp{sequía}, \esp{reciclaje}, \esp{deforestación}, \esp{energías renovables}, \esp{medio ambiente}.

1. \esp{La \_\_\_\_\_\_ destruye la selva amazónica.}

2. \esp{El sol y el viento son \_\_\_\_\_\_ \_\_\_\_\_\_.}

3. \esp{No llueve desde hace ocho meses : es la \_\_\_\_\_\_.}

4. \esp{El \_\_\_\_\_\_ del plástico reduce los residuos.}

5. \esp{Debemos proteger el \_\_\_\_\_\_ \_\_\_\_\_\_.}

\tcblower
Solution détaillée.

1. \esp{La \textbf{deforestación} destruye la selva amazónica.} — C'est la coupe des arbres qui détruit la forêt.

2. \esp{El sol y el viento son \textbf{energías renovables}.} — Soleil et vent se renouvellent naturellement.

3. \esp{No llueve desde hace ocho meses : es la \textbf{sequía}.} — L'absence de pluie, c'est la sécheresse (et non l'inondation !).

4. \esp{El \textbf{reciclaje} del plástico reduce los residuos.} — Recycler le plastique diminue les déchets.

5. \esp{Debemos proteger el \textbf{medio ambiente}.} — Attention au masculin : \esp{el medio ambiente}.
\end{exoresolu}

\begin{exercice}[À toi de jouer]
Traduis en espagnol : 1. « La fumée des usines pollue l'air des villes. » 2. « Nous devons économiser l'eau et l'électricité. » 3. « Les écologistes veulent sensibiliser les jeunes. » 4. « La sécheresse assèche les rivières et les sols. »
\end{exercice}

\begin{corrige}[Exercice « À toi de jouer »]
1. \esp{El humo de las fábricas contamina el aire de las ciudades.} — \esp{fábrica} = usine ; verbe \esp{contaminar}.

2. \esp{Debemos ahorrar agua y electricidad.} — \esp{ahorrar} = économiser.

3. \esp{Los ecologistas quieren concienciar a los jóvenes.} — ou \esp{sensibilizar a los jóvenes} ; jamais \esp{hacer sensibles}.

4. \esp{La sequía seca los ríos y los suelos.} — accent sur \esp{sequía}.
\end{corrige}

\begin{aretenir}[L'essentiel du lexique]
Quatre champs à retenir : les \textbf{problèmes} (\esp{contaminación, calentamiento global, deforestación, sequía, efecto invernadero}), les \textbf{milieux} (\esp{el aire, el agua, los suelos, los bosques}), les \textbf{solutions} (\esp{energías renovables, reciclaje, plantar árboles, ahorrar agua}) et les \textbf{actions} (\esp{proteger, reducir, reutilizar, reciclar, concienciar}). Pièges : \esp{el medio ambiente} est masculin, \esp{el agua} reste féminin, \esp{fábrica} = usine, « sensibiliser » = \esp{concienciar / sensibilizar}. Ce vocabulaire est ta boîte à outils pour le commentaire de texte et pour l'affiche de la section 5.
\end{aretenir}

\coursec{Méthode du commentaire de texte}

Dans la section précédente, nous avons constitué notre boîte à outils lexicale : \esp{contaminación}, \esp{calentamiento global}, \esp{deforestación}, \esp{capa de ozono}, \esp{energías renovables}... Posséder ces mots est indispensable, mais au baccalauréat, il ne suffit pas de connaître le vocabulaire : il faut savoir \cle{commenter un texte}, c'est-à-dire le présenter, en analyser les idées et les procédés, puis porter un jugement personnel justifié. Cette section vous donne la méthode complète, étape par étape, appliquée au texte « Un planeta en peligro » étudié en classe.

\courssub{Qu'est-ce qu'un commentaire de texte ?}

\begin{definition}[Le commentaire de texte]
Le \cle{commentaire de texte} (\esp{comentario de texto}) est un exercice écrit qui consiste à présenter un texte (sa nature, son thème, sa source), à en expliquer les idées principales, à analyser les moyens linguistiques employés par l'auteur (lexique, temps, ton, chiffres, impératifs) et à donner un avis personnel argumenté. Il se rédige \textbf{en espagnol}, de manière organisée, en trois parties : \esp{introducción}, \esp{desarrollo}, \esp{conclusión}.
\end{definition}

Attention à ne pas confondre le commentaire avec le \textbf{résumé} : résumer, c'est reformuler les idées ; commenter, c'est les \emph{expliquer}, montrer \emph{comment} l'auteur les exprime et \emph{pourquoi} son texte est efficace. Le mot d'ordre : ne jamais paraphraser sans analyser.

\courssub{Les trois parties du commentaire}

\begin{propriete}[Structure obligatoire du commentaire]
Tout commentaire de texte au baccalauréat comporte trois parties :
\begin{enumerate}
\item \textbf{Introducción} : présenter le texte — son \cle{type} (article, lettre, discours, poème), son \cle{thème}, sa \cle{source} (journal, revue, discours officiel), et formuler la \cle{problématique} (la question centrale à laquelle le commentaire va répondre).
\item \textbf{Desarrollo} : analyser d'abord les \cle{idées} (le problème, les causes, les solutions proposées), puis la \cle{langue} (lexique, procédés, ton), en citant et en traduisant les passages clés.
\item \textbf{Conclusión} : faire le \cle{bilan} de l'analyse, donner son avis personnel justifié et ouvrir sur une question plus large (\esp{apertura}).
\end{enumerate}
\end{propriete}

\begin{methode}[Plan type du commentaire au Bac]
\begin{enumerate}
\item \textbf{Présentation du texte} : type de document, thème, source, problématique (2 à 3 phrases).
\item \textbf{Résumé des idées principales} : le problème exposé, ses causes, les solutions proposées (avec citations courtes).
\item \textbf{Analyse d'un ou deux procédés} : lexique péjoratif, impératifs, chiffres, ton du texte — toujours appuyée sur une citation traduite.
\item \textbf{Avis personnel justifié} : êtes-vous d'accord avec l'auteur ? Pourquoi ? Terminez par une ouverture.
\end{enumerate}
\end{methode}

\begin{popfigure}
\begin{center}
\begin{tikzpicture}[
node distance=0.9cm,
etape/.style={rectangle, rounded corners=6pt, draw=popblue, thick, fill=popblueL, text width=8.2cm, align=center, inner sep=6pt},
fleche/.style={-{Stealth[length=3mm]}, thick, poporange}
]
\node[etape, align=center] (intro){\textbf{Introducción}\\ Tipo de texto + tema + fuente + \textbf{problemática}};
\node[etape, below=of intro, fill=popgreenL, draw=popgreen, align=center] (des){\textbf{Desarrollo}\\ \textbf{Ideas} : problema, causas, soluciones (citas traducidas)\\ \textbf{Lengua} : léxico, imperativos, cifras, tono};
\node[etape, below=of des, fill=popgoldL, draw=popgold, align=center] (conc){\textbf{Conclusión}\\ Balance + opinión personal + \textbf{apertura}};
\draw[fleche] (intro) -- (des);
\draw[fleche] (des) -- (conc);
\end{tikzpicture}
\end{center}
\legende{Organigramme des trois parties du commentaire de texte}
\end{popfigure}

\courssub{Formuler la problématique}

La \cle{problématique} est une question, formulée en espagnol, qui annonce l'axe de votre analyse. Elle commence souvent par \esp{¿Cómo...?} (comment), \esp{¿Qué...?} (que/quelle) ou \esp{¿Por qué...?} (pourquoi). Pour notre texte « Un planeta en peligro », la problématique peut être :

\esp{¿Cómo denuncia el autor la contaminación y qué soluciones propone?} « Comment l'auteur dénonce-t-il la pollution et quelles solutions propose-t-il ? »

Cette question a deux volets : le \textbf{comment} (les procédés : lexique, ton, impératifs) et le \textbf{quoi} (le contenu : la dénonciation et les solutions). Tout votre développement devra y répondre.

\begin{attention}[La ponctuation espagnole des questions]
En espagnol, toute question s'ouvre par \textbf{¿} et se ferme par \textbf{?} ; toute exclamation s'ouvre par \textbf{¡} et se ferme par \textbf{!}. Oublier le signe ouvrant est une faute pénalisée : écrivez \esp{¿Qué soluciones propone?} et \esp{¡Recicla!}, jamais \esp{Qué soluciones propone?} ni \esp{Recicla!}
\end{attention}

\courssub{Analyser les idées : problème, causes, solutions}

Premier travail du développement : relever dans le texte les trois niveaux d'idées et les citer.

\begin{itemize}
\item \textbf{Le problème} : \esp{El aire de las grandes ciudades está contaminado} — relevé du constat.
\item \textbf{Les causes} : \esp{las fábricas, los coches, la deforestación} — les responsables désignés.
\item \textbf{Les solutions} : \esp{debemos usar energías renovables y reciclar} — les propositions de l'auteur.
\end{itemize}

Chaque idée relevée doit être \textbf{citée entre guillemets} puis \textbf{traduite ou expliquée en français} si le correcteur le demande (ou simplement reformulée brièvement en espagnol). Citer sans commenter est une erreur ; commenter sans citer en est une autre.

\courssub{Analyser la langue : les procédés}

Deuxième travail du développement : montrer \emph{comment} l'auteur écrit. Voici les quatre procédés à repérer dans un texte de sensibilisation écologique.

\begin{propriete}[Les procédés d'un texte de sensibilisation]
\begin{enumerate}
\item \textbf{Le lexique péjoratif} : des mots négatifs qui dramatisent — \esp{peligro} (danger), \esp{destrucción} (destruction), \esp{amenaza} (menace), \esp{daño} (dommage), \esp{morir} (mourir).
\item \textbf{Les impératifs} : ils interpellent directement le lecteur — \esp{¡Recicla!}, \esp{¡Protege la naturaleza!}, \esp{¡Apaga la luz!}
\item \textbf{Les chiffres et les faits} : ils donnent un air scientifique et objectif au texte — \esp{millones de toneladas}, \esp{cada año}.
\item \textbf{Le ton} : souvent \cle{alarmiste} au début (pour inquiéter), puis \cle{optimiste} à la fin (pour donner envie d'agir) — \esp{Todavía podemos salvar el planeta.}
\end{enumerate}
\end{propriete}

\begin{exemplebox}[Phrases-modèles de commentaire, traduites]
\begin{itemize}
\item \esp{El texto trata de la contaminación del planeta.} « Le texte traite de la pollution de la planète. »
\item \esp{El autor denuncia la contaminación del aire.} « L'auteur dénonce la pollution de l'air. »
\item \esp{Según el autor, debemos proteger los bosques.} « Selon l'auteur, nous devons protéger les forêts. »
\item \esp{El tono del texto es alarmista al principio y optimista al final.} « Le ton du texte est alarmiste au début et optimiste à la fin. »
\item \esp{Este texto nos invita a actuar.} « Ce texte nous invite à agir. »
\item \esp{En conclusión, este texto nos invita a cuidar el medio ambiente.} « En conclusion, ce texte nous invite à prendre soin de l'environnement. »
\end{itemize}
\end{exemplebox}

\begin{attention}[Le verbe \esp{denunciar} et autres pièges]
\esp{Denunciar} signifie « dénoncer » au sens large : critiquer publiquement un problème, une injustice, un fléau. Il ne se limite pas à « dénoncer quelqu'un à la police ». Dans un commentaire, c'est un verbe en or : \esp{El autor denuncia la deforestación} « L'auteur dénonce la déforestation ». Autre piège : \esp{tratar de} signifie « traiter de » (thème) mais aussi « essayer de » selon le contexte ; dans \esp{El texto trata de...}, il s'agit toujours du thème. Enfin, \esp{actualmente} est un vrai ami : il signifie bien « actuellement ». Mais attention à \esp{una revista} : c'est une « revue, un magazine », et non une « révocation ».
\end{attention}

\courssub{Un commentaire modèle rédigé}

Voici un commentaire complet (environ 12 lignes) sur le texte « Un planeta en peligro » lu et annoté en classe. Observez les articulations : présentation, idées, langue, conclusion.

\begin{textebox}[Comentario modelo — « Un planeta en peligro »]
Este artículo, publicado en una revista ecológica, trata de la contaminación del planeta y de las medidas para proteger el medio ambiente. El autor plantea una pregunta esencial: ¿cómo podemos salvar la Tierra?

Primero, el texto describe el problema: el aire de las ciudades está contaminado, los ríos están sucios y los bosques desaparecen por la deforestación. Según el autor, las causas son las fábricas, los coches y el consumo excesivo. Después, propone soluciones concretas: usar energías renovables, reciclar y plantar árboles.

En cuanto a la lengua, el autor utiliza un léxico peyorativo: «peligro», «destrucción», «amenaza», que crea un tono alarmista. Sin embargo, al final, los imperativos «¡Recicla!», «¡Protege la naturaleza!» dan un tono optimista y llaman a la acción.

En conclusión, este texto nos invita a actuar para salvar el planeta. En mi opinión, es un artículo muy útil porque nos recuerda que todos somos responsables. ¿Seremos capaces de cambiar nuestros hábitos a tiempo?
\end{textebox}

\textbf{Glossaire} : \esp{plantear una pregunta} = poser une question ; \esp{peyorativo} = péjoratif ; \esp{llamar a la acción} = appeler à l'agir ; \esp{recordar} = rappeler ; \esp{ser capaz de} = être capable de ; \esp{a tiempo} = à temps.

Observez la mécanique : chaque procédé cité (« peligro », « destrucción ») est suivi de son \textbf{effet} (« crea un tono alarmista »). C'est cela, analyser : relier une forme à un effet de sens.

\begin{center}
\begin{tabularx}{\linewidth}{|X|X|X|}
\hline
\rowcolor{popblueL} \textbf{Français} & \textbf{Español} & \textbf{Remarque} \\
\hline
Le texte traite de... & El texto trata de... & Pour présenter le thème \\
\hline
L'auteur dénonce... & El autor denuncia... & Verbe clé du commentaire \\
\hline
Selon l'auteur... & Según el autor... & Suivi de l'indicatif \\
\hline
Le ton est alarmiste & El tono es alarmista & Analyse du ton \\
\hline
Ce texte nous invite à... & Este texto nos invita a... & + infinitif \\
\hline
En conclusion... & En conclusión... & Pour ouvrir le bilan \\
\hline
À mon avis... & En mi opinión... / A mi juicio... & Avis personnel \\
\hline
\end{tabularx}
\end{center}

\courssub{Entraînement méthodique}

\begin{exoresolu}[Construire une introduction de commentaire]
Énoncé : à partir des informations suivantes, rédigez en espagnol une introduction complète (3 phrases) : un article paru dans un journal camerounais hispanophone, sur la pollution de l'eau à Douala, qui propose des solutions. Utilisez : \esp{tratar de}, \esp{denunciar}, une problématique avec \esp{¿Cómo...?}.
\tcblower
Solution détaillée : on suit l'ordre type — type + source, thème, problématique.\\
\esp{Este artículo, publicado en un periódico de Douala, trata de la contaminación del agua en la ciudad. El autor denuncia esta situación y propone varias soluciones. ¿Cómo podemos proteger el agua, fuente de vida, en nuestras ciudades?}\\
Traduction : « Cet article, publié dans un journal de Douala, traite de la pollution de l'eau dans la ville. L'auteur dénonce cette situation et propose plusieurs solutions. Comment pouvons-nous protéger l'eau, source de vie, dans nos villes ? »\\
Points de réussite : verbe \esp{tratar de} + nom, verbe \esp{denunciar} + nom, problématique interrogative avec point d'interrogation ouvrant ¿ et fermant ?.
\end{exoresolu}

\begin{exercice}[À vous de commenter]
Relevez dans le commentaire modèle ci-dessus : 1) la phrase de présentation du texte ; 2) deux citations du lexique péjoratif ; 3) la phrase qui exprime l'avis personnel ; 4) la question d'ouverture. Puis rédigez en espagnol une conclusion de 3 phrases sur un texte traitant de la déforestation au Cameroun, avec \esp{En conclusión}, \esp{En mi opinión} et une question d'ouverture.
\end{exercice}

\begin{corrige}[Exercice — éléments de correction]
1) \esp{Este artículo, publicado en una revista ecológica, trata de la contaminación del planeta...} 2) « peligro », « destrucción », « amenaza ». 3) \esp{En mi opinión, es un artículo muy útil porque nos recuerda que todos somos responsables.} 4) \esp{¿Seremos capaces de cambiar nuestros hábitos a tiempo?}\\
Conclusion possible : \esp{En conclusión, este texto denuncia la deforestación de los bosques de Camerún y nos invita a plantar árboles. En mi opinión, es un texto importante porque los bosques son el pulmón del planeta. ¿Qué haremos mañana si no actuamos hoy?} « En conclusion, ce texte dénonce la déforestation des forêts du Cameroun et nous invite à planter des arbres. À mon avis, c'est un texte important car les forêts sont le poumon de la planète. Que ferons-nous demain si nous n'agissons pas aujourd'hui ? »
\end{corrige}

\begin{attention}[Les trois erreurs fatales du commentaire]
\begin{enumerate}
\item \textbf{Paraphraser sans analyser} : répéter le texte avec d'autres mots ne vaut rien. Chaque citation doit être suivie d'un commentaire sur son \emph{effet} (ton, émotion, conviction).
\item \textbf{Oublier les citations} : une analyse sans preuve textuelle est du bavardage. Règle d'or : une idée = une citation courte entre guillemets.
\item \textbf{Rédiger en français} : le commentaire se rédige entièrement \textbf{en espagnol}. Seule la traduction de citations peut être demandée en français, ponctuellement.
\end{enumerate}
\end{attention}

\begin{aretenir}[L'essentiel de la méthode]
Un commentaire = \textbf{introducción} (type, thème, source, problématique) + \textbf{desarrollo} (idées : problème, causes, solutions ; langue : lexique, impératifs, chiffres, ton) + \textbf{conclusión} (bilan, avis personnel, ouverture). On cite, on traduit ou explique, on analyse l'effet produit. Les phrases-modèles \esp{El texto trata de...}, \esp{El autor denuncia...}, \esp{Este texto nos invita a...} sont vos squelettes de rédaction : apprenez-les par cœur.
\end{aretenir}

Maintenant que vous savez analyser le texte d'un autre, il est temps de passer à la production : la section suivante vous apprendra à créer votre propre affiche de sensibilisation, avec slogans et impératifs.

\coursec{L'affiche de sensibilisation : impératif et slogans}

Dans la section précédente, nous avons appris à commenter un texte sur la pollution : repérer le thème, analyser les arguments, dégager l'intention de l'auteur. Mais comprendre un texte ne suffit pas : le programme de Terminale A demande aussi de \textbf{produire}. Or, dans la vraie vie, quand on veut défendre l'environnement, on n'écrit pas un commentaire de texte : on crée une \cle{affiche} (\esp{un cartel}), avec un slogan percutant et des consignes claires. Cette section vous donne les outils linguistiques pour cela : l'\textbf{impératif}, les \textbf{verbes d'appel à l'action} et l'art du \textbf{slogan}.

\courssub{Observer d'abord : une affiche réelle}

Avant toute règle, observons. Voici le texte d'une affiche imaginée pour une campagne de sensibilisation au lycée de Yaoundé.

\begin{textebox}[Cartel de la campaña « Vive Verde » — Lycée de Yaoundé]
\textbf{¡SALVA EL PLANETA!}\par
\emph{El planeta no es basura, cuídalo con cordura.}\par
¡Recicla el plástico y el papel!\par
¡No tires basura al río!\par
¡Ahorra agua y energía!\par
¡Planta un árbol cada año!\par
¡Cuidemos juntos nuestro barrio!\par
\emph{Asociación ecológica « Vive Verde » — Lycée de Yaoundé}
\end{textebox}

\textbf{Glossaire :} \esp{basura} = déchets, ordures ; \esp{cuídalo} = prends-en soin (\esp{cuidar} + pronom \esp{lo}) ; \esp{cordura} = bon sens, sagesse ; \esp{ahorrar} = économiser ; \esp{barrio} = quartier ; \esp{juntos} = ensemble.

\textbf{Questions d'observation :}
\begin{enumerate}
\item Combien de verbes à l'impératif comptez-vous sur cette affiche ?
\item Quelle phrase n'est \emph{pas} à l'impératif ? Pourquoi, à votre avis ?
\item Quel effet produit la rime \esp{basura} / \esp{cordura} ?
\end{enumerate}

\courssub{Les caractéristiques d'une affiche efficace}

\begin{definition}[El cartel de sensibilización]
Une affiche de sensibilisation (\esp{un cartel de sensibilización}) est un document visuel et bref qui cherche à \textbf{informer}, \textbf{alerter} et \textbf{faire agir} le public sur une cause (ici : la protection de l'environnement).
\end{definition}

Une bonne affiche comprend cinq éléments obligatoires :

\begin{enumerate}
\item \textbf{Un titre accrocheur} (\esp{un título llamativo}) : court, en grosses lettres, souvent à l'impératif : \esp{¡Salva el planeta!}
\item \textbf{Un slogan} (\esp{un eslogan}) : une phrase mémorable, souvent rimée ou rythmée.
\item \textbf{Trois à cinq consignes} à l'impératif : des actions concrètes que le lecteur peut faire.
\item \textbf{Un visuel} (\esp{una imagen}) : dessin, photo, symbole (arbre, planète, goutte d'eau).
\item \textbf{Une signature} (\esp{una firma}) : le nom de l'association, de l'école ou du club qui organise la campagne.
\end{enumerate}

\begin{attention}[Le slogan : court ou rien]
Un slogan efficace respecte trois critères : il est \textbf{court} (moins de 8 mots), il utilise l'\textbf{impératif} ou une \textbf{construction nominale} (sans verbe conjugué), et il est \textbf{mémorable} (rime, rythme, contraste). Évitez les phrases longues comme \esp{Es muy importante que todos nosotros protejamos el medio ambiente} : personne ne retient une phrase de douze mots. Préférez : \esp{Menos plástico, más vida.}
\end{attention}

\courssub{L'impératif affirmatif : donner des consignes}

\begin{propriete}[Formation de l'impératif affirmatif (tú)]
À la deuxième personne du singulier (\esp{tú}), l'impératif affirmatif est \textbf{identique à la 3\textsuperscript{e} personne du singulier du présent de l'indicatif} :
\begin{itemize}
\item \esp{proteger} → él \esp{protege} → \esp{¡Protege la naturaleza!} « Protège la nature ! »
\item \esp{reciclar} → él \esp{recicla} → \esp{¡Recicla el papel!} « Recycle le papier ! »
\item \esp{escribir} → él \esp{escribe} → \esp{¡Escribe un eslogan!} « Écris un slogan ! »
\end{itemize}
\end{propriete}

\begin{popfigure}
\begin{center}
\begin{tikzpicture}[node distance=0.5cm]
\node[draw=popblue, fill=popblueL, rounded corners, minimum width=4.2cm, minimum height=1.1cm, align=center] (ar) {\textbf{-ar : hablar}\\ \esp{¡habla!}};
\node[draw=popgreen, fill=popgreenL, rounded corners, minimum width=4.2cm, minimum height=1.1cm, align=center, right=0.8cm of ar] (er) {\textbf{-er : comer}\\ \esp{¡come!}};
\node[draw=poporange, fill=poporangeL, rounded corners, minimum width=4.2cm, minimum height=1.1cm, align=center, right=0.8cm of er] (ir) {\textbf{-ir : escribir}\\ \esp{¡escribe!}};
\node[draw=poppink, fill=poppinkL, rounded corners, minimum width=4.2cm, minimum height=1.1cm, align=center, below=0.7cm of ar] (arn) {\esp{¡no hables!}};
\node[draw=poppink, fill=poppinkL, rounded corners, minimum width=4.2cm, minimum height=1.1cm, align=center, below=0.7cm of er] (ern) {\esp{¡no comas!}};
\node[draw=poppink, fill=poppinkL, rounded corners, minimum width=4.2cm, minimum height=1.1cm, align=center, below=0.7cm of ir] (irn) {\esp{¡no escribas!}};
\draw[-{Stealth}, thick, popmuted] (ar) -- (arn);
\draw[-{Stealth}, thick, popmuted] (er) -- (ern);
\draw[-{Stealth}, thick, popmuted] (ir) -- (irn);
\node[align=center, text=popdark, font=\small, above=0.15cm of ar, yshift=0.35cm] {\textbf{Affirmatif : 3\textsuperscript{e} pers. sg. du présent} \qquad \textbf{Négatif : no + subjonctif}};
\end{tikzpicture}
\end{center}
\legende{Formation de l'impératif à \esp{tú} : affirmatif (en haut) et négatif (en bas).}
\end{popfigure}

\begin{attention}[Les irréguliers de l'impératif affirmatif]
Huit verbes très fréquents ont un impératif \esp{tú} irrégulier, à apprendre par cœur :
\begin{center}
\begin{tabular}{|l|l|l|}
\hline
\rowcolor{popblueL} \textbf{Infinitivo} & \textbf{Imperativo} & \textbf{Exemple} \\
\hline
\esp{hacer} & \esp{¡haz!} & \esp{¡Haz un cartel!} « Fais une affiche ! » \\
\hline
\esp{ir} & \esp{¡ve!} & \esp{¡Ve a la manifestación!} « Va à la manifestation ! » \\
\hline
\esp{tener} & \esp{¡ten!} & \esp{¡Ten cuidado!} « Fais attention ! » \\
\hline
\esp{ser} & \esp{¡sé!} & \esp{¡Sé responsable!} « Sois responsable ! » \\
\hline
\esp{poner} & \esp{¡pon!} & \esp{¡Pon la basura en el cubo!} « Mets les déchets à la poubelle ! » \\
\hline
\esp{salir} & \esp{¡sal!} & \esp{¡Sal a limpiar el barrio!} « Sors nettoyer le quartier ! » \\
\hline
\esp{decir} & \esp{¡di!} & \esp{¡Di no al plástico!} « Dis non au plastique ! » \\
\hline
\esp{venir} & \esp{¡ven!} & \esp{¡Ven a plantar árboles!} « Viens planter des arbres ! » \\
\hline
\end{tabular}
\end{center}
Piège classique du francophone : dire \esp{¡hace!} au lieu de \esp{¡haz!}. La forme régulière n'existe pas ici. Autre piège : \esp{¡sé!} (impératif de \esp{ser}) prend un accent pour ne pas être confondu avec le pronom \esp{se} ; de même \esp{¡ve!} (impératif de \esp{ir}) se distingue de \esp{ve} (« il voit », présent de \esp{ver}).
\end{attention}

\begin{propriete}[Place des pronoms avec l'impératif]
À l'impératif \textbf{affirmatif}, les pronoms se placent \textbf{après} le verbe et s'y \textbf{attachent} en un seul mot, ce qui provoque souvent un accent graphique : \esp{cuidar} + \esp{lo} → \esp{¡Cuídalo!} « Prends-en soin ! » ; \esp{hacer} + \esp{lo} → \esp{¡Hazlo!} « Fais-le ! » ; \esp{plantar} + \esp{los} → \esp{¡Plántalos!} « Plante-les ! »\par
À l'impératif \textbf{négatif}, les pronoms se placent \textbf{avant} le verbe, détachés : \esp{¡No lo tires!} « Ne le jette pas ! » ; \esp{¡No los cortes!} « Ne les coupe pas ! »
\end{propriete}

\courssub{L'impératif négatif : interdire poliment}

\begin{propriete}[Formation de l'impératif négatif]
L'impératif négatif se forme avec \esp{no} + le \textbf{subjonctif présent} :
\begin{itemize}
\item \esp{contaminar} → \esp{¡No contamines el aire!} « Ne pollue pas l'air ! »
\item \esp{tirar} → \esp{¡No tires basura al río!} « Ne jette pas de déchets dans la rivière ! »
\item \esp{desperdiciar} → \esp{¡No desperdicies el agua!} « Ne gaspille pas l'eau ! »
\item \esp{cortar} → \esp{¡No cortes los árboles!} « Ne coupe pas les arbres ! »
\end{itemize}
Rappel de formation du subjonctif : les verbes en \textbf{-ar} prennent les voyelles en \textbf{-e} (\esp{hables}), les verbes en \textbf{-er/-ir} prennent les voyelles en \textbf{-a} (\esp{comas}, \esp{escribas}). Les verbes à changement de radical le conservent : \esp{pensar} → \esp{¡no pienses!}, \esp{dormir} → \esp{¡no duermas!}, \esp{pedir} → \esp{¡no pidas!}.
\end{propriete}

\begin{exemplebox}[Affirmatif et négatif en regard]
\begin{itemize}
\item \esp{¡Protege la capa de ozono!} « Protège la couche d'ozone ! » / \esp{¡No destruyas la naturaleza!} « Ne détruis pas la nature ! »
\item \esp{¡Ahorra energía!} « Économise l'énergie ! » / \esp{¡No malgastes la luz!} « Ne gaspille pas la lumière ! »
\item \esp{¡Usa el transporte público!} « Utilise les transports en commun ! » / \esp{¡No uses el coche cada día!} « N'utilise pas la voiture tous les jours ! »
\end{itemize}
\end{exemplebox}

\courssub{Les formes collectives : agir ensemble}

Une campagne écologique s'adresse rarement à une seule personne : elle mobilise un groupe. Deux formes sont alors utiles.

\begin{propriete}[Impératif collectif]
\begin{itemize}
\item \textbf{\esp{nosotros}} (exhortation, « nous ») : on utilise le \textbf{subjonctif présent} : \esp{¡Cuidemos el planeta!} « Prenons soin de la planète ! » ; \esp{¡Luchemos contra la deforestación!} « Luttons contre la déforestation ! »
\item \textbf{\esp{vosotros}} (Espagne uniquement) : infinitif sans \textbf{-r} final + \textbf{-d} : \esp{¡Reciclad!} « Recyclez ! » ; \esp{¡Respetad la naturaleza!} « Respectez la nature ! » ; au négatif, \esp{no} + subjonctif : \esp{¡No contaminéis!} « Ne polluez pas ! »
\end{itemize}
En Amérique latine et en Guinée équatoriale, on préfère \esp{ustedes} + subjonctif : \esp{¡Reciclen!} « Recyclez ! » ; \esp{¡No tiren basura!} « Ne jetez pas de déchets ! »
\end{propriete}

\begin{exemplebox}[Exemples traduits]
\esp{¡No tires basura al río!} « Ne jette pas de déchets dans la rivière ! »\par
\esp{¡Cuidemos nuestro planeta!} « Prenons soin de notre planète ! »\par
\esp{¡Plantemos árboles en el colegio!} « Plantons des arbres au lycée ! »\par
\esp{¡Limpiad las playas!} « Nettoyez les plages ! » (Espagne, \esp{vosotros})
\end{exemplebox}

\courssub{Les verbes d'appel à l'action}

Voici le lexique essentiel des campagnes de sensibilisation, à conjuguer à l'impératif.

\begin{center}
\begin{tabularx}{\linewidth}{|X|X|X|}
\hline
\rowcolor{popblueL} \textbf{Verbo} & \textbf{Français} & \textbf{Ejemplo (imperativo tú)} \\
\hline
\esp{proteger} & protéger & \esp{¡Protege los bosques!} \\
\hline
\esp{preservar} & préserver & \esp{¡Preserva el agua potable!} \\
\hline
\esp{cuidar} & prendre soin de & \esp{¡Cuida tu barrio!} \\
\hline
\esp{salvar} & sauver & \esp{¡Salva el planeta!} \\
\hline
\esp{reciclar} & recycler & \esp{¡Recicla el plástico!} \\
\hline
\esp{reutilizar} & réutiliser & \esp{¡Reutiliza las bolsas!} \\
\hline
\esp{reducir} & réduire & \esp{¡Reduce el consumo!} \\
\hline
\esp{ahorrar} & économiser & \esp{¡Ahorra energía!} \\
\hline
\esp{plantar} & planter & \esp{¡Planta un árbol!} \\
\hline
\esp{limpiar} & nettoyer & \esp{¡Limpia la playa!} \\
\hline
\esp{respetar} & respecter & \esp{¡Respeta la naturaleza!} \\
\hline
\end{tabularx}
\end{center}

\begin{attention}[Faux amis et pièges lexicaux]
\begin{itemize}
\item \esp{salvar} signifie « sauver », jamais « saluer » : \esp{¡Salva los océanos!} « Sauve les océans ! »
\item \esp{ahorrar} = économiser (argent, eau, énergie) ; ne pas confondre avec \esp{ahogar} (étouffer, noyer).
\item \esp{reducir} a un impératif régulier : \esp{¡reduce!} (et non *\esp{¡reduz!}).
\item \esp{tirar basura} = jeter des ordures ; le mot \esp{basura} s'emploie au singulier dans les slogans : \esp{¡No tires basura!}
\item Attention à l'orthographe de \esp{¡protege!} : le verbe \esp{proteger} garde le \textbf{g} à l'impératif affirmatif, mais au négatif le subjonctif exige \textbf{j} : \esp{¡No protejas solo tu casa!} « Ne protège pas seulement ta maison ! »
\end{itemize}
\end{attention}

\courssub{Construire un slogan : la recette}

\begin{methode}[Fabriquer un slogan en quatre étapes]
\begin{enumerate}
\item \textbf{Choisir un verbe d'action} et le mettre à l'impératif : \esp{salvar}, \esp{reciclar}, \esp{plantar}.
\item \textbf{Ajouter un complément court} (2 à 4 mots) : \esp{¡Salva el planeta!}
\item \textbf{Ajouter du rythme, une rime ou un contraste} si possible : \esp{¡Recicla hoy para vivir mañana!} repose sur le parallélisme temporel \esp{hoy} / \esp{mañana}.
\item \textbf{Vérifier la longueur} : moins de 8 mots. Sinon, couper.
\end{enumerate}
\end{methode}

Les slogans hispaniques réels utilisent souvent trois procédés :

\begin{itemize}
\item \textbf{La rime} : \esp{El planeta no es basura, cuídalo con cordura.} « La planète n'est pas une poubelle, prends-en soin avec bon sens. »
\item \textbf{Le contraste} : \esp{Menos plástico, más vida.} « Moins de plastique, plus de vie. » (construction nominale, sans verbe)
\item \textbf{Le parallélisme temporel} : \esp{¡Recicla hoy para vivir mañana!} « Recycle aujourd'hui pour vivre demain ! »
\end{itemize}

\begin{savaistu}[Les 3 R : un slogan planétaire]
\esp{Reduce, reutiliza, recicla} — « Réduis, réutilise, recycle » — est le slogan international des \textbf{3 R}, adopté par les campagnes écologiques de tout le monde hispanophone, de Madrid à Mexico et jusqu'à Malabo. Sa force ? Trois verbes à l'impératif, trois mots, une allitération en \textbf{r} : impossible à oublier. C'est un modèle parfait de ce qu'on vous demande de produire à l'examen.
\end{savaistu}

\courssub{Maquette d'affiche : organiser l'espace}

\begin{popfigure}
\begin{center}
\begin{tikzpicture}
\draw[draw=popdark, very thick, rounded corners] (0,0) rectangle (9,10);
\draw[draw=popblue, thick, fill=popblueL, rounded corners] (0.4,8.4) rectangle (8.6,9.6);
\node[align=center, text=popdark] at (4.5,9) {\textbf{TÍTULO}\\ \esp{¡SALVA EL PLANETA!}};
\draw[draw=poppurple, thick, fill=poppurpleL, rounded corners] (0.4,7.2) rectangle (8.6,8.1);
\node[align=center, text=popdark] at (4.5,7.65) {\textbf{ESLOGAN}\\ \esp{El planeta no es basura, cuídalo con cordura}};
\draw[draw=popgreen, thick, fill=popgreenL, rounded corners] (0.4,3.2) rectangle (5.2,6.8);
\node[align=center, text=popdark] at (2.8,5.9) {\textbf{CONSIGNAS}};
\node[align=left, text=popdark] at (2.8,4.8) {\esp{¡Recicla el papel!}\\ \esp{¡No tires basura!}\\ \esp{¡Planta un árbol!}};
\draw[draw=poporange, thick, fill=poporangeL, rounded corners] (5.6,3.2) rectangle (8.6,6.8);
\node[align=center, text=popdark] at (7.1,5) {\textbf{IMAGEN}\\ (un árbol,\\ la Tierra...)};
\draw[draw=popgold, thick, fill=popgoldL, rounded corners] (0.4,2) rectangle (8.6,2.8);
\node[align=center, text=popdark] at (4.5,2.4) {\textbf{FIRMA :} \esp{Asociación « Vive Verde »}};
\end{tikzpicture}
\end{center}
\legende{Maquette simplifiée d'une affiche de sensibilisation : cinq zones obligatoires.}
\end{popfigure}

\courssub{Entraînement}

\begin{exoresolu}[Transformer des consignes en slogans]
Énoncé : transformez les phrases suivantes en consignes d'affiche à l'impératif (\esp{tú}), affirmatif ou négatif selon le sens.
\begin{enumerate}
\item Il faut économiser l'eau. (\esp{ahorrar})
\item Il ne faut pas jeter de déchets par terre. (\esp{tirar})
\item Il faut utiliser les énergies renouvelables. (\esp{usar})
\item Il ne faut pas couper les forêts. (\esp{cortar})
\end{enumerate}
\tcblower
Solution détaillée :
\begin{enumerate}
\item \esp{¡Ahorra agua!} — impératif affirmatif = 3\textsuperscript{e} pers. sg. du présent (\esp{él ahorra}).
\item \esp{¡No tires basura al suelo!} — impératif négatif = \esp{no} + subjonctif (\esp{tires}, verbe en -ar → voyelle \textbf{e}).
\item \esp{¡Usa energías renovables!} — affirmatif régulier.
\item \esp{¡No cortes los bosques!} — négatif : subjonctif \esp{cortes}.
\end{enumerate}
Attention : pour la phrase 2, on ne dit jamais *\esp{¡No tira!} : la négation exige le subjonctif.
\end{exoresolu}

\begin{exercice}[À vous de créer]
Rédigez trois slogans pour une campagne contre la déforestation au Cameroun : un avec un impératif affirmatif, un avec un impératif négatif, un avec une forme collective en \esp{nosotros}. Puis esquissez la maquette de votre affiche en respectant les cinq zones vues plus haut.
\end{exercice}

\begin{corrige}[Exercice « À vous de créer »]
Éléments de réponse possibles (toute formulation correcte est acceptée) :
\begin{itemize}
\item Affirmatif : \esp{¡Planta un árbol cada año!} « Plante un arbre chaque année ! »
\item Négatif : \esp{¡No destruyas los bosques!} « Ne détruis pas les forêts ! »
\item Collectif : \esp{¡Protejamos juntos nuestros bosques!} « Protégeons ensemble nos forêts ! »
\end{itemize}
La maquette doit comporter : un titre accrocheur en haut, un slogan court, 3 à 5 consignes à l'impératif, un emplacement pour le visuel et la signature du club écologique du lycée.
\end{corrige}

\begin{aretenir}[L'essentiel de la section]
\begin{itemize}
\item Une affiche = titre + slogan + 3 à 5 consignes à l'impératif + visuel + signature.
\item Impératif affirmatif (\esp{tú}) = 3\textsuperscript{e} pers. sg. du présent : \esp{¡recicla!}, \esp{¡protege!}
\item Impératif négatif = \esp{no} + subjonctif : \esp{¡no contamines!}, \esp{¡no tires!}
\item Pronoms : attachés après à l'affirmatif (\esp{¡cuídalo!}), détachés avant au négatif (\esp{¡no lo tires!}).
\item Collectif : \esp{¡cuidemos!} (nosotros), \esp{¡reciclad!} (vosotros, Espagne), \esp{¡reciclen!} (ustedes, Amérique latine).
\item Irréguliers à mémoriser : \esp{haz, ve, ten, sé, pon, sal, di, ven}.
\item Un bon slogan : moins de 8 mots, rime ou contraste, mémorable : \esp{Reduce, reutiliza, recicla}.
\end{itemize}
\end{aretenir}

\coursec{Production guidée : créer son affiche}

Après avoir étudié la structure de l'affiche de sensibilisation, l'emploi de l'impératif et les techniques du slogan (\emph{cf. section précédente}), nous passons à la réalisation concrète. Cette production guidée vous accompagne pas à pas, de l'idée au produit final, en mobilisant le lexique de la leçon (\cle{contaminación}, \cle{calentamiento global}, \cle{deforestación}, \cle{sequía}, \cle{capa de ozono}, \cle{energías renovables}) et les formes grammaticales vues (impératif affirmatif/négatif, accord des adjectifs).

\courssub{Étape 1 : choisir un problème environnemental local ou global}

\begin{textebox}[Contexte camerounais : trois pistes pour votre affiche]
En Cameroun, les défis environnementaux sont quotidiens. Choisissez \textbf{un} seul thème pour que le message reste percutant.

\textbf{Piste A — Pollution de l'air à Douala :} Les émissions des véhicules anciens, la combustion des déchets à ciel ouvert et les activités industrielles créent un brouillard persistant. Mots-clés : \esp{\cle{contaminación} del aire}, \esp{humo}, \esp{vehículos viejos}, \esp{quema de basura}, \esp{enfermedades respiratorias}.

\textbf{Piste B — Plastiques à Yaoundé :} Les sachets \emph{« Callibox »} et bouteilles bouchent les caniveaux, provoquent des inondations et finissent dans le fleuve Wouri. Mots-clés : \esp{\cle{contaminación} por plásticos}, \esp{bolsas de plástico}, \esp{inundaciones}, \esp{reciclaje}, \esp{economía circular}.

\textbf{Piste C — Déforestation et eau potable :} L'exploitation forestière illégale et le bois de chauffe menacent les bassins versants qui alimentent Yaoundé et Douala. Mots-clés : \esp{\cle{deforestación}}, \cle{sequía}, \esp{cuencas hidrográficas}, \esp{agua potable}, \esp{reforestación}.
\tcblower
\textbf{Glossaire :} \esp{humo} = fumée ; \esp{quema de basura} = brûlage d'ordures ; \esp{enfermedades respiratorias} = maladies respiratoires ; \esp{bolsas de plástico} = sacs en plastique ; \esp{cuencas hidrográficas} = bassins versants ; \esp{reforestación} = reboisement.
\end{textebox}

\begin{methode}[Comment choisir son sujet ?]
\begin{enumerate}
\item \textbf{Listez} 3 problèmes visibles dans votre quartier (déchets, fumée, eau trouble, arbres coupés).
\item \textbf{Sélectionnez} celui qui vous touche le plus et pour lequel vous connaissez des solutions simples.
\item \textbf{Notez} 5 mots du lexique de la leçon qui s'y rattachent (ex. pour Piste B : \esp{plástico}, \esp{reciclar}, \esp{basura}, \esp{mar}, \esp{futuro}).
\item \textbf{Validez} avec votre binôme : le sujet est-il clair ? Le vocabulaire suffit-il pour 3 consignes ?
\end{enumerate}
\end{methode}

\courssub{Étape 2 : rédiger un titre et un slogan percutants}

Le titre nomme le problème ; le slogan (phrase courte, rythmée, mémorable) appelle à l'action. En espagnol, le slogan utilise souvent l'impératif, l'infinitif d'affiche ou une phrase nominale.

\begin{exemplebox}[Titres et slogans — modèles traduits]
\begin{itemize}
\item \esp{¡Douala respira! — ¡Alto a la \cle{contaminación} del aire!} \quad (« Douala respire ! — Halte à la pollution de l'air ! »)
\item \esp{Yaoundé sin plástico — ¡Tu bolsa, tu responsabilidad!} \quad (« Yaoundé sans plastique — Ton sac, ta responsabilité ! »)
\item \esp{Árboles para la vida — ¡Planta uno hoy!} \quad (« Des arbres pour la vie — Plante-en un aujourd'hui ! »)
\item \esp{Agua es vida — ¡No la desperdicies!} \quad (« L'eau est vie — Ne la gaspille pas ! »)
\item \esp{\cle{Energías renovables} para Camerún — ¡El sol nos ilumina!} \quad (« Énergies renouvelables pour le Cameroun — Le soleil nous illumine ! »)
\end{itemize}
\end{exemplebox}

\begin{propriete}[Règles d'or du slogan efficace]
\begin{itemize}
\item \textbf{Brièveté :} 5 à 8 mots maximum.
\item \textbf{Rythme :} allitérations (\esp{plástico / planeta}), rimes (\esp{agua / futura}), parallélisme (\esp{¡Planta un árbol, regala vida!}).
\item \textbf{Vocabulaire de la leçon :} au moins \textbf{3 mots} parmi \cle{contaminación}, \cle{calentamiento global}, \cle{deforestación}, \cle{sequía}, \cle{capa de ozono}, \cle{energías renovables} ou leurs dérivés (\esp{contaminar}, \esp{deforestador}, \esp{renovable}).
\item \textbf{Accord des adjectifs :} \esp{energía limpia} (fém. sg.), \esp{planeta verde} (masc. sg.), \esp{aguas potables} (fém. pl.), \esp{bosques protegidos} (masc. pl.).
\end{itemize}
\end{propriete}

\begin{attention}[Piège fréquent : l'accord des adjectifs dans les slogans]
En français, l'adjectif s'accorde souvent à l'oral comme à l'écrit. En espagnol, l'accord est \textbf{obligatoire et visible} à l'écrit.
\begin{center}
\begin{tabularx}{\linewidth}{|X|X|X|}
\hline
\rowcolor{popblueL} \textbf{Français} & \textbf{Espagnol correct} & \textbf{Erreur fréquente} \\
\hline
Énergie propre & \esp{energía limpia} & \esp{energía limpio} \\
Planète verte & \esp{planeta verde} & \esp{planeta verda} (genre incorrect) \\
Eau potable & \esp{agua potable} & \esp{agua potables} (nombre incorrect) \\
Forêts protégées & \esp{bosques protegidos} & \esp{bosques protegida} (genre/nombre incorrects) \\
\hline
\end{tabularx}
\end{center}
\emph{Règle :} l'adjectif s'accorde en \textbf{genre} (masculin/féminin) et en \textbf{nombre} (singulier/pluriel) avec le nom qu'il qualifie. \esp{Agua} est féminin (même si \esp{el agua} au singulier pour l'euphonie) → \esp{agua limpia}, \esp{aguas limpias}. Les adjectifs en \emph{-e} (\esp{verde}, \esp{potable}, \esp{grande}) ne varient pas en genre, seulement en nombre (\esp{verde} / \esp{verdes}).
\end{attention}

\courssub{Étape 3 : écrire 3 à 5 consignes à l'impératif (affirmatif et négatif)}

L'impératif est le mode de l'ordre, du conseil, de l'invitation. Sur une affiche, on alterne \textbf{affirmatif} (faites ceci) et \textbf{négatif} (ne faites pas cela) pour varier le ton.

\begin{propriete}[Rappel : formation de l'impératif (tú / vosotros / ustedes)]
\begin{center}
\begin{tabular}{|l|c|c|c|}
\hline
\rowcolor{popblueL} \textbf{Verbe} & \textbf{Tú (affirmatif)} & \textbf{Tú (négatif)} & \textbf{Ustedes (affirmatif / négatif)} \\
\hline
\esp{hablar} (ar) & \esp{habla} & \esp{no hables} & \esp{hablen / no hablen} \\
\esp{comer} (er) & \esp{come} & \esp{no comas} & \esp{coman / no coman} \\
\esp{vivir} (ir) & \esp{vive} & \esp{no vivas} & \esp{vivan / no vivan} \\
\hline
\rowcolor{popblueL} \textbf{Irréguliers (tú affirmatif)} & \textbf{Forme} & \textbf{Négatif (subjonctif)} & \textbf{Ustedes} \\
\hline
\esp{ir} & \esp{ve} & \esp{no vayas} & \esp{vayan / no vayan} \\
\esp{ser} & \esp{sé} & \esp{no seas} & \esp{sean / no sean} \\
\esp{estar} & \esp{está} & \esp{no estés} & \esp{estén / no estén} \\
\esp{tener} & \esp{ten} & \esp{no tengas} & \esp{tengan / no tengan} \\
\esp{poner} & \esp{pon} & \esp{no pongas} & \esp{pongan / no pongan} \\
\esp{hacer} & \esp{haz} & \esp{no hagas} & \esp{hagan / no hagan} \\
\esp{decir} & \esp{di} & \esp{no digas} & \esp{digan / no digan} \\
\esp{salir} & \esp{sal} & \esp{no salgas} & \esp{salgan / no salgan} \\
\hline
\end{tabular}
\end{center}
\emph{Choix du registre :} \textbf{Tú} = proximité, jeunesse, ton direct (\esp{¡Recicla!}). \textbf{Ustedes} = public large, formel, institutionnel (\esp{¡Reciclen!}). Pour une affiche scolaire, \textbf{tú} est souvent plus engageant.
\end{propriete}

\begin{exemplebox}[Consignes types pour chaque piste — Tú affirmatif / négatif]
\textbf{Piste A (Air) :}
\begin{itemize}
\item \esp{¡Usa transporte público o comparte coche!} (affirmatif)
\item \esp{¡No quemes basura en la calle!} (négatif)
\item \esp{¡Exige control de emisiones!} (affirmatif)
\item \esp{¡Planta árboles en tu barrio!} (affirmatif)
\item \esp{¡No uses moto vieja sin catalizador!} (négatif)
\end{itemize}

\textbf{Piste B (Plastique) :}
\begin{itemize}
\item \esp{¡Lleva tu propia bolsa al mercado!} (affirmatif)
\item \esp{¡No tires botellas en la calle!} (négatif)
\item \esp{¡Separa los residuos en casa!} (affirmatif)
\item \esp{¡Recicla el plástico!} (affirmatif)
\item \esp{¡No compres productos con exceso de envoltorio!} (négatif)
\end{itemize}

\textbf{Piste C (Eau / Arbres) :}
\begin{itemize}
\item \esp{¡Cuida el agua, cierra el grifo!} (affirmatif)
\item \esp{¡No talar árboles sin permiso!} (négatif — infinitif d'affiche)
\item \esp{¡Denuncia la tala ilegal!} (affirmatif)
\item \esp{¡Participa en la reforestación!} (affirmatif)
\item \esp{¡No contamines los ríos!} (négatif)
\end{itemize}
\end{exemplebox}

\begin{exoresolu}[Construire ses 4 consignes pour l'affiche « Yaoundé sin plástico »]
\textbf{Consigne :} Rédigez 4 consignes (2 affirmatives, 2 négatives) à l'impératif \textbf{tú}, en utilisant au moins 3 mots du lexique de la leçon.

\textbf{Démarche :}
\begin{enumerate}
\item \textbf{Ciblez les verbes d'action :} \esp{llevar}, \esp{tirar}, \esp{separar}, \esp{reciclar}, \esp{comprar}, \esp{reutilizar}, \esp{rechazar}.
\item \textbf{Conjuguez à l'impératif tú :} \esp{lleva}, \esp{no tires}, \esp{separa}, \esp{recicla}, \esp{no compres}, \esp{reutiliza}, \esp{rechaza}.
\item \textbf{Intégrez le lexique :} \esp{bolsas de plástico}, \esp{reciclaje}, \esp{basura}, \esp{medio ambiente}, \esp{futuro}.
\item \textbf{Vérifiez l'orthographe et les accents.}
\end{enumerate}

\textbf{Production finale :}
\begin{enumerate}
\item \esp{¡Lleva tu bolsa de tela al mercado!} (affirmatif — \cle{contaminación} évitée)
\item \esp{¡No tires plástico en las cunetas!} (négatif — \cle{contaminación} directe)
\item \esp{¡Separa y recicla tus residuos!} (affirmatif — \esp{recicla} dérivé de \cle{energías renovables} / économie circulaire)
\item \esp{¡Rechaza los envoltorios innecesarios!} (affirmatif — vocabulaire associé)
\end{enumerate}
\tcblower
\textbf{Analyse :} 4 consignes variées, 2 négatives (\esp{no tires}, \esp{rechaza} = sens négatif), lexique présent (\esp{plástico}, \esp{residuos}, \esp{recicla}), accents corrects (\esp{lleva} pas \esp{lléva} ; \esp{tires} pas \esp{tíres} ; \esp{recicla} pas \esp{recícla}).
\end{exoresolu}

\courssub{Étape 4 : ajouter la phrase d'appel final}

La phrase de clôture ancre l'affiche dans l'urgence collective. Le modèle imposé : \esp{« ¡Actúa ahora! El planeta te necesita. »} (« Agis maintenant ! La planète a besoin de toi. »).

\begin{aretenir}[Variantes autorisées pour l'oral du Bac]
Vous pouvez adapter la seconde phrase si elle reste courte et mobilisatrice :
\begin{itemize}
\item \esp{¡Actúa ahora! Tu futuro está en juego.}
\item \esp{¡Actúa ahora! Camerún te mira.}
\item \esp{¡Actúa ahora! Cada gesto cuenta.}
\end{itemize}
\emph{Structure :} Impératif \esp{actúa} (tú de \esp{actuar}) + adverbe \esp{ahora} + phrase déclarative courte (sujet + verbe + complément).
\end{aretenir}

\courssub{Étape 5 : mise en page, signature et présentation orale (Bac)}

\begin{methode}[Checklist avant de rendre l'affiche — \textbf{à cocher en binôme}]
\begin{enumerate}
\item \textbf{Grand titre} en haut, lisible à 3 m (majuscules, police sans empattement, couleur contrastée).
\item \textbf{Slogan} central, en couleur, encadré ou en bulle.
\item \textbf{3 à 5 consignes} numérotées ou à puces, impératifs \textbf{tú} corrects (accents, négation \esp{no} + subjonctif).
\item \textbf{Au moins 3 mots} du lexique de la leçon (\cle{contaminación}, \cle{calentamiento global}, \cle{deforestación}, \cle{sequía}, \cle{capa de ozono}, \cle{energías renovables} ou dérivés).
\item \textbf{Phrase finale :} \esp{¡Actúa ahora! El planeta te necesita.}
\item \textbf{Illustration / dessin} fait main ou image libre de droit, en lien direct avec le thème.
\item \textbf{Signature du groupe :} \esp{Grupo 3 — 2nde A — Lycée de Nkolbisson — 2025}.
\item \textbf{Relecture croisée :} orthographe, accents, accords (adjectifs, participes), ponctuation espagnole (¡ ! ¿ ?).
\item \textbf{Photo numérique} pour l'oral (format paysage, bonne lumière).
\end{enumerate}
\end{methode}

\begin{savaistu}[Le savais-tu ? — « Clean City » et l'espagnol au Cameroun]
À Douala, la campagne \emph{« Clean City »} (ville propre) affiche des consignes en \textbf{français, anglais, bassa, douala}… et parfois en \textbf{espagnol} ! La communauté hispanophone (Guinée équatoriale voisine, étudiants, commerçants) est visible aux marchés Central et Congo. Une affiche bien faite en espagnol, avec un slogan clair (\esp{¡Douala limpia, Douala viva!}) et des pictogrammes, pourrait être affichée sur les panneaux municipaux. Votre travail de classe a donc une portée citoyenne réelle !
\end{savaistu}

\courssub{Préparation à l'oral du Bac : présenter son affiche en 5 phrases}

\begin{methode}[Structure type de la présentation orale (1 min 30 max)]
\begin{enumerate}
\item \textbf{Accroche + thème :} \esp{Buenos días. Nuestro cartel denuncia la contaminación por plásticos en Yaoundé.}
\item \textbf{Problème localisé :} \esp{Las bolsas y botellas tapan los desagües y causan inundaciones en los barrios como Mvog-Ada.}
\item \textbf{Solutions proposées (impératifs) :} \esp{Propone tres acciones: llevar tu bolsa, separar la basura y reciclar.}
\item \textbf{Slogan + appel :} \esp{El lema es: « ¡Yaoundé sin plástico, tu bolsa, tu responsabilidad! » y termina con « ¡Actúa ahora! El planeta te necesita. »}
\item \textbf{Engagement personnel :} \esp{Nosotros, el Grupo 3, nos comprometemos a reciclar y a sensibilizar a nuestras familias.}
\end{enumerate}
\end{methode}

\begin{exemplebox}[Modèle complet d'affiche — Texte seul (mise en page à faire sur papier A3)]
\textbf{TÍTULO :} \esp{¡YAOUNDÉ SIN PLÁSTICO!}

\textbf{SLOGAN :} \esp{¡Tu bolsa, tu responsabilidad!}

\textbf{CONSIGNAS :}
\begin{enumerate}
\item \esp{¡Lleva tu bolsa de tela al mercado!}
\item \esp{¡No tires botellas en las cunetas!}
\item \esp{¡Separa y recicla tus residuos!}
\item \esp{¡Rechaza los envoltorios innecesarios!}
\end{enumerate}

\textbf{LLAMADA FINAL :} \esp{¡Actúa ahora! El planeta te necesita.}

\textbf{FIRMA :} \esp{Grupo 4 — Tle A — Lycée Bilingue d'Essos — Mayo 2025}
\end{exemplebox}

\begin{center}
\begin{popfigure}
\begin{tikzpicture}[
  node distance=1.2cm and 1.5cm,
  stepnode/.style={rectangle, rounded corners, draw=popblue, thick, fill=popblueL, text width=2.8cm, align=center, font=\small\bfseries, minimum height=1.6cm},
  arrow/.style={-Stealth, thick, popdark}
]
\node[stepnode, align=center] (tema){1. TEMA\\Elección del problema};
\node[stepnode, right=of tema, align=center] (slogan){2. SLOGAN\\Título + frase gancho};
\node[stepnode, right=of slogan, align=center] (consignas){3. CONSIGNAS\\3–5 imperativos};
\node[stepnode, right=of consignas, align=center] (llamada){4. LLAMADA FINAL\\¡Actúa ahora!};
\node[stepnode, right=of llamada, align=center] (presentacion){5. PRESENTACIÓN\\Cartel + Oral};

\draw[arrow] (tema) -- (slogan);
\draw[arrow] (slogan) -- (consignas);
\draw[arrow] (consignas) -- (llamada);
\draw[arrow] (llamada) -- (presentacion);

\foreach \n/\l in {tema/T, slogan/S, consignas/C, llamada/L, presentacion/P} {
  \node[circle, draw=poporange, fill=poporangeL, thick, minimum size=0.7cm, font=\tiny\bfseries, anchor=south] at (\\.north) {\l};
}
\end{tikzpicture}
\legende{Les 5 étapes de la production de l'affiche : du choix du thème à la présentation orale.}
\end{popfigure}
\end{center}

\courssub{Grille d'auto-évaluation (à remplir avant le dépôt)}

\begin{center}
\begin{tabularx}{\linewidth}{|X|c|c|c|}
\hline
\rowcolor{popblueL} \textbf{Critère} & \textbf{Non acquis (0)} & \textbf{En cours (1)} & \textbf{Acquis (2)} \\
\hline
Lexique de la leçon (≥ 3 mots corrects) & 0–1 mot & 2 mots & ≥ 3 mots bien employés \\
\hline
Impératifs (forme + accent + négation) & ≥ 3 erreurs & 1–2 erreurs & 0 erreur \\
\hline
Slogan percutant (court, rythmé, vocabulaire juste) & Absent ou confus & Correct mais faible & Percutant, mémorable \\
\hline
Présentation (titre, lisibilité, illustration, signature) & Incomplète / illisible & Complète, perfectible & Soignée, prête à afficher \\
\hline
Présentation orale (5 phrases, fluidité, vocabulaire) & < 3 phrases, hésitant & 4–5 phrases, quelques fautes & 5 phrases, fluide, précis \\
\hline
\end{tabularx}
\end{center}
\textbf{Score ≥ 8/10 = affiche validée pour l'oral du Bac.}

\begin{exercice}[Entraînement : transformer en impératif tú (affirmatif et négatif)]
Mettez les verbes entre parenthèses à l'impératif \textbf{tú} (affirmatif ou négatif selon le sens).
\begin{enumerate}
\item \esp{\_\_\_\_\_\_\_\_\_\_\_\_ (usar) bolsas de tela, \_\_\_\_\_\_\_\_\_\_\_\_ (tirar) plástico al suelo.}
\item \esp{\_\_\_\_\_\_\_\_\_\_\_\_ (separar) la basura, \_\_\_\_\_\_\_\_\_\_\_\_ (mezclar) orgánico e inorgánico.}
\item \esp{\_\_\_\_\_\_\_\_\_\_\_\_ (apagar) la luz al salir, \_\_\_\_\_\_\_\_\_\_\_\_ (dejar) el grifo abierto.}
\item \esp{\_\_\_\_\_\_\_\_\_\_\_\_ (plantar) un árbol cada año, \_\_\_\_\_\_\_\_\_\_\_\_ (cortar) árboles sin permiso.}
\item \esp{\_\_\_\_\_\_\_\_\_\_\_\_ (denunciar) la tala ilegal, \_\_\_\_\_\_\_\_\_\_\_\_ (callar) ante la contaminación.}
\end{enumerate}
\end{exercice}

\begin{corrige}[Exercice n°1]
\begin{enumerate}
\item \esp{¡Usa bolsas de tela, no tires plástico al suelo!}
\item \esp{¡Separa la basura, no mezcles orgánico e inorgánico!}
\item \esp{¡Apaga la luz al salir, no dejes el grifo abierto!}
\item \esp{¡Planta un árbol cada año, no cortes árboles sin permiso!}
\item \esp{¡Denuncia la tala ilegal, no calles ante la contaminación!}
\end{enumerate}
\emph{Note :} \esp{no tires} (subj. \esp{tirar}), \esp{no mezcles} (subj. \esp{mezclar}), \esp{no dejes} (subj. \esp{dejar} — irrégulier), \esp{no cortes} (subj. \esp{cortar}), \esp{no calles} (subj. \esp{callar}). Accents vérifiés.
\end{corrige}

\begin{exercice}[Production libre : votre affiche en 30 minutes]
En binôme, réalisez l'affiche complète (papier A3 fourni) sur l'un des trois thèmes de l'Étape 1. Respectez la checklist de la méthode. Prenez une photo pour l'oral blanc de la semaine prochaine.
\end{exercice}

\newpage

\coursec{Activité d'intégration}

\coursec{Lectura y comentario : contaminación y medidas de preservación ; creación de carteles de sensibilización}

\courssub{Objectifs de l'activité}

Cette activité d'intégration vise à mobiliser l'ensemble des savoirs et savoir-faire travaillés au cours de la leçon E09 dans une tâche finale complexe, authentique et motivante. Les élèves devront :

\begin{itemize}
    \item \textbf{Comprendre} une situation de communication réelle ancrée dans leur environnement immédiat (quartier de Yaoundé, Douala ou toute ville camerounaise) : la dégradation de l'environnement local.
    \item \textbf{Mobiliser} le lexique thématique de la pollution et de la préservation (\cle{contaminación}, \cle{calentamiento global}, \cle{deforestación}, \cle{sequía}, \cle{capa de ozono}, \cle{energías renovables}, \cle{residuos plásticos}, \cle{aire insalubre}, \cle{agua contaminada}, etc.).
    \item \textbf{Appliquer} la morphologie et la syntaxe de l'\textbf{impératif} (affirmatif et négatif, formes \emph{ustedes} / \emph{vosotros} selon la norme choisie, pronoms enclitiques) pour formuler des consignes d'action concrètes.
    \item \textbf{Maîtriser} les techniques du \textbf{commentaire de texte} (identification du thème, structure, lexique, visée argumentative) pour présenter leur affiche à l'oral comme une mini-conférence.
    \item \textbf{Développer} des compétences transverses : travail collaboratif (groupes de 4), créativité (slogan, dessin), expression orale continue, évaluation par les pairs (grille critériée), citoyenneté écoresponsable.
\end{itemize}

\courssub{Situation}

\begin{textebox}[Contexte : Appel à projets « Yaoundé Propre, Yaoundé Vert »]
\emph{Contexto real :} El Ayuntamiento de Yaoundé, en colaboración con el Ministerio de Medio Ambiente (MINEPDED) y varias ONG ecologistas (como \emph{ACDIC} o \emph{Green Cameroon}), lanza una campaña de sensibilización ciudadana titulada \textbf{« ¡Salvemos nuestro barrio! »}. El objetivo es combatir la proliferación de residuos plásticos en los arroyos (Mfoundi, Mingoa), la quema de basura a cielo abierto que emite humos tóxicos, la tala anárquica de árboles para leña y la contaminación de los puntos de agua potable.

Como alumnos de Terminale A del Liceo Bilingüe de Essos (o de cualquier establecimiento de la capital), ustedes forman parte del \textbf{Club Verde Escolar}. Su profesor de español les propone participar en el concurso de carteles: cada grupo diseña un cartel completo en español que será expuesto en el pasillo del liceo y, los tres mejores, en la sala de actos durante la \emph{Semana del Medio Ambiente} (primer semana de junio).

\medskip
\noindent \textbf{Documento de partida (extracto del comunicado municipal) :}
\medskip

« \emph{La Municipalidad de Yaoundé invita a la juventud estudiantil a convertirse en embajadores del cambio. Los barrios de Mvog-Ada, Etoug-Ebe, Nkolbisson y Obili sufren una alarmante acumulación de botellas PET, bolsas plásticas y neumáticos usados. El aire se vuelve irrespirable al atardecer por la quema de desechos. Los niños juegan cerca de arroyos convertidos en cloacas a cielo abierto. ¡Basta! Cada gesto cuenta: separar, reciclar, plantar, denunciar. Diseñen un cartel que despierte conciencias y mueva a la acción.} » 
\end{textebox}

\courssub{Consignes}

Les consignes sont organisées en \textbf{quatre phases successives}. Chaque phase fait l'objet d'un contrôle intermédiaire par l'enseignant (validation du lexique, de la correction de l'impératif, de la cohérence du slogan, de la qualité de l'oral).

\begin{enumerate}
    \item \textbf{Phase 1 — Analyse du contexte et choix du problème (15 min).}
    \begin{enumerate}
        \item En groupe, lisez le communiqué municipal (texte ci-dessus). Identifiez \textbf{trois problèmes écologiques concrets} mentionnés (ex. : \emph{acumulación de botellas PET}, \emph{quema de desechos}, \emph{arroyos convertidos en cloacas}).
        \item Choisissez \textbf{UN seul problème} que votre groupe va traiter. Justifiez votre choix en une phrase en espagnol (ex. : \esp{Hemos elegido la contaminación del arroyo Mfoundi porque afecta directamente a nuestro barrio.} — « Nous avons choisi la pollution du Mfoundi car elle affecte directement notre quartier. »).
        \item Dressez une \textbf{liste de 10--12 mots-clés} (noms, verbes, adjectifs) indispensables pour votre affiche. Utilisez le lexique de la leçon + dictionnaire si besoin. Ex. : \cle{vertederos ilegales}, \cle{humo tóxico}, \cle{reciclar}, \cle{plantar árboles}, \cle{denunciar}, \cle{concienciar}, \cle{hábitos sostenibles}.
    \end{enumerate}

    \item \textbf{Phase 2 — Conception de l'affiche (30 min).}
    L'affiche doit comporter \textbf{cinq zones obligatoires} (disposition libre, format A3 ou A2) :
    \begin{enumerate}
        \item \textbf{Titre accrocheur} (max 6 mots, majuscules, centré). Ex. : \esp{¡BASTA PLÁSTICO EN NUESTROS ARROYOS!}
        \item \textbf{Slogan percutant} (phrase courte, rythme, rime ou allitération bienvenue). Ex. : \esp{Un río limpio, un barrio sano.} / \esp{No quemes tu futuro, recicla tu presente.}
        \item \textbf{Quatre consignes à l'impératif} (forme \emph{ustedes} pour la politesse et l'universalité, ou \emph{vosotros} si l'établissement suit la norme péninsulaire ; \textbf{au moins une forme négative}). Chaque consigne commence par le verbe à l'impératif. Exemples :
        \begin{itemize}
            \item \esp{Separen los residuos en casa.}
            \item \esp{No quemen basura en la calle.}
            \item \esp{Planten un árbol cada año.}
            \item \esp{Denuncien los vertederos ilegales al 119.}
        \end{itemize}
        \item \textbf{Appel final / signature} : phrase mobilisatrice + mention du groupe / lycée. Ex. : \esp{¡Únete al cambio! Club Verde — Liceo Bilingüe de Essos.}
        \item \textbf{Illustration} : dessin, pictogrammes, collage, code QR vers une vidéo courte (optionnel). L'image doit renforcer le message, pas le décorer seulement.
    \end{enumerate}
    \begin{attention}[Pièges grammaticaux fréquents]
    \begin{itemize}
        \item \textbf{Impératif affirmatif + pronom} : le pronom s'enclitique et l'accent écrit se déplace : \esp{Sepárenlos} (pas *separenlos), \esp{Plántenlo}, \esp{Denúncienlo}. Règle : accent sur la voyelle tonique du verbe original.
        \item \textbf{Impératif négatif + pronom} : le pronom précède le verbe : \esp{No los quemen}, \esp{No lo tiren}.
        \item \textbf{Verbes irréguliers à l'impératif \emph{ustedes}} : \esp{Sean} (ser), \esp{Estén} (estar), \esp{Tengan} (tener), \esp{Vayan} (ir), \esp{Sepan} (saber), \esp{Hagan} (hacer), \esp{Pongan} (poner), \esp{Salgan} (salir), \esp{Traigan} (traer), \esp{Digan} (decir).
        \item \textbf{Verbes à changement orthographique (réguliers morphologiquement)} : \esp{reduzcan} (c $\to$ zc devant \emph{a/o}), \esp{protejan} (g $\to$ j devant \emph{a/o}), \esp{reciclen} (z $\to$ c devant \emph{e/i} — mais \emph{reciclen} s'écrit sans accent, voir tableau).
        \item \textbf{Orthographe} : pas de *\emph{reciclen} avec \emph{z} ; \emph{reduzcan} (c $\to$ zc) ; \emph{protejan} (g $\to$ j).
    \end{itemize}
    \end{attention}

    \item \textbf{Phase 3 — Préparation de la présentation orale (15 min).}
    Chaque groupe prépare une \textbf{présentation de 5 phrases exactement} (chronométrée : 1 min 30 max), structurée comme un \textbf{mini-commentaire de texte} appliqué à leur propre affiche. Modèle imposé :
    \begin{enumerate}
        \item \textbf{Thème et visée} : \esp{Nuestro cartel denuncia la contaminación del arroyo Mfoundi y llama a la acción ciudadana.}
        \item \textbf{Analyse du slogan} : \esp{El eslogan « Un río limpio, un barrio sano » usa la rima y el paralelismo para fijar el mensaje en la memoria.}
        \item \textbf{Justification des consignes} : \esp{Las cuatro consignes en imperativo (separen, no quemen, planten, denuncien) cubren la prevención, la prohibición, la reparación y la vigilancia.}
        \item \textbf{Choix lexicaux} : \esp{Hemos empleado términos técnicos como « vertederos ilegales », « humo tóxico » y « residuos PET » para dar credibilidad científica.}
        \item \textbf{Appel final} : \esp{Esperamos que este cartel, expuesto en el pasillo, motive a cada estudiante a cambiar un hábito hoy mismo.}
    \end{enumerate}
    \textbf{Contraintes} : pas de lecture papier, regard vers le public, prononciation soignée (accent tonique, \emph{r} simple vs \emph{rr}, \emph{ñ}, \emph{ll}, \emph{j} / \emph{g} géodésique). L'enseignant note sur 5 (contenu, structure, correction, prononciation, engagement).

    \item \textbf{Phase 4 — Exposition, vote et bilan (20 min).}
    \begin{enumerate}
        \item Les affiches sont affichées au tableau / sur les murs du couloir. Chaque élève (sauf son propre groupe) dispose de \textbf{3 gommettes} (or, argent, bronze) à coller au dos de ses 3 affiches préférées selon la \textbf{grille d'évaluation} ci-dessous.
        \item \textbf{Grille d'évaluation (sur 20 points, simplifiée pour le vote pair) :}
        \begin{center}
        \begin{tabularx}{\linewidth}{|X|c|X|}
        \hline
        \rowcolor{popblueL} \textbf{Critère} & \textbf{Points} & \textbf{Indicateurs observables} \\
        \hline
        Pertinence du problème (local, réel) & 4 & Le problème touche vraiment le quartier ; données concrètes. \\
        Qualité du slogan (court, mémorable, espagnol correct) & 4 & Rime, allitération, parallélisme ; pas de faute. \\
        Impératifs (4, variés, corrects, pronoms bien placés) & 6 & Affirmatif + négatif ; verbes irréguliers maîtrisés ; enclitiques. \\
        Richesse lexicale (environnement + action) & 3 & Minimum 8 mots du lexique ciblé utilisés à bon escient. \\
        Esthétique / lisibilité / originalité & 3 & Titre visible, illustration porteuse de sens, mise en page soignée. \\
        \hline
        \end{tabularx}
        \end{center}
        \item Dépouillement public, annonce des 3 groupes lauréats. Remise de « diplômes \emph{Embajadores Verdes} » + exposition en salle de actes.
        \item \textbf{Bilan réflexif collectif (5 min)} : l'enseignant anime un tour de table rapide : \esp{¿Qué han aprendido? ¿Qué les costó más? ¿Cambiarán un hábito?} Réponses en espagnol si possible, sinon français accepté.
    \end{enumerate}
\end{enumerate}

\courssub{Ressources à mobiliser}

Cette fiche de révision rapide est distribuée aux groupes (ou projetée) pendant les phases 2 et 3.

\begin{propriete}[Impératif \emph{ustedes} — Formes régulières]
\begin{center}
\begin{tabular}{|l|l|l|l|}
\hline
\rowcolor{popblueL} \textbf{Verbe} & \textbf{Racine} & \textbf{Terminaison \emph{-ar}} & \textbf{Terminaison \emph{-er/-ir}} \\
\hline
Hablar & habl- & \textbf{Hablen} & — \\
Comer & com- & — & \textbf{Coman} \\
Vivir & viv- & — & \textbf{Vivan} \\
\hline
\end{tabular}
\end{center}
\emph{Rappel} : l'impératif \emph{ustedes} = subjonctif présent 3e personne du pluriel. Pour les verbes réguliers, \textbf{pas d'accent écrit} car la tonicité tombe sur l'avant-dernière syllabe (mots en \emph{-n}) : \esp{reciclen}, \esp{reduzcan}, \esp{protejan}, \esp{separen}, \esp{quemen}, \esp{planten}, \esp{denuncien}.
\end{propriete}

\begin{propriete}[Impératif \emph{ustedes} — Verbes irréguliers fréquents (à connaître par cœur)]
\begin{center}
\begin{tabular}{|l|l|l|}
\hline
\rowcolor{popblueL} \textbf{Infinitif} & \textbf{Impératif \emph{ustedes}} & \textbf{Traduction} \\
\hline
Ser & \textbf{Sean} & Soyez \\
Estar & \textbf{Estén} & Soyez (état/lieu) \\
Tener & \textbf{Tengan} & Ayez \\
Ir & \textbf{Vayan} & Allez \\
Saber & \textbf{Sepan} & Sachez \\
Hacer & \textbf{Hagan} & Faites \\
Poner & \textbf{Pongan} & Mettez \\
Salir & \textbf{Salgan} & Sortez \\
Traer & \textbf{Traigan} & Apportez \\
Decir & \textbf{Digan} & Dites \\
\hline
\end{tabular}
\end{center}
\end{propriete}

\begin{propriete}[Verbes à changement orthographique à l'impératif \emph{ustedes} (réguliers morphologiquement)]
\begin{center}
\begin{tabular}{|l|l|l|l|}
\hline
\rowcolor{popblueL} \textbf{Infinitif} & \textbf{Changement} & \textbf{Impératif \emph{ustedes}} & \textbf{Traduction} \\
\hline
Reducir & c $\to$ zc & \textbf{Reduzcan} & Réduisez \\
Proteger & g $\to$ j & \textbf{Protejan} & Protégez \\
Reciclar & z $\to$ c (devant e) & \textbf{Reciclen} & Recyclez \\
Conseguir & g $\to$ j & \textbf{Consigan} & Obtenez \\
\hline
\end{tabular}
\end{center}
\emph{Note} : ces verbes suivent la règle régulière du subjonctif ; le changement d'orthographe sert à conserver le son du radical.
\end{propriete}

\begin{definition}[Lexique clé — Environnement \& Action (niveau Terminale A)]
\begin{itemize}
    \item \textbf{Noms :} \cle{la contaminación} (pollution), \cle{el calentamiento global} (réchauffement climatique), \cle{la deforestación} (déforestation), \cle{la sequía} (sécheresse), \cle{la capa de ozono} (couche d'ozone), \cle{las energías renovables} (énergies renouvelables), \cle{los residuos / desechos} (déchets), \cle{el plástico / el PET}, \cle{el vertedero (ilegal)} (décharge sauvage), \cle{el humo tóxico} (fumée toxique), \cle{el arroyo / el río} (ruisseau / rivière), \cle{el aire insalubre} (air insalubre), \cle{el agua potable / contaminada}, \cle{la tala (ilegal)} (abattage), \cle{la reforestación} (reboisement), \cle{el reciclaje} (recyclage), \cle{la concienciación} (sensibilisation).
    \item \textbf{Verbes d'action (infinitif) :} \cle{contaminar}, \cle{quemar}, \cle{tirar (basura)}, \cle{separar}, \cle{reciclar}, \cle{reutilizar}, \cle{reducir}, \cle{plantar}, \cle{reforestar}, \cle{denunciar}, \cle{proteger}, \cle{conservar}, \cle{concienciar}, \cle{movilizar}, \cle{cambiar (hábitos)}.
    \item \textbf{Adjectifs :} \cle{contaminado/a}, \cle{insalubre}, \cle{tóxico/a}, \cle{sostenible}, \cle{ecológico/a}, \cle{verde}, \cle{urgente}, \cle{necesario/a}, \cle{responsable}.
    \item \textbf{Connecteurs pour l'oral :} \cle{por eso} (c'est pourquoi), \cle{por tanto} (par conséquent), \cle{además} (de plus), \cle{en primer lugar / en segundo lugar}, \cle{finalmente}, \cle{en conclusión}.
\end{itemize}
\end{definition}

\begin{popfigure}
\begin{tikzpicture}[
    mindmap,
    grow cyclic,
    every node/.style={concept, execute at begin node=\hskip0pt},
    root concept/.append style={concept color=popblue, minimum size=2.5cm, text=white, font=\bfseries\large},
    level 1 concept/.append style={concept color=popgreen, sibling angle=90, minimum size=2cm, text=white, font=\bfseries\small},
    level 2 concept/.append style={concept color=poporange, sibling angle=45, minimum size=1.6cm, text=white, font=\small},
    level 3 concept/.append style={concept color=poppurple, sibling angle=30, minimum size=1.3cm, text=white, font=\scriptsize},
    edge from parent/.style={draw=popline, thick}
]
\node {Affiche de sensibilisation}
    child { node {Titre accrocheur}
        child { node {Max 6 mots} }
        child { node {Majuscules} }
        child { node {Centré} }
    }
    child { node {Slogan percutant}
        child { node {Rime / parallélisme} }
        child { node {Court, mémorable} }
        child { node {Espagnol correct} }
    }
    child { node {4 Consignes (impératif)}
        child { node {Forme ustedes} }
        child { node {1 négative min.} }
        child { node {Variées : prévention, interdiction, réparation, vigilance} }
        child { node {Pronoms bien placés} }
    }
    child { node {Appel final + signature}
        child { node {Mobilisateur} }
        child { node {Groupe / Lycée} }
    }
    child { node {Illustration porteuse}
        child { node {Avant / Après} }
        child { node {Pictogrammes} }
        child { node {Code QR (option)} }
    };
\end{tikzpicture}
\legende{Structure mentale de l'affiche de sensibilisation (5 zones obligatoires)}
\end{popfigure}

\begin{methode}[Mini-commentaire oral en 5 phrases — Structure type]
\begin{enumerate}
    \item \textbf{Annonce thème + visée} : \esp{Nuestro cartel denuncia / propone / llama a... + [problème précis] + [objectif].}
    \item \textbf{Analyse slogan} : \esp{El eslogan « ... » usa [procédé : rime, parallélisme, métaphore, impératif, question rhétorique] para [effet visé].}
    \item \textbf{Justification consignes} : \esp{Las [número] consignes en imperativo ([verbe1], [verbe2], [verbe3], [verbe4]) cubren [dimensions : prevención, prohibición, reparación, vigilancia / participación].}
    \item \textbf{Choix lexicaux} : \esp{Hemos empleado términos como « ... », « ... » y « ... » para dar [credibilidad / fuerza / precisión] al mensaje.}
    \item \textbf{Appel / portée} : \esp{Esperamos que este cartel, expuesto en [lieu], motive a [public cible] a [action concrète].}
\end{enumerate}
\end{methode}

\courssub{Proposition de production et corrigé commenté}

\begin{exoresolu}[Affiche modèle — Groupe « Guardianes del Mfoundi » (Lycée Bilingüe d'Essos)]
\textbf{Énoncé :} Concevoir l'affiche complète (titre, slogan, 4 impératifs, appel, illustration décrite) pour le problème « Pollution du ruisseau Mfoundi par les déchets plastiques et brûlage à ciel ouvert », puis rédiger la présentation orale en 5 phrases.

\tcblower

\textbf{PRODUCTION MODÈLE (Affiche)}

\begin{center}
\fbox{\begin{minipage}{0.9\linewidth}
\centering
\textbf{\LARGE ¡SALVEMOS EL MFOUNDI, NUESTRA VIDA!} \\[0.5em]
\textit{\large « Un arroyo limpio, un barrio sano. »} \\[1em]
\makebox[\linewidth]{\rule{0.8\linewidth}{1pt}} \\[0.5em]
\textbf{\large 1. Separen los residuos plásticos en casa.} \\
\textbf{\large 2. No quemen basura ni neumáticos en la calle.} \\
\textbf{\large 3. Planten árboles nativos en las orillas.} \\
\textbf{\large 4. Denuncien los vertederos ilegales al 119.} \\[1em]
\makebox[\linewidth]{\rule{0.8\linewidth}{1pt}} \\[0.5em]
\textbf{\large ¡TÚ ERES EL CAMBIO! \quad Club Verde — Liceo Bilingüe de Essos}
\end{minipage}}
\end{center}

\medskip
\noindent \textbf{Description de l'illustration (au dos de l'affiche ou en légende) :}
\emph{En haut à gauche : un ruisseau marron, encombré de bouteilles PET, sacs noirs, pneus ; fumée grise montant d'un feu. En haut à droite : le même ruisseau bleu clair, berges arborées (acacias, manguiers), enfants qui jouent, poubelles de tri (jaune/vert/bleu). Au centre : une grande flèche verte « ANTES $\to$ DESPUÉS ». En bas : logo « Club Verde » + petit code QR dessiné.}

\medskip
\noindent \textbf{PRÉSENTATION ORALE MODÈLE (5 phrases)}

\begin{enumerate}
    \item \esp{Nuestro cartel denuncia la grave contaminación del arroyo Mfoundi por residuos plásticos y humo tóxico, y llama a la acción inmediata de los vecinos de Essos.}
    \item \esp{El eslogan « Un arroyo limpio, un barrio sano » usa el paralelismo sintáctico para vincular la salud del río a la del barrio, fijando el mensaje en la memoria.}
    \item \esp{Las cuatro consignes en imperativo — \emph{separen} (prevención), \emph{no quemen} (prohibición), \emph{planten} (reparación ecológica) y \emph{denuncien} (vigilancia ciudadana) — cubren el ciclo completo de la gestión ambiental.}
    \item \esp{Hemos empleado términos técnicos como « residuos PET », « vertederos ilegales », « humo tóxico » y « especies nativas » para dar credibilidad científica y anclar el problema en la realidad local.}
    \item \esp{Esperamos que este cartel, expuesto en el pasillo principal, motive a cada estudiante a separar su basura hoy mismo y a convertirse en guardián del Mfoundi.}
\end{enumerate}

\medskip
\noindent \textbf{COMMENTAIRE DES CHOIX DE LANGUE (pour l'enseignant)}

\begin{itemize}
    \item \textbf{Titre} : \esp{¡Salvemos el Mfoundi, nuestra vida!} — Impératif \emph{nosotros} (\emph{salvemos}) = appel à l'action collective inclusive ; \emph{nuestra vida} = appropriation identitaire forte. Ponctuation espagnole correcte (\emph{¡ !}).
    \item \textbf{Slogan} : \esp{Un arroyo limpio, un barrio sano.} — Parallélisme nominal (SN + adj.), rythme binaire (4+4 syllabes). Pas de verbe = universalité intemporelle. \emph{Ser} sous-entendu (caractéristique essentielle), pas \emph{estar} (état passager).
    \item \textbf{Impératif 1} : \esp{Separen los residuos plásticos en casa.} — Verbe régulier \emph{-ar} $\to$ \emph{separen}. COD précis \emph{los residuos plásticos} (pas \emph{basura} générique). CCL \emph{en casa} (action domestique accessible).
    \item \textbf{Impératif 2 (négatif)} : \esp{No quemen basura ni neumáticos en la calle.} — Négation \emph{no} + impératif \emph{quemen} (régulier \emph{-ar}, pas de diphtongue au subjonctif). \emph{Ni} = coordination négative. \emph{En la calle} = espace public.
    \item \textbf{Impératif 3} : \esp{Planten árboles nativos en las orillas.} — \emph{Planten} (régulier \emph{-ar}). \emph{Árboles nativos} = vocabulaire précis (espèces locales : \emph{acacia, mango, iroko}). \emph{En las orillas} = localisation technique (berges, fém. pl.).
    \item \textbf{Impératif 4} : \esp{Denuncien los vertederos ilegales al 119.} — \emph{Denuncien} (régulier \emph{-er} $\to$ \emph{-an}). \emph{Vertederos ilegales} = terme juridique/administratif. \emph{Al 119} = \emph{a el 119} (numéro d'urgence environnementale \textbf{fictif} au Cameroun, mais réaliste) ; contraction obligatoire \emph{al}.
    \item \textbf{Appel final} : \esp{¡Tú eres el cambio!} — Pronom personnel accentué \emph{tú} (insistance sur la responsabilité individuelle) + \emph{ser} (identité essentielle). Signature institutionnelle.
    \item \textbf{Oral} : Phrase 1 = thème + visée (verbes \emph{denuncia}, \emph{llama}). Phrase 2 = analyse rhétorique (métalangage : \emph{paralelismo sintáctico}). Phrase 3 = justification grammaticale/fonctionnelle (terminologie : \emph{imperativo}, \emph{prevención}, \emph{prohibición}, \emph{reparación}, \emph{vigilancia}). Phrase 4 = lexique spécialisé (crédibilité). Phrase 5 = portée pragmatique (lieu, public, action visée).
    \item \textbf{Erreurs évitées} : pas de *\emph{separenlos} (pronom enclitique inutile car COD non pronominalisé) ; pas de *\emph{no quemen la basura} (article défini évité pour généralité) ; pas de *\emph{planten los árboles} (même raison) ; accord \emph{nativos} (masc. pl.) avec \emph{árboles} ; \emph{orillas} fém. pl. $\to$ \emph{las orillas}.
\end{itemize}
\end{exoresolu}

\courssub{Bilan de l'activité}

\begin{aretenir}[Ce que cette activité a permis de valider]
\begin{itemize}
    \item \textbf{Compétence « Compréhension de l'écrit »} : lecture analytique d'un communiqué institutionnel authentique (registre formel, vocabulaire technique, structure argumentative).
    \item \textbf{Compétence « Expression écrite »} : production d'un texte court à visée persuasive (affiche) respectant des contraintes de format, de registre (impératif \emph{ustedes}), de correction morphosyntaxique (accents, enclitiques, irréguliers, changements orthographiques) et de cohérence thématique.
    \item \textbf{Compétence « Expression orale en continu »} : prise de parole préparée, structurée (méthode du commentaire), chronométrée, avec exigence de prononciation et de regard public.
    \item \textbf{Compétence « Interaction / Médiation »} : négociation dans le groupe (choix du problème, rédaction collaborative), évaluation par les pairs (grille critériée), retour réflexif final.
    \item \textbf{Savoirs linguistiques} : impératif \emph{ustedes} (réguliers, 10 irréguliers fréquents, verbes à changement orthographique), placement des pronoms (affirmatif/négatif), accord adjectif/nom, lexique environnemental (noms, verbes, adjectifs, connecteurs), ponctuation espagnole (\emph{¿ ¡}).
    \item \textbf{Savoirs culturels / citoyens} : ancrage dans la réalité camerounaise (Yaoundé, Mfoundi, quartiers, ONG locales), lien avec la Guinée équatoriale (seul pays hispanophone d'Afrique, même enjeux de déforestation / pollution plastique), ouverture sur les Objectifs de Développement Durable (ODD 6, 11, 13, 15).
\end{itemize}
\end{aretenir}

\begin{savaistu}[Le Cameroun et la Guinée équatoriale : deux poumons verts menacés]
Le Cameroun et la Guinée équatoriale partagent la forêt du bassin du Congo, deuxième massif forestier tropical après l'Amazonie. La \cle{deforestación} y avance à cause de l'agriculture sur brûlis, de l'exploitation forestière illégale et de la demande en bois de chauffe (90\% de l'énergie domestique à Yaoundé). Les \cle{residuos plásticos} bouchent les caniveaux, provoquant des inondations (Douala, 2023) et tuant la faune aquatique. Le \cle{calentamiento global} raccourcit la saison des pluies, menaçant la culture du cacao (premier export agricole). L'Espagne et l'Union européenne financent des projets \cle{energías renovables} (solaire au Nord-Cameroun, hydroélectrique à Memve'ele). La \emph{Semana del Medio Ambiente} (5 juin) est célébrée dans les lycées des deux pays avec des campagnes \emph{« Un árbol, una vida »}. Votre affiche s'inscrit dans cette dynamique transfrontalière !
\end{savaistu}

\begin{attention}[Erreurs fréquentes observées lors des sessions précédentes — Check-list pour l'auto-correction]
\begin{itemize}
    \item \textbf{Accents manquants sur l'impératif affirmatif + pronom} : *\emph{separenlos} $\to$ \esp{sepárenlos} ; *\emph{plantenlo} $\to$ \esp{plántenlo} ; *\emph{denuncienlo} $\to$ \esp{denúncienlo}.
    \item \textbf{Confusion \emph{ser} / \emph{estar} dans le slogan} : *\emph{Un arroyo está limpio} (état passager) vs \esp{Un arroyo es limpio} (caractéristique) — ici \emph{limpio} comme idéal = \emph{ser} ou adjectif épithète sans verbe (meilleur).
    \item \textbf{Faux ami \emph{contaminación} / \emph{contamination}} : en espagnol \emph{contaminación} = pollution (air, eau, sol) ; \emph{contaminación radiactiva} = contamination radioactive. Ne pas confondre avec \emph{infección} (médical).
    \item \textbf{Anglicisme \emph{realizar} au sens de « réaliser / comprendre »} : \emph{realizar} = exécuter, mener à bien. Pour « prendre conscience » : \esp{tomar conciencia}, \esp{darse cuenta}.
    \item \textbf{Oubli de la contraction \emph{al} / \emph{del}} : *\emph{a el 119} $\to$ \esp{al 119} ; *\emph{de el barrio} $\to$ \esp{del barrio}.
    \item \textbf{Majuscule systématique au début des vers du slogan} : en espagnol, seule la première lettre de la phrase prend majuscule (sauf noms propres) : \esp{« Un arroyo limpio, un barrio sano »} (pas *\emph{Un Arroyo Limpio...}).
    \item \textbf{Accents écrits injustifiés sur impératifs réguliers} : *\emph{recíclen}, *\emph{reduzcán}, *\emph{protèjan} $\to$ \esp{reciclen}, \esp{reduzcan}, \esp{protejan} (mots en \emph{-n}, tonique sur l'avant-dernière syllabe = pas d'accent).
    \item \textbf{Verbe \emph{reforestrar}} : préférer \esp{reforestar} (plus fréquent et standard).
\end{itemize}
\end{attention}

\newpage

\coursec{Méthodes et exercices résolus}

\courssub{Méthodes et exercices résolus}

\begin{methode}[Commenter un texte d'actualité : pollution et environnement]
\textbf{Objectif :} Répondre aux questions de compréhension, analyser la structure et le lexique, puis produire un commentaire personnel structuré.
\begin{enumerate}
\item \textbf{Lecture globale :} Lire le texte une première fois sans s'arrêter. Identifier le \cle{type de texte} (article de presse, rapport, tribune), la \cle{source}, le \cle{thème} principal et la \cle{thèse} défendue.
\item \textbf{Lecture détaillée :} Souligner les \cle{mots-clés} du thème (\esp{contaminación}, \esp{energías renovables}, etc.), les \cle{connecteurs logiques} (\esp{por tanto}, \esp{sin embargo}, \esp{además}) et les \cle{chiffres/dates}.
\item \textbf{Analyse du vocabulaire :} Repérer les champs lexicaux (pollution / solutions), les \cle{faux amis} et les \cle{verbes d'action} (\esp{reducir}, \esp{fomentar}, \esp{proteger}).
\item \textbf{Structure du commentaire :}
\begin{itemize}
\item \textbf{Introduction :} Accroche, présentation du document (titre, auteur, date, source), thème, thèse, annonce du plan (2 ou 3 parties).
\item \textbf{Développement (2 parties) :} Ex. : I. Les causes et conséquences alarmantes. II. Les pistes de solutions et l'engagement citoyen. Citer le texte (guillemets \esp{« ... »}) et expliquer.
\item \textbf{Conclusion :} Bilan synthétique + ouverture personnelle (contexte camerounais, rôle des jeunes, avenir de la planète).
\end{itemize}
\item \textbf{Relecture :} Vérifier l'orthographe (accents !), les accords, l'usage du \cle{subjonctif} après les expressions impersonnelles (\esp{es necesario que} + subj.), la cohérence des temps.
\end{enumerate}
\end{methode}

\begin{exoresolu}[Commentaire guidé : « Un planeta en peligro »]
\textbf{Énoncé :} Lisez l'extrait suivant et répondez aux questions en espagnol. Puis rédigez une introduction de commentaire type Bac.

\begin{textebox}[Extrait original -- Article d'opinion simulé]
\esp{El calentamiento global ya no es una predicción futura, es una realidad palpable. En las grandes ciudades como Madrid, Ciudad de México o Yaoundé, la contaminación del aire supera los límites recomendados por la OMS. La deforestación avanza a un ritmo vertiginoso: cada año desaparecen millones de hectáreas de bosque. Sin embargo, la transición hacia energías renovables ofrece una esperanza real. Invertir en energía solar y eólica no solo reduce la huella de carbono, sino que crea empleo verde. ¿Qué esperamos para actuar?}
\end{textebox}

\textbf{Questions :}
\begin{enumerate}
\item Identifiez le thème principal et la thèse de l'auteur.
\item Relevez 4 termes du champ lexical de la \cle{dégradation} et 3 du champ des \cle{solutions}.
\item Quel temps verbal est utilisé dans \esp{es una realidad} ? Justifiez.
\item Transformez \esp{¿Qué esperamos para actuar?} en style indirect : \esp{El autor se pregunta...}
\item Traduisez en français : \esp{Invertir en energía solar y eólica no solo reduce la huella de carbono, sino que crea empleo verde.}
\end{enumerate}

\tcblower
\textbf{Solution détaillée :}

\textbf{1. Thème et thèse.}
\begin{itemize}
\item \textbf{Thème :} Le réchauffement climatique et la pollution urbaine (air, forêts) face à la transition énergétique.
\item \textbf{Thèse :} L'urgence est réelle (constat alarmant), mais des solutions concrètes existent (énergies renouvelables) ; l'auteur appelle à l'action immédiate.
\end{itemize}

\textbf{2. Champs lexicaux.}
\begin{center}
\begin{tabularx}{\linewidth}{|X|X|}
\hline
\rowcolor{popblueL} \textbf{Dégradation / Problème} & \textbf{Solutions / Espoir} \\ \hline
\esp{calentamiento global} & \esp{transición} \\ \hline
\esp{contaminación del aire} & \esp{energías renovables} \\ \hline
\esp{deforestación} & \esp{esperanza real} \\ \hline
\esp{ritmo vertiginoso} & \esp{empleo verde} \\ \hline
\esp{huella de carbono} (conséquence) & \esp{actuar} \\ \hline
\end{tabularx}
\end{center}

\textbf{3. Analyse verbale : \esp{es una realidad}.}
Verbe \esp{ser}, \textbf{présent de l'indicatif}, 3\textsuperscript{e} pers. sg. \textbf{Justification :} \esp{ser} identifie une caractéristique essentielle, une vérité générale. Le réchauffement est présenté comme un fait établi. (Contraste : \esp{está} exprimerait un état passager ou une localisation.)

\textbf{4. Style indirect.}
\esp{El autor se pregunta \textbf{qué esperamos para actuar}.} « L'auteur se demande ce que nous attendons pour agir. »
\textit{Règle :} pas de changement de temps (présent $\to$ présent) car le verbe introducteur (\esp{se pregunta}) est au présent. L'interrogatif \esp{qué} garde son accent en question indirecte. Pas de \esp{si}, car c'est une question partielle (mot interrogatif).

\textbf{5. Traduction (Version).}
\begin{quote}
« Investir dans l'énergie solaire et éolienne non seulement réduit l'empreinte carbone, mais crée aussi de l'emploi vert. »
\end{quote}
\textit{Points de vigilance :}
\begin{itemize}
\item \esp{invertir en} = investir dans (jamais \esp{invertir a}).
\item \esp{no solo ... sino que} = non seulement ... mais (aussi). Structure corrélative suivie de l'indicatif (fait réel).
\item \esp{huella de carbono} = empreinte carbone (terme figé).
\item \esp{empleo verde} = emploi vert (singulier collectif en espagnol).
\end{itemize}

\textbf{6. Proposition d'introduction (Bac).}
\end{exoresolu}

\begin{exemplebox}[Introduction type]
\esp{El artículo de opinión « Un planeta en peligro », publicado en un diario digital hispano, aborda la problemática urgente del calentamiento global y de la contaminación urbana. El autor denuncia la gravedad de la situación -- aire irrespirable, deforestación masiva -- pero destaca que la transición energética hacia las renovables constituye una solución viable y generadora de empleo. ¿Estamos a tiempo de revertir el daño? En primer lugar, analizaremos las causas y consecuencias expuestas; en segundo lugar, estudiaremos las soluciones propuestas y el llamamiento a la acción ciudadana.}
\end{exemplebox}
\textit{Traduction :} « L'article d'opinion "Une planète en danger", publié dans un quotidien numérique hispanophone, aborde le problème urgent du réchauffement climatique et de la pollution urbaine. L'auteur dénonce la gravité de la situation -- air irrespirable, déforestation massive -- mais souligne que la transition énergétique vers les renouvelables constitue une solution viable et créatrice d'emploi. Sommes-nous encore à temps de réparer les dégâts ? Nous analyserons d'abord les causes et conséquences exposées, puis les solutions proposées et l'appel à l'action citoyenne. »


\begin{methode}[Maîtriser le lexique thématique : pollution, climat, préservation]
\textbf{Objectif :} Mémoriser, orthographier (accents !) et employer le vocabulaire clé du module dans des phrases complètes.
\begin{enumerate}
\item \textbf{Classer par familles :} \cle{Causes} (\esp{industria}, \esp{transporte}, \esp{deforestación}), \cle{Conséquences} (\esp{calentamiento global}, \esp{sequía}, \esp{efecto invernadero}, \esp{agujero en la capa de ozono}), \cle{Solutions} (\esp{energías renovables}, \esp{reciclaje}, \esp{desarrollo sostenible}).
\item \textbf{Travailler les collocations (associations naturelles) :} ne pas apprendre les mots isolés.
\begin{itemize}
\item \esp{emitir gases} / \esp{reducir emisiones} (émettre des gaz / réduire les émissions)
\item \esp{proteger la capa de ozono} / \esp{agotar los recursos} (protéger la couche d'ozone / épuiser les ressources)
\item \esp{fomentar el reciclaje} / \esp{concienciar a la población} (encourager le recyclage / sensibiliser la population)
\end{itemize}
\item \textbf{Attention aux faux amis et pièges orthographiques :}
\begin{itemize}
\item \esp{contaminación} (pollution générale) $\neq$ \esp{infección} (infection médicale).
\item \esp{residuos} (déchets) $\neq$ \esp{desperdicios} (gaspillage, restes de nourriture).
\item \esp{calentamiento} : \textbf{sans accent} (mot terminé par -o, syllabe tonique \esp{-mien-}).
\item \esp{sequía} (sécheresse) : accent sur \textbf{í} (hiatus \esp{i-a}).
\item \esp{medio ambiente} (environnement) : traditionnellement en deux mots ; la graphie soudée \esp{medioambiente} est admise par la RAE, mais la forme en deux mots reste la plus sûre à l'examen.
\end{itemize}
\item \textbf{Entraînement :} traduire 10 phrases Fr $\to$ Esp imposant ce vocabulaire. Vérifier le genre (\esp{la sequía}, \esp{el calentamiento}) et l'accord des adjectifs (\esp{energías renovables}, \esp{desarrollo sostenible}).
\end{enumerate}
\end{methode}

\begin{exoresolu}[Lexique : du mot à la phrase contextuelle]
\textbf{Énoncé :}
\begin{enumerate}
\item Complétez les phrases avec le mot exact du lexique (accents obligatoires).
\begin{itemize}
\item[a)] La \_\_\_\_\_\_\_\_ de la \_\_\_\_\_\_\_\_ de ozono aumenta el riesgo de cáncer de piel.
\item[b)] Los gobiernos deben \_\_\_\_\_\_\_\_ el uso de \_\_\_\_\_\_\_\_ renovables.
\item[c)] La \_\_\_\_\_\_\_\_ prolongada destruye las cosechas en el norte de Camerún.
\item[d)] Cada ciudadano puede reducir su \_\_\_\_\_\_\_\_ ecológica reciclando.
\end{itemize}
\item Traduisez en espagnol (soignez les accords et les prépositions) :
\begin{itemize}
\item[a)] La déforestation menace la biodiversité amazonienne.
\item[b)] Nous devons sensibiliser les jeunes au changement climatique.
\item[c)] L'industrie pollue l'air et l'eau par ses déchets toxiques.
\item[d)] Les énergies propres sont l'avenir de l'humanité.
\end{itemize}
\item Corrigez les erreurs dans ces phrases d'élèves :
\begin{itemize}
\item[a)] \esp{La contaminación es un problema muy grave para el medioambiente.} (1 erreur de graphie)
\item[b)] \esp{Debemos proteger la capa de ozono para evitar el calentamiento global.} (1 erreur conceptuelle)
\item[c)] \esp{La sequia causa muchos problemas en la agricultura.} (1 erreur d'orthographe)
\item[d)] \esp{Usamos energias renovables para reducir la huella de carbon.} (2 erreurs d'accent)
\end{itemize}
\end{enumerate}

\tcblower
\textbf{Solution détaillée :}

\textbf{1. Complétion.}
\begin{itemize}
\item[a)] \esp{La \textbf{destrucción} de la \textbf{capa} de ozono...} (collocations possibles : \esp{el deterioro de la capa de ozono}, \esp{el agujero en la capa de ozono}).
\item[b)] \esp{... deben \textbf{fomentar} el uso de \textbf{energías} renovables.} (\esp{fomentar} + nom : registre soutenu ; synonymes : \esp{promover}, \esp{impulsar}.)
\item[c)] \esp{La \textbf{sequía} prolongada...} Accent obligatoire sur \textbf{í} (hiatus).
\item[d)] \esp{... su \textbf{huella} ecológica...} Expression figée : \esp{huella ecológica} / \esp{huella de carbono}.
\end{itemize}

\textbf{2. Thème (Français $\to$ Espagnol).}
\begin{itemize}
\item[a)] \esp{La deforestación amenaza la biodiversidad amazónica.}
\textit{Note :} \esp{amenazar} + COD (sans préposition). \esp{Biodiversidad} : féminin singulier.
\item[b)] \esp{Debemos concienciar a los jóvenes sobre el cambio climático.}
\textit{Note :} \esp{concienciar a alguien \textbf{sobre/de} algo}. Terme officiel : \esp{cambio climático} (pas \esp{cambio de clima}).
\item[c)] \esp{La industria contamina el aire y el agua con sus residuos tóxicos.}
\textit{Note :} \esp{contaminar} + COD + \esp{con} (instrument). \esp{Residuos tóxicos} est préférable à \esp{desechos} pour les déchets industriels.
\item[d)] \esp{Las energías limpias son el futuro de la humanidad.}
\textit{Note :} \esp{energías limpias} = synonyme courant d'\esp{energías renovables}. \esp{Ser} pour une définition. \esp{La humanidad} : féminin singulier.
\end{itemize}

\textbf{3. Correction d'erreurs types.}
\begin{itemize}
\item[a)] \esp{... para el \textbf{medio ambiente}.} Graphie recommandée en deux mots (la forme soudée \esp{medioambiente} est admise mais moins attendue à l'examen).
\item[b)] \esp{Debemos proteger la capa de ozono para evitar \textbf{la radiación ultravioleta} (o \textbf{el agujero en la capa de ozono}).}
\textit{Explication :} la couche d'ozone filtre les rayons UV. Le \textbf{réchauffement climatique} (\esp{calentamiento global}) est causé par l'effet de serre (CO\textsubscript{2}, méthane), pas par le trou dans la couche d'ozone. Confusion conceptuelle classique à éviter dans un commentaire.
\item[c)] \esp{La \textbf{sequía} causa muchos problemas en la agricultura.} Accent sur \textbf{í} : hiatus \esp{i-a}, la voyelle fermée accentuée rompt la diphtongue.
\item[d)] \esp{Usamos \textbf{energías} renovables para reducir la huella de \textbf{carbono}.}
\textit{Erreurs :} \esp{energías} (accent sur \textbf{í}, hiatus comme \esp{sequía}) ; \esp{carbono} (accent sur \textbf{ó}, mot esdrújulo : \esp{car-bo-no}, syllabe tonique \esp{car-}). En revanche \esp{renovables} ne prend pas d'accent (mot terminé par -s, tonique sur l'avant-dernière).
\end{itemize}
\end{exoresolu}

\begin{methode}[L'impératif affirmatif et négatif : l'outil de l'affiche]
\textbf{Objectif :} Former correctement l'impératif (\esp{tú}, \esp{usted}, \esp{nosotros}, \esp{vosotros}, \esp{ustedes}) pour écrire des slogans percutants (affirmatif = ordre/conseil ; négatif = interdiction).
\begin{enumerate}
\item \textbf{Rappel de formation (verbes réguliers) :}
\begin{center}
\begin{tabular}{|l|c|c|c|}
\hline
\rowcolor{popblueL} & \textbf{-AR (hablar)} & \textbf{-ER (comer)} & \textbf{-IR (vivir)} \\ \hline
\textbf{Tú (affirmatif)} & habla & come & vive \\ \hline
\textbf{Usted} & hable & coma & viva \\ \hline
\textbf{Nosotros} & hablemos & comamos & vivamos \\ \hline
\textbf{Vosotros} & hablad & comed & vivid \\ \hline
\textbf{Ustedes} & hablen & coman & vivan \\ \hline
\end{tabular}
\end{center}
\textit{Astuce :} l'impératif \esp{tú} affirmatif = 3\textsuperscript{e} pers. sg du présent de l'indicatif ; toutes les autres formes = subjonctif présent.
\item \textbf{Forme négative (tous les sujets) :} \esp{no} + \textbf{subjonctif présent}.
Ex. : \esp{no hables}, \esp{no hable}, \esp{no hablemos}, \esp{no habléis}, \esp{no hablen}.
\item \textbf{Irréguliers fréquents (affirmatif \esp{tú}) :} \esp{ir $\to$ ve}, \esp{ser $\to$ sé}, \esp{tener $\to$ ten}, \esp{venir $\to$ ven}, \esp{poner $\to$ pon}, \esp{salir $\to$ sal}, \esp{hacer $\to$ haz}, \esp{decir $\to$ di}. \textbf{Moyen mnémotechnique :} la phrase anglaise « \cle{Vin Diesel Has Ten Weapons} » rappelle \esp{Ven, Di, Sal, Haz, Ten, Ve, Pon, Sé}.
\item \textbf{Pronoms objets (position) :}
\begin{itemize}
\item Affirmatif : pronoms \textbf{enclitiques} (accrochés derrière, accent écrit pour conserver la tonique) : \esp{Recíclalo}, \esp{Cuídalo}, \esp{¡Hagámoslo!}
\item Négatif : pronoms \textbf{proclitiques} (devant le verbe) : \esp{No lo tires}, \esp{No lo hagamos}.
\end{itemize}
\item \textbf{Choix du sujet pour l'affiche :}
\begin{itemize}
\item \esp{Tú / Vosotros} : ton direct, jeune, mobilisateur (\esp{¡Recicla!}, \esp{¡No contamines!}).
\item \esp{Usted / Ustedes} : ton formel, institutionnel (\esp{Por favor, reduzca su consumo}).
\item \esp{Nosotros} : engagement collectif, « nous » inclusif (\esp{¡Protejamos nuestro planeta!}).
\end{itemize}
\end{enumerate}
\end{methode}

\begin{exoresolu}[Transformer des verbes en slogans d'affiche]
\textbf{Énoncé :} Pour une campagne « Yaoundé Verde », mettez les verbes entre parenthèses à la forme d'impératif adaptée au contexte indiqué. Justifiez le choix du sujet.

\begin{enumerate}
\item Contexte : affiche dans un lycée, s'adressant aux élèves (tutoiement). \esp{(\_\_\_\_\_\_\_\_ -- separar) la basura en el recreo.}
\item Contexte : panneau officiel municipal, ton formel. \esp{(\_\_\_\_\_\_\_\_ -- no / tirar) papeles al suelo.}
\item Contexte : manifeste d'une association de jeunes, engagement collectif. \esp{(\_\_\_\_\_\_\_\_ -- plantar) árboles en el barrio.}
\item Contexte : conseil santé/environnement à la radio (général). \esp{(\_\_\_\_\_\_\_\_ -- usar) transporte público.}
\item Contexte : slogan choc « Ne gaspillez pas l'eau ! » (pluriel familier, Espagne) puis (formel, Amérique latine / Cameroun).
\item Placez le pronom au bon endroit : \esp{(reciclar + lo)} $\to$ \_\_\_\_\_\_\_\_ (affirmatif \esp{tú}) ; \esp{(no / contaminar + la)} $\to$ \_\_\_\_\_\_\_\_ (négatif \esp{usted}).
\end{enumerate}

\tcblower
\textbf{Solution détaillée :}

\begin{enumerate}
\item \esp{\textbf{Separa} la basura en el recreo.} « Trie les déchets à la récréation. »
\textit{Justification :} \textbf{\esp{tú} affirmatif}. Verbe régulier en -AR $\to$ 3\textsuperscript{e} pers. sg du présent. Ton proche, éducatif.
\item \esp{\textbf{No tire} papeles al suelo.} « Ne jetez pas de papiers par terre. »
\textit{Justification :} \textbf{\esp{usted} négatif} = \esp{no} + subjonctif présent (3\textsuperscript{e} sg) : \esp{tire}. Registre formel, autorité.
\item \esp{\textbf{Plantemos} árboles en el barrio.} « Plantons des arbres dans le quartier. »
\textit{Justification :} \textbf{\esp{nosotros} affirmatif} = forme de subjonctif \esp{plantemos} (sans accent : la tonique tombe sur \esp{-te-}). Sens : « plantons ensemble ».
\item \esp{\textbf{Use} transporte público.} (ou \esp{\textbf{Usen}} pour \esp{ustedes}).
\textit{Justification :} \esp{usted/ustedes} pour un conseil impersonnel institutionnel. \textit{Alternative :} l'infinitif, fréquent sur les panneaux : \esp{Usar transporte público.}
\item \textbf{Espagne (\esp{vosotros}) :} \esp{\textbf{No malgastéis} el agua.} \textbf{Amérique latine / Cameroun (\esp{ustedes}) :} \esp{\textbf{No malgasten} el agua.}
\textit{Justification :} négatif $\to$ subjonctif : \esp{malgastéis} (2\textsuperscript{e} pl, accent sur \textbf{é}) / \esp{malgasten} (3\textsuperscript{e} pl).
\item \textbf{Pronoms :}
\begin{itemize}
\item Affirmatif \esp{tú} : \esp{\textbf{Recíclalo}.} (Accent sur \textbf{í} pour conserver la tonique : \esp{re-cí-cla-lo}.)
\item Négatif \esp{usted} : \esp{\textbf{No la contamine}.} (Pronom proclitique \esp{la} devant le subjonctif \esp{contamine}.)
\end{itemize}
\end{enumerate}

\end{exoresolu}

\begin{aretenir}[Check-list du slogan efficace]
\begin{itemize}
\item Verbe d'action à l'impératif (court, fort).
\item Vocabulaire précis (\esp{Reducir, Reutilizar, Reciclar -- las 3R}).
\item Image mentale évoquée (\esp{¡Apaga la luz!} $\to$ l'interrupteur).
\item Ponctuation espagnole : \esp{¡ ... !} (signe ouvrant ET fermant).
\end{itemize}
\end{aretenir}


\begin{methode}[Traduction Thème / Version : pièges grammaticaux du thème « environnement »]
\textbf{Objectif :} Réussir les exercices de traduction (thème = Fr $\to$ Esp ; version = Esp $\to$ Fr) en évitant les calques et en maîtrisant les points grammaticaux « stars » du Bac : \cle{por/para}, \cle{ser/estar}, \cle{subjonctif}, \cle{temps du passé}.
\begin{enumerate}
\item \textbf{Analyser la phrase source :} identifier le sujet, le verbe principal, les compléments, le temps, l'aspect, le mode.
\item \textbf{\esp{Por} vs \esp{Para} (le duel éternel) :}
\begin{itemize}
\item \esp{POR} : cause (\esp{por la contaminación}), moyen (\esp{por autobús}), durée (\esp{por dos horas}), échange, mouvement à travers.
\item \esp{PARA} : but/finalité (\esp{para proteger}), destinataire (\esp{para los jóvenes}), échéance (\esp{para mañana}), opinion (\esp{para mí}).
\end{itemize}
\item \textbf{\esp{Ser} vs \esp{Estar} :}
\begin{itemize}
\item \esp{SER} : identité, nature, origine, heure, possession, définition (\esp{El calentamiento \textbf{es} un problema grave}).
\item \esp{ESTAR} : état, localisation, résultat d'une action, action en cours (\esp{El río \textbf{está} contaminado}, \esp{Estamos \textbf{destruyendo} el bosque}).
\end{itemize}
\item \textbf{Subjonctif déclenché par :} volonté (\esp{querer que}), émotion (\esp{me preocupa que}), doute/négation (\esp{no creo que}), tournures impersonnelles (\esp{es necesario que}, \esp{es urgente que}), \esp{para que} (but).
\item \textbf{Temps du passé :} \esp{indefinido} (action ponctuelle, terminée, datée : \esp{Ayer plantamos árboles}) vs \esp{imperfecto} (description, habitude, contexte : \esp{Antes el aire era puro}).
\item \textbf{Relecture ciblée :} vérifier \textbf{uniquement} : accords, accents, \esp{por/para}, \esp{ser/estar}, subjonctif, temps du passé.
\end{enumerate}
\end{methode}

\begin{exoresolu}[Thème et Version « Spécial Environnement »]
\textbf{Énoncé :}
\textbf{A. Version (Espagnol $\to$ Français).} Traduisez fidèlement.
\esp{« Aunque las ciudades africanas crecen rápido, es fundamental que los gobiernos \textbf{inviertan} en energías limpias \textbf{para} que la calidad del aire \textbf{mejore}. La contaminación \textbf{no es} solo un problema estético; \textbf{está} matando a miles de personas al año. \textbf{Por} tanto, \textbf{debemos actuar} ya. »}

\textbf{B. Thème (Français $\to$ Espagnol).} Traduisez en espagnol correct.
\begin{enumerate}
\item « Il est nécessaire que nous réduisions nos déchets plastiques. »
\item « Cette rivière est polluée par les usines depuis des années. »
\item « Nous avons planté des arbres pour lutter contre l'érosion. » (action passée, terminée)
\item « Pour que la forêt survive, il ne faut pas la couper. »
\item « Pourquoi ne recyclez-vous pas ? -- Parce que c'est difficile. »
\end{enumerate}

\tcblower
\textbf{Solution détaillée :}

\textbf{A. Version.}
\begin{quote}
« Bien que les villes africaines grandissent vite, il est fondamental que les gouvernements \textbf{investissent} dans les énergies propres \textbf{afin que} la qualité de l'air \textbf{s'améliore}. La pollution \textbf{n'est} pas seulement un problème esthétique ; elle \textbf{est en train de tuer} des milliers de personnes par an. \textbf{Par conséquent}, \textbf{nous devons agir} dès maintenant. »
\end{quote}
\textit{Justifications grammaticales :}
\begin{itemize}
\item \esp{Aunque} + indicatif (\esp{crecen}) : fait réel, concession admise.
\item \esp{Es fundamental que} + \textbf{subjonctif} (\esp{inviertan}) : tournure impersonnelle de nécessité. \esp{Invertir} : verbe à changement radical \textbf{e $\to$ ie} (\esp{invierto}, \esp{inviertan}).
\item \esp{Para que} + \textbf{subjonctif} (\esp{mejore}) : but. \esp{Mejorar} : régulier en -AR.
\item \esp{No es} : \esp{ser} pour la nature/définition du problème.
\item \esp{Está matando} : \esp{estar} + gérondif, processus en cours. (Accepté aussi : \esp{mata}, fait général.)
\item \esp{Por tanto} : locution conjonctive = « par conséquent ». \esp{Debemos actuar} : obligation morale, présent de l'indicatif.
\end{itemize}

\textbf{B. Thème.}
\begin{enumerate}
\item \esp{\textbf{Es necesario que reduzcamos} nuestros residuos plásticos.}
\textit{Justif. :} \esp{es necesario que} + \textbf{subjonctif présent} (1\textsuperscript{re} pl). \esp{Reducir} : changement \textbf{c $\to$ zc} devant a/o (\esp{reduzco}, \esp{reduzcamos}).
\item \esp{\textbf{Este río está contaminado por las fábricas desde hace años.}}
\textit{Justif. :} \esp{está contaminado} : \textbf{\esp{estar}} + participe passé = état résultant d'une action. \esp{Por} introduit l'agent. \esp{Desde hace años} = « depuis des années » (présent + durée écoulée).
\item \esp{\textbf{Plantamos árboles para combatir la erosión.}}
\textit{Justif. :} \esp{plantamos} : \textbf{indefinido} (action ponctuelle, terminée). \textit{Piège :} l'imparfait \esp{plantábamos} signifierait « nous plantions » (habitude). \esp{Para} + infinitif = but.
\item \esp{\textbf{Para que el bosque sobreviva, no hay que talarlo.}}
\textit{Justif. :} \esp{para que} + \textbf{subjonctif} : \esp{sobreviva} (\esp{sobrevivir}, verbe régulier en -IR, \textbf{sans} changement radical). \esp{No hay que} + infinitif = obligation impersonnelle. \esp{Talar} = abattre (des arbres). \textbf{Attention au genre :} \esp{el bosque} est \textbf{masculin} $\to$ pronom \esp{lo} (\esp{talarlo}), et non \esp{la}.
\item \esp{-- \textbf{¿Por qué no reciclan?} (ou \esp{¿Por qué no recicláis?}) -- \textbf{Porque es difícil.}}
\textit{Justif. :} \esp{por qué} (deux mots, accent) = question ; \esp{porque} (un mot, sans accent) = réponse/cause. \esp{Es difícil} : \esp{ser} pour une caractéristique intrinsèque.
\end{enumerate}
\end{exoresolu}

\begin{methode}[Produire une affiche de sensibilisation : de l'idée au produit fini]
\textbf{Objectif :} Concevoir une affiche complète (titre, visuel décrit, slogan principal, 3 actions concrètes, signature) en espagnol correct, adaptée à un public cible (lycéens, marché, mairie).
\begin{enumerate}
\item \textbf{Analyser la commande :} sujet précis (ex. : \esp{residuos plásticos en el río Wouri}), public cible (ex. : \esp{comerciantes del mercado Central}), ton (alarmiste / positif / humoristique), format.
\item \textbf{Structurer l'affiche (zonage visuel) :}
\begin{enumerate}
\item \textbf{TITRE (\esp{titular}) :} court, choc, nominal ou interrogatif. \esp{¡Río Wouri: SOS!} / \esp{¿Dónde va tu bolsa?}
\item \textbf{CORPS CENTRAL (visuel + slogan) :} verbe à l'impératif. \esp{¡No tires plástico!} / \esp{¡Recicla tu botella!}
\item \textbf{ACTIONS CONCRÈTES (3 max) :} puces ou numérotation. Infinitif (consigne générale) ou impératif \esp{ustedes}. \esp{1. Separar la basura. 2. Usar bolsas de tela. 3. Denunciar los vertidos ilegales.}
\item \textbf{PIED D'AFFICHE (signature) :} logo ONG / mairie / lycée, hashtag \esp{\#YaoundéLimpio}, éventuel QR code.
\end{enumerate}
\item \textbf{Relire :} orthographe (accents !), accords, clarté du message, impact visuel (mots en \textbf{gras}, MAJUSCULES pour le slogan).
\end{enumerate}
\end{methode}

\begin{exoresolu}[Création d'affiche : « Salvemos el Manglar de Douala »]
\textbf{Énoncé :} Vous êtes membre du club « Jóvenes por el Clima » au Lycée de Douala. Concevez le \textbf{contenu textuel complet} d'une affiche A3 pour sensibiliser les pêcheurs et les riverains à la protection de la mangrove. Contraintes : titre, slogan principal (impératif), 3 actions (infinitif), message final (subjonctif avec \esp{para que}), signature. Rédigez en espagnol.

\tcblower
\textbf{Proposition de production écrite (modèle type Bac).}

\begin{textebox}[Contenu de l'affiche -- Texte original]
\textbf{¡MANGLAR VIVO, PUEBLO VIVO!} \\[0.5em]
\esp{\textbf{¡PROTEJAMOS NUESTRAS RAÍCES!}} \\[1em]
\esp{$\bullet$ \quad Evitar la tala indiscriminada de mangle.} \\
\esp{$\bullet$ \quad No arrojar redes ni plásticos al agua.} \\
\esp{$\bullet$ \quad Denunciar la contaminación industrial.} \\[1em]
\esp{\textbf{Es vital que todos colaboremos para que el ecosistema sobreviva.}} \\[1em]
\esp{Club « Jóvenes por el Clima » -- Lycée de Douala} \\
\esp{\#ManglarDouala \quad \#CamerúnVerde}
\end{textebox}

\textit{Traduction :} « Mangrove vivante, peuple vivant ! Protégeons nos racines ! -- Éviter l'abattage aveugle du palétuvier. -- Ne pas jeter de filets ni de plastiques dans l'eau. -- Dénoncer la pollution industrielle. -- Il est vital que nous collaborions tous pour que l'écosystème survive. -- Club "Jeunes pour le Climat" -- Lycée de Douala. »

\textbf{Analyse linguistique du modèle :}
\begin{itemize}
\item \textbf{Titre :} nominal, parallélisme \esp{vivo/vivo}, vocabulaire fort (\esp{raíces} = racines de la mangrove et identité).
\item \textbf{Slogan :} \esp{¡Protejamos!} $\to$ \textbf{\esp{nosotros} affirmatif} (engagement collectif, inclut pêcheurs et lycéens). \esp{Proteger} : changement \textbf{g $\to$ j} devant a/o (\esp{protejo}, \esp{protejamos}).
\item \textbf{Actions :} \textbf{infinitif} (\esp{evitar}, \esp{no arrojar}, \esp{denunciar}) : la norme pour les listes de consignes, neutre et impersonnel. Lexique précis : \esp{tala indiscriminada}, \esp{arrojar} (jeter volontairement), \esp{contaminación industrial}.
\item \textbf{Message final :} structure complexe \textbf{subjonctif + \esp{para que} + subjonctif}.
\begin{itemize}
\item \esp{Es vital que} (impersonnel de nécessité) $\to$ \esp{colaboremos} (subjonctif présent, 1\textsuperscript{re} pl de \esp{colaborar}).
\item \esp{Para que} (but) $\to$ \esp{sobreviva} (subjonctif présent, 3\textsuperscript{e} sg de \esp{sobrevivir}, verbe régulier en -IR).
\end{itemize}
\item \textbf{Signature :} identité claire, hashtags pour la visibilité sur les réseaux sociaux.
\end{itemize}

\end{exoresolu}

\begin{exemplebox}[Variante « ton humoristique / jeunes »]
\esp{¡NO SEAS « PLASTIC-ADICTO »!} \\
\esp{¡Rompe el ciclo! Usa tu botella. \#RefillDouala} \\
\esp{Liceo Bilingüe de Bonamoussadi}
\end{exemplebox}
\textit{Traduction :} « Ne sois pas "plastique-addict" ! Brise le cycle ! Utilise ta gourde. »


\begin{attention}[Les 5 erreurs qui coûtent des points au Bac (Spécial Environnement)]
\begin{enumerate}
\item \textbf{Accents oubliés ou mal placés sur les mots-clés du thème :} \esp{sequía} (accent sur \textbf{í}), \esp{energía} (accent sur \textbf{í}), \esp{atmósfera} (accent sur \textbf{ó}, esdrújulo), \esp{carbono} (accent sur \textbf{ó}, esdrújulo), \esp{contaminación} (accent sur \textbf{ó}). En revanche, \textbf{pas d'accent} sur \esp{calentamiento}, \esp{medio ambiente}, \esp{renovables}. \textit{Règles :} mots aigus terminés par voyelle, -n ou -s = accent ; mots esdrújulos = accent toujours ; hiatus avec voyelle fermée tonique (\esp{i}, \esp{u}) = accent.
\item \textbf{Faux ami \esp{contaminación} / \esp{infección} :} \esp{contaminación} = pollution (air, eau, sol, bruit) ; \esp{infección} = infection (médical : virus, bactéries). Ne jamais dire \esp{infección del aire}.
\item \textbf{Calque \esp{por} / \esp{para} :} \esp{Lucho \textbf{por} el medio ambiente} (cause, motif). \esp{Lucho \textbf{para} salvar el planeta} (but). \esp{Es importante \textbf{para} mí} (opinion). \esp{Lo hago \textbf{por} ti} (pour toi, à cause de toi). \textit{Astuce :} \esp{para} = flèche vers l'avant (but, destination) ; \esp{por} = flèche vers l'arrière (cause) ou à travers (moyen, lieu).
\item \textbf{\esp{Ser} / \esp{Estar} avec les adjectifs d'état :} \esp{El clima \textbf{es} tropical} (caractère permanent). \esp{El clima \textbf{está} cambiante} (état actuel, évolution). \esp{El río \textbf{está} contaminado} (état résultant d'une action). \esp{La contaminación \textbf{es} grave} (nature du problème) vs \esp{La ciudad \textbf{está} muy contaminada} (état actuel).
\item \textbf{Subjonctif oublié après \esp{es necesario que}, \esp{para que}, \esp{no creer que} :} \esp{Es necesario que \textbf{reduzcamos}} (subj.) $\neq$ \esp{reducimos} (ind.). \esp{Para que \textbf{mejore}} (subj.) $\neq$ \esp{mejora} (ind.). \textit{Réflexe :} repérer le déclencheur (volonté, émotion, impersonnel, doute, but) $\to$ \textbf{subjonctif obligatoire}.
\end{enumerate}
\end{attention}

\newpage

\coursec{Exercices}

\courssub{Je vérifie mes connaissances}

\begin{exercice}[QCM : Lexique et grammaire de l'environnement]
Chaque question ne possède qu'une seule réponse correcte. Choisissez la bonne option.
\begin{enumerate}
    \item Comment dit-on « la couche d'ozone » en espagnol ?
    \begin{enumerate}
        \item La capa de ozono
        \item El capa de ozono
        \item La capa de ozona
    \end{enumerate}
    \item Quel verbe signifie « recycler » ?
    \begin{enumerate}
        \item Reducir
        \item Reutilizar
        \item Reciclar
    \end{enumerate}
    \item Complétez l'impératif affirmatif pour \emph{tú} : \esp{\_\_\_\_\_\_\_\_\_\_ (proteger) los bosques.}
    \begin{enumerate}
        \item Protege
        \item Proteja
        \item Protegemos
    \end{enumerate}
    \item Le mot \cle{sequía} correspond à quel mot français ?
    \begin{enumerate}
        \item Inondation
        \item Sécheresse
        \item Pollution
    \end{enumerate}
    \item Quelle forme est l'impératif négatif pour \emph{vosotros} du verbe \emph{contaminar} ?
    \begin{enumerate}
        \item No contaminéis
        \item No contaminan
        \item No contamineis
    \end{enumerate}
\end{enumerate}
\end{exercice}

\begin{exercice}[Vrai ou Faux justifié : Lecture rapide]
Lisez le texte suivant et dites si les affirmations sont vraies (V) ou fausses (F). Justifiez chaque réponse par une phrase du texte.
\begin{textebox}[Extrait : La situación del río Wouri]
La contaminación del río Wouri en Douala es un problema grave. Las fábricas vierten residuos tóxicos sin tratamiento y la población tira basura plástica. Esto provoca la muerte de peces y enfermedades en las comunidades ribereñas. Es urgente implementar plantas de tratamiento y educar a los ciudadanos sobre el reciclaje.
\end{textebox}
\begin{enumerate}
    \item El río Wouri está limpio y sano.
    \item Las fábricas tratan sus residuos antes de verterlos.
    \item La basura plástica afecta la fauna del río.
    \item La solución pasa solo por la acción del gobierno.
\end{enumerate}
\end{exercice}

\begin{exercice}[Vocabulaire : Association et définition]
Associez chaque mot espagnol à sa définition française, puis écrivez une phrase originale en espagnol pour chacun.
\begin{center}
\begin{tabularx}{\linewidth}{|X|X|}\hline
\rowcolor{popblueL} \textbf{Español} & \textbf{Français} \\ \hline
1. Calentamiento global & A. Action de planter des arbres \\ \hline
2. Deforestación & B. Augmentation de la température terrestre \\ \hline
3. Energías renovables & C. Disparition des forêts \\ \hline
4. Reforestación & D. Énergies inépuisables (solaire, éolien) \\ \hline
5. Huella de carbono & E. Quantité de CO2 émise par une activité \\ \hline
\end{tabularx}
\end{center}
\end{exercice}

\begin{exercice}[Conjugaison : L'impératif de sensibilisation]
Conjuguez les verbes entre parenthèses à l'impératif affirmatif (pour \emph{tú}) ou négatif (pour \emph{ustedes}) selon le contexte.
\begin{enumerate}
    \item \esp{\_\_\_\_\_\_\_\_\_\_ (ahorrar) agua! El grifo gotea.} (Affirmatif \emph{tú})
    \item \esp{\_\_\_\_\_\_\_\_\_\_ (no / tirar) basura al suelo!} (Négatif \emph{ustedes})
    \item \esp{\_\_\_\_\_\_\_\_\_\_ (usar) transporte público para reducir el smog.} (Affirmatif \emph{tú})
    \item \esp{\_\_\_\_\_\_\_\_\_\_ (no / gastar) electricidad innecesariamente.} (Négatif \emph{ustedes})
    \item \esp{\_\_\_\_\_\_\_\_\_\_ (plantar) un árbol hoy.} (Affirmatif \emph{tú})
    \item \esp{\_\_\_\_\_\_\_\_\_\_ (separar) los residuos: orgánicos, plástico, papel.} (Affirmatif \emph{vosotros})
    \item \esp{\_\_\_\_\_\_\_\_\_\_ (no / contaminar) el aire con humos tóxicos.} (Négatif \emph{ustedes})
    \item \esp{\_\_\_\_\_\_\_\_\_\_ (cuidar) el planeta, es nuestra casa.} (Affirmatif \emph{tú})
\end{enumerate}
\end{exercice}

\courssub{Je m'entraîne}

\begin{exercice}[Traduction : Du français vers l'espagnol (Thème)]
Traduisez les phrases suivantes en espagnol en mobilisant le lexique de la leçon et l'impératif.
\begin{enumerate}
    \item Nous devons réduire la pollution de l'air dans nos villes.
    \item Le réchauffement climatique cause des sécheresses prolongées.
    \item Recyclez le papier et le plastique ! (Adresse à un ami)
    \item Ne coupez pas les arbres de la forêt ! (Adresse à un groupe / \emph{ustedes})
    \item Les énergies renouvelables protègent la couche d'ozone.
    \item Il y a beaucoup de pollution plastique dans l'océan Atlantique.
    \item Économisons l'eau : fermez le robinet ! (Nous / \emph{nosotros} + infinitif d'ordre ou impératif)
    \item Le gouvernement doit arrêter la déforestation illégale.
\end{enumerate}
\end{exercice}

\begin{exercice}[Transformation : De l'infinitif au slogan (Impératif)]
Transformez chaque infinitif en deux slogans : un à l'impératif affirmatif (\emph{tú}) et un à l'impératif négatif (\emph{ustedes}). Exemple : \esp{proteger los bosques $\rightarrow$ ¡Protege los bosques! / ¡No destruyan los bosques!}
\begin{enumerate}
    \item ahorrar energía
    \item reciclar botellas
    \item usar bolsas de tela
    \item plantar árboles nativos
    \item limpiar las playas
    \item reducir el consumo de plástico
    \item respetar la naturaleza
    \item denunciar la contaminación industrial
    \item educar a los niños
    \item cambiar los hábitos
\end{enumerate}
\end{exercice}

\begin{exercice}[Texte à trous : La pollution plastique]
Complétez le texte avec les mots de la liste ci-dessous. Adaptez le genre et le nombre si nécessaire.
\begin{center}
\textbf{Liste :} \esp{contaminación, océanos, microplásticos, fauna, reciclaje, degradar, conciencia, residuos, sequía, renovables}
\end{center}
\begin{textebox}[La amenaza invisible]
La \esp{\_\_\_\_\_\_\_\_\_\_} por plásticos es una crisis global. Cada año, millones de toneladas de \esp{\_\_\_\_\_\_\_\_\_\_} llegan a los \esp{\_\_\_\_\_\_\_\_\_\_}. Los plásticos tardan siglos en \esp{\_\_\_\_\_\_\_\_\_\_} y se convierten en \esp{\_\_\_\_\_\_\_\_\_\_} que entran en la cadena alimentaria. La \esp{\_\_\_\_\_\_\_\_\_\_} marina sufre: tortugas, peces y aves mueren por ingestión. La solución requiere \esp{\_\_\_\_\_\_\_\_\_\_} ciudadana y un \esp{\_\_\_\_\_\_\_\_\_\_} efectivo. También debemos apostar por energías \esp{\_\_\_\_\_\_\_\_\_\_} y reducir nuestra huella. La \esp{\_\_\_\_\_\_\_\_\_\_} en algunas zonas empeora por la falta de gestión de aguas.
\end{textebox}
\end{exercice}

\begin{exercice}[Appariement : Slogans et actions ciblées]
Reliez chaque slogan à l'action concrète qu'il promeut et au lieu d'affichage idéal (école, marché, bureau, arrêt de bus, domicile).
\begin{center}
\begin{tabularx}{\linewidth}{|X|X|X|}\hline
\rowcolor{popblueL} \textbf{Slogan} & \textbf{Action} & \textbf{Lieu idéal} \\ \hline
¡Apaga la luz al salir! & A. Réduire les déchets plastiques & 1. École / Bureau \\ \hline
¡Lleva tu propia bolsa! & B. Économiser l'électricité & 2. Marché / Supermarché \\ \hline
¡Cierra el grifo! & C. Éviter le gaspillage d'eau & 3. Domicile / Salle de bain \\ \hline
¡Separa la basura! & D. Trier pour recycler & 4. Cuisine / Bureau \\ \hline
¡Usa la bici! & E. Réduire émissions CO2 & 5. Arrêt de bus / Rue \\ \hline
\end{tabularx}
\end{center}
Écrivez ensuite 3 nouveaux slogans originaux pour 3 autres lieux (ex. : bord de rivière, station-service, cantine).
\end{exercice}

\courssub{Je me prépare au Bac}

\begin{exercice}[Compréhension écrite type Bac : « El plástico ahoga el Golfo de Guinea »]
Lisez attentivement le texte suivant et répondez aux questions en espagnol. Respectez la limite de mots indiquée.
\begin{textebox}[El plástico ahoga el Golfo de Guinea]
Las playas de Kribi y Limbe, antes paradisíacas, muestran hoy una cara triste: montañas de botellas, bolsas y redes de pesca cubren la arena. La corriente del Golfo de Guinea trae residuos de todo el Atlántico, pero gran parte proviene de la gestión deficiente de desechos en Douala y Yaoundé. Los pescadores capturan más plástico que peces; las tortugas marinas confunden bolsas con medusas y mueren asfixiadas. Los microplásticos entran en la cadena alimentaria y llegan a nuestros platos.

Ante esta emergencia, la ONG \emph{Mares Limpios Camerún} organiza limpiezas mensuales y talleres en escuelas. « La clave es la educación y la economía circular », explica su directora, Ana Mbarga. « Debemos exigir leyes que prohíban plásticos de un solo uso y fomentar envases retornables ». El gobierno anunció un plan nacional para 2030, pero los activistas piden acción inmediata: « ¡No hay planeta B! », gritan los jóvenes en las marchas climáticas.
\end{textebox}
\begin{enumerate}
    \item \textbf{Vrai ou Faux ? Justifiez par une citation du texte.}
    \begin{enumerate}
        \item La pollution des plages vient uniquement des déchets locaux.
        \item Les pêcheurs attrapent davantage de plastique que de poisson.
        \item Le gouvernement a déjà résolu le problème avec le plan 2030.
    \end{enumerate}
    \item \textbf{Questions ouvertes (réponse en 30 mots max en espagnol).}
    \begin{enumerate}
        \item ¿Qué consecuencias tiene la contaminación plástica en la fauna marina según el texto?
        \item ¿Qué dos medidas concretas propone la directora de la ONG?
        \item ¿Cuál es el lema de los jóvenes en las marchas y qué significa?
    \end{enumerate}
    \item \textbf{Lexique en contexte.}
    \begin{enumerate}
        \item Expliquez en espagnol le sens de \cle{economía circular} d'après le texte.
        \item Trouvez dans le texte un synonyme de \esp{basura} (déchet).
        \item Donnez l'infinitif des verbes conjugués : \esp{trae, confunden, llegan, exige, fomentar}.
    \end{enumerate}
    \item \textbf{Expression personnelle (50 mots env.).} Selon vous, quelle action citoyenne est la plus urgente pour sauver le littoral camerounais ? Justifiez.
\end{enumerate}
\end{exercice}

\begin{exercice}[Thème et Version : Traduction spécialisée]
\textbf{A. Version (Français $\rightarrow$ Espagnol).} Traduisez en espagnol.
\begin{enumerate}
    \item La déforestation en Amazonie accélère le réchauffement climatique.
    \item Nous devons interdire les plastiques à usage unique dès maintenant.
    \item Les énergies renouvelables sont la solution pour l'avenir de la planète.
    \item N'achetez pas de produits avec trop d'emballages ! (Impératif \emph{ustedes})
    \item Les enfants apprennent à recycler à l'école primaire.
\end{enumerate}

\textbf{B. Thème (Espagnol $\rightarrow$ Français).} Traduisez en français.
\begin{textebox}[Fragmento para traducir]
« La sequía prolongada en el norte de Camerún amenaza la seguridad alimentaria. Los agricultores no pueden cultivar maíz ni mijo. Los pozos se secan y el ganado muere. Es vital construir presas y usar riego por goteo. Además, la reforestación con especies autóctonas frena el avance del desierto. Cada ciudadano puede actuar: plantar un árbol, no quemar la basura, denunciar la tala ilegal. El cambio empieza en tu barrio. »
\end{textebox}
\end{exercice}

\begin{exercice}[Sujet de rédaction type Bac : Production écrite]
\textbf{Consigne :} Vous participez à un concours d'affiches organisé par le Ministère de l'Environnement du Cameroun. Rédigez le texte complet d'une affiche de sensibilisation intitulée : \esp{« ¿Cómo proteger el medio ambiente en tu ciudad? »}

\textbf{Contraintes :}
\begin{itemize}
    \item Longueur : 15 lignes environ (120--150 mots).
    \item Structure : Accroche (titre/slogan principal) $\rightarrow$ Diagnostic court (problèmes locaux) $\rightarrow$ 5 à 6 mesures concrètes à l'impératif (affirmatif et négatif, \emph{tú} ou \emph{ustedes}) $\rightarrow$ Slogan final mobilisateur.
    \item Lexique imposé (au moins 5) : \cle{contaminación}, \cle{reciclaje}, \cle{energías renovables}, \cle{deforestación}, \cle{sequía}, \cle{capa de ozono}, \cle{huella de carbono}.
    \item Grammaire : Minimum \textbf{3 impératifs} différents (affirmatifs et/ou négatifs). Utilisez des connecteurs logiques (\esp{por tanto, además, es urgente que + subjuntivo}).
    \item Contexte : Ancrage camerounais (Douala, Yaoundé, Wouri, Mont Cameroun, marchés, bendskins, etc.) bienvenu.
\end{itemize}
\textbf{Éléments de réflexion :} Pensez aux transports (bendskins, bus), aux déchets (marchés, caniveaux), à l'énergie (générateurs, solaire), à la végétation (reboisement urbain).
\end{exercice}

\newpage

\coursec{Corrigés des exercices}

\begin{corrige}[Exercice 1 — QCM : Lexique et grammaire de l'environnement]
\begin{enumerate}
    \item \textbf{Réponse a :} \esp{La capa de ozono}. \esp{Capa} est féminin (\esp{la capa}), et le mot correct est \esp{ozono} (masculin, sans -a finale).
    \item \textbf{Réponse c :} \esp{Reciclar} = recycler. \esp{Reducir} = réduire ; \esp{reutilizar} = réutiliser. Retenez la triade écologique : \esp{reducir, reutilizar, reciclar}.
    \item \textbf{Réponse a :} \esp{¡Protege los bosques!} Impératif affirmatif \emph{tú} = 3\up{e} personne du singulier du présent : \esp{protege}. \esp{Proteja} est le subjonctif (impératif \emph{usted}) ; \esp{protegemos} est l'indicatif \emph{nosotros}.
    \item \textbf{Réponse b :} \esp{Sequía} = sécheresse. « Inondation » se dit \esp{inundación}.
    \item \textbf{Réponse a :} \esp{¡No contaminéis!} L'impératif négatif se construit avec le subjonctif présent : \esp{no contaminéis} (avec accent sur le \emph{é}). La forme \esp{contamineis} sans accent est fautive ; \esp{no contaminan} est un indicatif.
\end{enumerate}
\textbf{Point de vigilance :} à l'impératif négatif \emph{vosotros}, les verbes en \esp{-ar} prennent \esp{-éis} et les verbes en \esp{-er/-ir} prennent \esp{-áis} : \esp{no contaminéis} mais \esp{no bebáis}.
\end{corrige}

\begin{corrige}[Exercice 2 — Vrai ou Faux justifié : Lecture rapide]
\begin{enumerate}
    \item \textbf{F.} Le texte dit : \esp{« La contaminación del río Wouri en Douala es un problema grave »}. Le fleuve est donc pollué, pas propre.
    \item \textbf{F.} Le texte dit : \esp{« Las fábricas vierten residuos tóxicos sin tratamiento »}. « Sans traitement » : les usines ne traitent pas leurs déchets.
    \item \textbf{V.} Le texte dit : \esp{« la población tira basura plástica. Esto provoca la muerte de peces »}. Le plastique tue les poissons, donc il affecte la faune.
    \item \textbf{F.} Le texte propose deux solutions complémentaires : \esp{« implementar plantas de tratamiento y educar a los ciudadanos sobre el reciclaje »}. L'éducation des citoyens montre que la solution ne dépend pas que du gouvernement.
\end{enumerate}
\textbf{Méthode :} au Bac, une justification sans citation exacte du texte ne rapporte pas tous les points. Recopiez toujours le segment justificatif entre guillemets.
\end{corrige}

\begin{corrige}[Exercice 3 — Vocabulaire : Association et définition]
\textbf{Associations :} 1-B ; 2-C ; 3-D ; 4-A ; 5-E.
\begin{center}
\begin{tabularx}{\linewidth}{|X|X|}\hline
\rowcolor{popblueL} \textbf{Español} & \textbf{Français} \\ \hline
1. Calentamiento global & B. Augmentation de la température terrestre \\ \hline
2. Deforestación & C. Disparition des forêts \\ \hline
3. Energías renovables & D. Énergies inépuisables (solaire, éolien) \\ \hline
4. Reforestación & A. Action de planter des arbres \\ \hline
5. Huella de carbono & E. Quantité de CO2 émise par une activité \\ \hline
\end{tabularx}
\end{center}
\textbf{Phrases originales possibles (exemples) :}
\begin{enumerate}
    \item \esp{El calentamiento global derrite los glaciares y sube el nivel del mar.} « Le réchauffement climatique fait fondre les glaciers et monte le niveau de la mer. »
    \item \esp{La deforestación destruye el hábitat de los gorilas en el sur de Camerún.} « La déforestation détruit l'habitat des gorilles dans le sud du Cameroun. »
    \item \esp{Las energías renovables, como la solar, no contaminan el aire.} « Les énergies renouvelables, comme le solaire, ne polluent pas l'air. »
    \item \esp{La reforestación de las colinas de Yaoundé mejora la calidad del aire.} « Le reboisement des collines de Yaoundé améliore la qualité de l'air. »
    \item \esp{Viajar en bicicleta reduce tu huella de carbono.} « Voyager à vélo réduit ton empreinte carbone. »
\end{enumerate}
\end{corrige}

\begin{corrige}[Exercice 4 — Conjugaison : L'impératif de sensibilisation]
\begin{enumerate}
    \item \esp{¡Ahorra agua! El grifo gotea.} — Impératif \emph{tú} de \esp{ahorrar} : \esp{ahorra}.
    \item \esp{¡No tiren basura al suelo!} — Impératif négatif \emph{ustedes} = subjonctif présent : \esp{no tiren}.
    \item \esp{Usa transporte público para reducir el smog.} — \esp{usa} (\emph{tú}).
    \item \esp{No gasten electricidad innecesariamente.} — \esp{no gasten} (\emph{ustedes}).
    \item \esp{Planta un árbol hoy.} — \esp{planta} (\emph{tú}).
    \item \esp{Separad los residuos: orgánicos, plástico, papel.} — Impératif affirmatif \emph{vosotros} : infinitif moins \emph{-r} plus \emph{-d} : \esp{separad}.
    \item \esp{No contaminen el aire con humos tóxicos.} — \esp{no contaminen} (\emph{ustedes}).
    \item \esp{Cuida el planeta, es nuestra casa.} — \esp{cuida} (\emph{tú}).
\end{enumerate}
\textbf{Point de vigilance :} ne confondez pas \esp{tira} (impératif \emph{tú}) et \esp{tiren} (impératif négatif \emph{ustedes}). Les verbes en \esp{-er/-ir} ont un impératif \emph{tú} en \emph{-e} : \esp{come, escribe, protege}.
\end{corrige}

\begin{corrige}[Exercice 5 — Traduction : Du français vers l'espagnol (Thème)]
\begin{enumerate}
    \item \esp{Debemos reducir la contaminación del aire en nuestras ciudades.}
    \item \esp{El calentamiento global causa sequías prolongadas.} (On accepte aussi \esp{provoca}.)
    \item \esp{¡Recicla el papel y el plástico!} — Impératif \emph{tú} : \esp{recicla}.
    \item \esp{¡No corten los árboles del bosque!} — Impératif négatif \emph{ustedes} : subjonctif \esp{no corten}. Attention : \esp{cortar} est régulier, pas de changement de voyelle.
    \item \esp{Las energías renovables protegen la capa de ozono.}
    \item \esp{Hay mucha contaminación plástica en el océano Atlántico.}
    \item \esp{Ahorremos agua: ¡cierra el grifo!} — Impératif \emph{nosotros} = subjonctif : \esp{ahorremos} ; puis \esp{cierra} (\emph{tú}, verbe à changement \emph{e $\rightarrow$ ie}).
    \item \esp{El gobierno debe detener la deforestación ilegal.} (On accepte \esp{parar} ou \esp{frenar}.)
\end{enumerate}
\textbf{Points de vigilance :} « pollution » se traduit par \esp{contaminación} (et non \esp{polución}, qui existe mais est moins courant) ; « réchauffement climatique » = \esp{calentamiento global} ou \esp{cambio climático} ; n'oubliez pas les accents : \esp{sequías, Atlántico, árboles}.
\end{corrige}

\begin{corrige}[Exercice 6 — Transformation : De l'infinitif au slogan (Impératif)]
\begin{enumerate}
    \item \esp{¡Ahorra energía! / ¡No desperdicien energía!}
    \item \esp{¡Recicla botellas! / ¡No tiren botellas al suelo!}
    \item \esp{¡Usa bolsas de tela! / ¡No compren bolsas de plástico!}
    \item \esp{¡Planta árboles nativos! / ¡No talen árboles nativos!}
    \item \esp{¡Limpia las playas! / ¡No ensucien las playas!}
    \item \esp{¡Reduce el consumo de plástico! / ¡No aumenten el consumo de plástico!}
    \item \esp{¡Respeta la naturaleza! / ¡No destruyan la naturaleza!}
    \item \esp{¡Denuncia la contaminación industrial! / ¡No ignoren la contaminación industrial!}
    \item \esp{¡Educa a los niños! / ¡No descuiden la educación de los niños!}
    \item \esp{¡Cambia los hábitos! / ¡No mantengan hábitos contaminantes!}
\end{enumerate}
\textbf{Justification grammaticale :} l'impératif affirmatif \emph{tú} reprend la 3\up{e} personne du singulier du présent (\esp{ahorra, recicla, usa, planta, limpia, respeta, denuncia, educa}) ; attention aux verbes à changement vocalique : \esp{reducir $\rightarrow$ reduce}, \esp{cambiar $\rightarrow$ cambia} (ici \esp{cambiar} est régulier à l'impératif). L'impératif négatif \emph{ustedes} = \esp{no} + subjonctif présent en \esp{-en} : \esp{no tiren, no compren, no talen, no destruyan} (attention : \esp{destruir $\rightarrow$ destruyan}, avec \emph{y} intercalée).
\end{corrige}

\begin{corrige}[Exercice 7 — Texte à trous : La pollution plastique]
\begin{enumerate}
    \item \esp{contaminación} : \esp{La contaminación por plásticos es una crisis global.}
    \item \esp{residuos} : \esp{millones de toneladas de residuos} (pluriel, masculin).
    \item \esp{océanos} : \esp{llegan a los océanos}.
    \item \esp{degradar} : \esp{tardan siglos en degradar(se)} — infinitif après \esp{en}. La forme pronominale \esp{degradarse} est aussi correcte.
    \item \esp{microplásticos} : \esp{se convierten en microplásticos}.
    \item \esp{fauna} : \esp{La fauna marina sufre}.
    \item \esp{conciencia} : \esp{requiere conciencia ciudadana}.
    \item \esp{reciclaje} : \esp{un reciclaje efectivo}.
    \item \esp{renovables} : \esp{energías renovables} (adjectif pluriel).
    \item \esp{sequía} : \esp{La sequía en algunas zonas empeora}.
\end{enumerate}
\textbf{Texte complété (traduction) :} « La pollution par les plastiques est une crise mondiale. Chaque année, des millions de tonnes de déchets arrivent dans les océans. Les plastiques mettent des siècles à se dégrader et se transforment en microplastiques qui entrent dans la chaîne alimentaire. La faune marine souffre : tortues, poissons et oiseaux meurent par ingestion. La solution exige une prise de conscience citoyenne et un recyclage efficace. Nous devons aussi miser sur les énergies renouvelables et réduire notre empreinte. La sécheresse dans certaines zones s'aggrave faute de gestion des eaux. »
\end{corrige}

\begin{corrige}[Exercice 8 — Appariement : Slogans et actions ciblées]
\textbf{Appariements :}
\begin{center}
\begin{tabularx}{\linewidth}{|X|X|X|}\hline
\rowcolor{popblueL} \textbf{Slogan} & \textbf{Action} & \textbf{Lieu idéal} \\ \hline
¡Apaga la luz al salir! & B. Économiser l'électricité & 1. École / Bureau \\ \hline
¡Lleva tu propia bolsa! & A. Réduire les déchets plastiques & 2. Marché / Supermarché \\ \hline
¡Cierra el grifo! & C. Éviter le gaspillage d'eau & 3. Domicile / Salle de bain \\ \hline
¡Separa la basura! & D. Trier pour recycler & 4. Cuisine / Bureau \\ \hline
¡Usa la bici! & E. Réduire émissions CO2 & 5. Arrêt de bus / Rue \\ \hline
\end{tabularx}
\end{center}
\textbf{Trois slogans originaux (exemples) :}
\begin{itemize}
    \item Bord de rivière (Wouri) : \esp{¡El río no es un basurero: no tires nada al agua!} « La rivière n'est pas une poubelle : ne jette rien dans l'eau ! »
    \item Station-service : \esp{¡Apaga el motor mientras esperas: menos humo, más aire limpio!} « Éteins le moteur en attendant : moins de fumée, plus d'air pur ! »
    \item Cantine scolaire : \esp{¡Termina tu plato: no desperdicies comida!} « Finis ton assiette : ne gaspille pas la nourriture ! »
\end{itemize}
\textbf{Point de vigilance :} \esp{apaga} (\esp{apagar}), \esp{cierra} (\esp{cerrar}, \emph{e $\rightarrow$ ie}), \esp{lleva}, \esp{separa}, \esp{usa} sont des impératifs \emph{tú} réguliers ou à changement vocalique ; vérifiez toujours la voyelle radical.
\end{corrige}

\begin{corrige}[Exercice 9 — Compréhension écrite type Bac : « El plástico ahoga el Golfo de Guinea »]
\textbf{Barème indicatif (sur 20) :} V/F justifié 3 $\times$ 2 = 6 pts ; questions ouvertes 3 $\times$ 2 = 6 pts ; lexique 3 $\times$ 1,5 = 4,5 pts (arrondi 4) ; expression personnelle 4 pts.

\textbf{1. Vrai ou Faux :}
\begin{enumerate}
    \item \textbf{F.} \esp{« La corriente del Golfo de Guinea trae residuos de todo el Atlántico »} : une partie des déchets vient de l'Atlantique entier, pas seulement du local.
    \item \textbf{V.} \esp{« Los pescadores capturan más plástico que peces »}.
    \item \textbf{F.} \esp{« El gobierno anunció un plan nacional para 2030, pero los activistas piden acción inmediata »} : le plan est seulement annoncé, le problème n'est pas résolu.
\end{enumerate}

\textbf{2. Questions ouvertes (réponses modèles) :}
\begin{enumerate}
    \item \esp{Las tortugas confunden bolsas con medusas y mueren asfixiadas; los microplásticos entran en la cadena alimentaria y llegan a nuestros platos.} (24 mots)
    \item \esp{Propone exigir leyes que prohíban los plásticos de un solo uso y fomentar los envases retornables, además de educar a la población.} (22 mots)
    \item \esp{El lema es « ¡No hay planeta B! ». Significa que la Tierra es nuestro único hogar y que debemos protegerla ahora.} (23 mots)
\end{enumerate}

\textbf{3. Lexique en contexte :}
\begin{enumerate}
    \item \esp{La economía circular es un sistema donde los residuos se reutilizan y reciclan para no desperdiciar recursos, por ejemplo con envases retornables.}
    \item Synonyme de \esp{basura} dans le texte : \esp{residuos} (ou \esp{desechos}).
    \item \esp{trae $\rightarrow$ traer ; confunden $\rightarrow$ confundir ; llegan $\rightarrow$ llegar ; exige $\rightarrow$ exigir ; fomentar} est déjà à l'infinitif.
\end{enumerate}

\textbf{4. Expression personnelle (modèle, environ 50 mots) :}
\esp{En mi opinión, la acción más urgente es prohibir los plásticos de un solo uso y organizar limpiezas regulares de las playas de Kribi y Limbe. Además, es urgente que el gobierno eduque a los ciudadanos sobre el reciclaje. Si cada barrio separa su basura, el litoral camerunés volverá a estar limpio.} (50 mots)

\textbf{Points de vigilance :} citez le texte entre guillemets pour justifier ; respectez la limite de mots ; attention à \esp{es urgente que + subjuntivo} (\esp{eduque, actúe}) ; ne traduisez pas mot à mot « il n'y a pas de planète B » : le slogan espagnol est figé : \esp{¡No hay planeta B!}.
\end{corrige}

\begin{corrige}[Exercice 10 — Thème et Version : Traduction spécialisée]
\textbf{A. Version (Français $\rightarrow$ Espagnol) :}
\begin{enumerate}
    \item \esp{La deforestación en la Amazonia acelera el calentamiento global.}
    \item \esp{Debemos prohibir los plásticos de un solo uso desde ahora.} (ou \esp{a partir de ahora}.)
    \item \esp{Las energías renovables son la solución para el futuro del planeta.}
    \item \esp{¡No compren productos con demasiados envases!} — Impératif négatif \emph{ustedes} : \esp{no compren}.
    \item \esp{Los niños aprenden a reciclar en la escuela primaria.}
\end{enumerate}
\textbf{Points de vigilance :} « plastiques à usage unique » = \esp{plásticos de un solo uso} (expression figée) ; « emballages » = \esp{envases} ou \esp{embalajes} ; « dès maintenant » = \esp{desde ya / a partir de ahora}.

\textbf{B. Thème (Espagnol $\rightarrow$ Français) :}
« La sécheresse prolongée dans le nord du Cameroun menace la sécurité alimentaire. Les agriculteurs ne peuvent cultiver ni maïs ni mil. Les puits s'assèchent et le bétail meurt. Il est vital de construire des barrages et d'utiliser l'irrigation au goutte-à-goutte. De plus, le reboisement avec des espèces autochtones freine l'avancée du désert. Chaque citoyen peut agir : planter un arbre, ne pas brûler les ordures, dénoncer la coupe illégale. Le changement commence dans ton quartier. »

\textbf{Points de vigilance :} \esp{pozos} = puits (et non « puits de pétrole » ici) ; \esp{ganado} = bétail ; \esp{riego por goteo} = irrigation au goutte-à-goutte ; \esp{tala ilegal} = coupe/abattage illégal ; \esp{especies autóctonas} = espèces autochtones ; \esp{barrio} = quartier (faux ami : pas « barrière »).
\end{corrige}

\begin{corrige}[Exercice 11 — Sujet de rédaction type Bac : Production écrite]
\textbf{Barème indicatif (sur 20) :} respect de la structure et des contraintes 4 pts ; richesse du lexique imposé (au moins 5 termes) 4 pts ; correction grammaticale (3 impératifs minimum, connecteurs, subjonctif) 6 pts ; ancrage camerounais et originalité du slogan 3 pts ; présentation et longueur 3 pts.

\textbf{Proposition de rédaction modèle :}

\esp{¿CÓMO PROTEGER EL MEDIO AMBIENTE EN TU CIUDAD?}

\esp{¡Douala y Yaoundé se ahogan! La contaminación del aire, de los ríos como el Wouri y de nuestros mercados amenaza nuestra salud. La deforestación avanza, la sequía castiga el norte y la capa de ozono sufre. Por tanto, actuar es un deber de todos.}

\esp{¡Tú puedes cambiar las cosas! ¡Recicla el plástico y el papel en tu barrio! ¡No tires basura en los caniveaux! ¡Apaga el generador y usa energías renovables como el sol! ¡Planta un árbol frente a tu casa! ¡No quemes los residuos! Además, prefiere el transporte público o la bicicleta para reducir tu huella de carbono. Es urgente que cada ciudadano participe en el reciclaje y que las empresas respeten la naturaleza.}

\esp{¡Protege tu ciudad hoy: no hay planeta B!}

(Environ 130 mots ; lexique imposé : \esp{contaminación, deforestación, sequía, capa de ozono, reciclaje, energías renovables, huella de carbono} ; impératifs : \esp{recicla, no tires, apaga, usa, planta, no quemes} ; connecteurs : \esp{por tanto, además, es urgente que + subjuntivo}.)

\textbf{Points de vigilance :}
\begin{itemize}
    \item Vérifiez chaque impératif : \esp{apaga} (\esp{apagar}), \esp{no quemes} (subjonctif de \esp{quemar}), \esp{planta}.
    \item \esp{Es urgente que} exige le subjonctif : \esp{que cada ciudadano participe}, et non \esp{participa}.
    \item Accents obligatoires : \esp{contaminación, sequía, energías, árbol, bicicleta}.
    \item Évitez le calque « environnement » : on dit \esp{medio ambiente} ; « sensibiliser » = \esp{sensibilizar} ou \esp{concienciar}.
    \item Le slogan final doit être court, rythmé et mémorable : c'est la signature de l'affiche.
\end{itemize}
\end{corrige}

\newpage

\coursec{Fiche bilan}

\courssub{Carte mentale de la leçon}
\begin{popfigure}
\begin{tikzpicture}[
    mindmap,
    grow cyclic,
    every node/.style={concept, execute at begin node=\hskip0pt},
    concept color=popblue,
    level 1/.append style={level distance=4.5cm, sibling angle=360/6},
    level 2/.append style={level distance=3cm, sibling angle=360/4},
    text=white,
    font=\small\bfseries
]
\node[concept, scale=1.1, align=center] {E09 : Contaminación \\ y preservación}
    child[concept color=poporange] { node[concept, align=center] {Texte support\\ \esp{« Un planeta}\\\esp{en peligro »}} }
    child[concept color=popgreen] { node[concept, align=center] {Lexique clé\\ \esp{contaminación,}\\\esp{calentamiento global,}\\\esp{deforestación,\\ energías renovables}} }
    child[concept color=poppurple] { node[concept, align=center] {Commentaire\\ Structure : idée directrice,\\ arguments, connecteurs} }
    child[concept color=poppink] { node[concept, align=center] {Grammaire\\ Impératif (affirmatif / négatif),\\ verbes d'action, \esp{por / para}} }
    child[concept color=popgold] { node[concept, align=center] {Affiche de sensibilisation\\ Slogan, image, message clair,\\ appel à l'action} }
    child[concept color=popdark] { node[concept, align=center] {Production\\ Rédaction d'une affiche,\\ exposé oral, débat} };
\end{tikzpicture}
\legende{Synthèse visuelle de la leçon E09 — Module 2 : Environnement, santé et bien-être}
\end{popfigure}

\courssub{Bilan des connaissances et savoir-faire}

\textbf{1. Lexique thématique essentiel (Espagnol $\leftrightarrow$ Français)}

\begin{center}
\begin{tabularx}{\linewidth}{|X|X|X|}
\hline
\rowcolor{popblueL} \textbf{Nom / Verbe} & \textbf{Traduction} & \textbf{Collocations / Exemples} \\
\hline
\esp{la contaminación} & la pollution & \esp{la contaminación del aire / del agua / del suelo} \\
\hline
\esp{el calentamiento global} & le réchauffement climatique & \esp{luchar contra el calentamiento global} \\
\hline
\esp{la deforestación} & la déforestation & \esp{la deforestación masiva en la Amazonía} \\
\hline
\esp{la sequía} & la sécheresse & \esp{una sequía prolongada ; la sequía causa hambruna} \\
\hline
\esp{la capa de ozono} & la couche d'ozone & \esp{el agujero de la capa de ozono ; proteger la capa de ozono} \\
\hline
\esp{las energías renovables} & les énergies renouvelables & \esp{la energía solar, eólica, hidráulica} \\
\hline
\esp{el efecto invernadero} & l'effet de serre & \esp{los gases de efecto invernadero} \\
\hline
\esp{la huella de carbono} & l'empreinte carbone & \esp{reducir nuestra huella de carbono} \\
\hline
\esp{reciclar / el reciclaje} & recycler / le recyclage & \esp{reciclar el plástico ; el reciclaje selectivo} \\
\hline
\esp{conservar / preservar} & conserver / préserver & \esp{conservar la biodiversidad ; preservar los recursos naturales} \\
\hline
\esp{la sostenibilidad} & la durabilité & \esp{el desarrollo sostenible ; metas de sostenibilidad} \\
\hline
\esp{concienciar / tomar conciencia} & sensibiliser / prendre conscience & \esp{concienciar a la población ; tomar conciencia del problema} \\
\hline
\esp{el medio ambiente} & l'environnement & \esp{proteger el medio ambiente ; cuidar el medio ambiente} \\
\hline
\esp{los residuos / la basura} & les déchets / les ordures & \esp{tirar la basura ; separar los residuos} \\
\hline
\esp{reutilizar / reducir} & réutiliser / réduire & \esp{reducir, reutilizar y reciclar (las tres erres)} \\
\hline
\end{tabularx}
\end{center}

\begin{attention}[Faux amis et pièges lexicaux]
\begin{itemize}
    \item \esp{el medio ambiente} signifie « l'environnement » ; ne dites jamais \esp{el ambiente} seul pour « l'environnement » (\esp{el ambiente} = « l'ambiance »).
    \item \esp{concienciar} = « sensibiliser » ; le verbe \esp{realizar} ne signifie pas « réaliser » au sens de « se rendre compte » : on dit \esp{darse cuenta de}.
    \item \esp{la sequía} (sécheresse) s'écrit avec un accent sur le \emph{í} ; ne pas confondre avec \esp{la sequedad} (sécheresse de la peau, d'un lieu).
    \item « Protéger » se dit \esp{proteger} (verbe en \emph{-ger} : \esp{yo protejo}, jamais \esp{protego}).
\end{itemize}
\end{attention}

\vspace{0.3cm}
\textbf{2. Règles de grammaire à connaître par cœur}

\begin{center}
\begin{tabularx}{\linewidth}{|X|X|}
\hline
\rowcolor{poporangeL} \textbf{Notion} & \textbf{Règle synthétique et exemples} \\
\hline
\textbf{Impératif affirmatif} (\esp{tú / usted / nosotros / vosotros / ustedes}) &
\textbf{Forme :} \esp{tú} = 3\textsuperscript{e} pers. sg. de l'indicatif présent (\esp{habla, come, vive}). \esp{vosotros} : infinitif sans \emph{-r} + \emph{-d} (\esp{hablad, comed, vivid}). \esp{usted, ustedes, nosotros} : subjonctif présent (\esp{¡hable! ¡hablen! ¡hablemos!}). \par
\textbf{Irréguliers (\esp{tú}) :} \esp{di, haz, ve, pon, sal, sé, ten, ven} (\esp{decir, hacer, ir, poner, salir, ser, tener, venir}). \par
\textbf{Pronoms :} enclitiques (accrochés au verbe) : \esp{¡Cómpralo! ¡Dímelo! ¡Reciclémoslo!} (accent maintenu sur la voyelle tonique). \\
\hline
\textbf{Impératif négatif} (\esp{no} + impératif) &
\textbf{Toutes les personnes = subjonctif présent :} \esp{no hables, no hable, no hablemos, no habléis, no hablen}. \par
\textbf{Pronoms :} proclitiques (placés devant le verbe) : \esp{No lo compres. No me lo digas. No lo tires al río.} \\
\hline
\textbf{Infinitif comme impératif impersonnel} &
Usage typique des affiches, notices et consignes générales : \esp{No fumar. Prohibido tirar basura. Preservar el agua. Reciclar siempre.} \\
\hline
\textbf{\esp{Por} / \esp{Para} (rappel ciblé environnement)} &
\textbf{\esp{Por} :} cause, motif, échange, durée, moyen : \esp{Lo hago \textbf{por} el planeta. \textbf{Por} culpa de la contaminación, el río está sucio.} \par
\textbf{\esp{Para} :} but, destination, destinataire, échéance : \esp{Esta bolsa es \textbf{para} reciclar. Trabajamos \textbf{para} un futuro mejor. El cartel estará listo \textbf{para} el lunes.} \\
\hline
\textbf{\esp{Se} impersonnel / passif réfléchi} &
Très fréquent dans les affiches et les consignes officielles : \esp{\textbf{Se prohíbe} tirar basura. \textbf{Se recomienda} usar el transporte público. \textbf{Se reciclan} los plásticos.} \\
\hline
\end{tabularx}
\end{center}

\begin{propriete}[Tableau récapitulatif de l'impératif : \esp{reciclar}]
\begin{center}
\begin{tabular}{|l|l|l|}
\hline
\rowcolor{popgreenL} \textbf{Personne} & \textbf{Affirmatif} & \textbf{Négatif} \\
\hline
\esp{tú} & \esp{¡recicla!} & \esp{¡no recicles!} \\
\hline
\esp{usted} & \esp{¡recicle!} & \esp{¡no recicle!} \\
\hline
\esp{nosotros/-as} & \esp{¡reciclemos!} & \esp{¡no reciclemos!} \\
\hline
\esp{vosotros/-as} & \esp{¡reciclad!} & \esp{¡no recicléis!} \\
\hline
\esp{ustedes} & \esp{¡reciclen!} & \esp{¡no reciclen!} \\
\hline
\end{tabular}
\end{center}
\end{propriete}

\vspace{0.3cm}
\textbf{3. Méthodes express}

\begin{methode}[Commentaire de texte en trois étapes]
\begin{enumerate}
    \item \textbf{Introduction :} présenter la nature du texte (article, rapport, discours), sa source, sa date, le thème général (la pollution) et l'idée directrice (la thèse de l'auteur).
    \item \textbf{Développement (deux ou trois parties) :} suivre la structure du texte. Pour chaque partie : citer entre guillemets \esp{« ... »}, puis analyser (vocabulaire du champ lexical de l'environnement, connecteurs \esp{por tanto, sin embargo, además, en consecuencia}, temps verbaux). Relier le texte au contexte (Cameroun, Guinée équatoriale, monde hispanophone).
    \item \textbf{Conclusion :} résumer les arguments, dégager la portée du texte et proposer une ouverture personnelle (engagement citoyen, solutions locales).
\end{enumerate}
\end{methode}

\begin{methode}[Création d'une affiche de sensibilisation : la checklist]
\begin{enumerate}
    \item \textbf{Accroche (slogan) :} courte, percutante, à l'impératif ou à l'infinitif : \esp{¡Salva tu río! / ¡No al plástico!}
    \item \textbf{Visuel :} une image forte, des couleurs contrastées (vert et bleu pour la nature, rouge et noir pour le danger).
    \item \textbf{Message informatif :} un chiffre clé ou un fait court : \esp{Gran parte de la basura marina es plástico.}
    \item \textbf{Appel à l'action :} verbe à l'impératif ou à l'infinitif : \esp{Reduce, reutiliza, recicla. / ¡Planta un árbol hoy!}
    \item \textbf{Signature ou logo :} ONG, école, ministère (ex. \esp{MINESEC}, club d'espagnol du lycée).
\end{enumerate}
\end{methode}

\begin{attention}[Traduction thème (français $\to$ espagnol) : vigilance maximale]
\begin{itemize}
    \item \textbf{\esp{Ser} / \esp{estar} :} \esp{El río \textbf{está} contaminado} (état résultant d'une action) mais \esp{La contaminación \textbf{es} un problema grave} (caractérisation).
    \item \textbf{\esp{Por} / \esp{para} :} « pour la planète » (cause) = \esp{\textbf{por} el planeta} ; « pour protéger la nature » (but) = \esp{\textbf{para} proteger la naturaleza}.
    \item \textbf{Subjonctif} après les expressions de volonté ou de nécessité : \esp{Es urgente que \textbf{reduzcamos} los residuos.} « Il est urgent que nous réduisions les déchets. »
    \item « Il faut » + infinitif = \esp{Hay que} + infinitif : \esp{Hay que cuidar el agua.}
\end{itemize}
\end{attention}

\courssub{Checklist avant le Bac}
\begin{aretenir}[Checklist E09 : environnement et affiche]
\begin{itemize}
    \item $\square$ Je sais définir en espagnol \esp{calentamiento global, deforestación, capa de ozono, energías renovables}.
    \item $\square$ Je conjugue sans erreur l'\textbf{impératif affirmatif} (formes régulières + les huit irréguliers \esp{di, haz, ve, pon, sal, sé, ten, ven}) et l'\textbf{impératif négatif} (subjonctif présent).
    \item $\square$ Je place correctement les pronoms : \textbf{enclitiques} à l'affirmatif (\esp{¡cuídalo!}), \textbf{proclitiques} au négatif (\esp{no lo cuides}).
    \item $\square$ Je distingue \textbf{\esp{por}} (cause : \esp{por la contaminación}) et \textbf{\esp{para}} (but : \esp{para proteger la naturaleza}).
    \item $\square$ J'utilise l'\textbf{infinitif comme impératif impersonnel} pour les consignes générales (\esp{No tirar basura. Reciclar.}).
    \item $\square$ Je maîtrise le \textbf{\esp{se} impersonnel / passif réfléchi} pour l'affichage (\esp{Se prohíbe, se recomienda, se reciclan}).
    \item $\square$ Je structure un \textbf{commentaire de texte} : introduction (présentation + thèse), développement (arguments + citations analysées), conclusion (bilan + ouverture).
    \item $\square$ Je repère les \textbf{connecteurs logiques} (\esp{por tanto, en consecuencia, sin embargo, por el contrario, además}).
    \item $\square$ Je sais concevoir une \textbf{affiche} : slogan (impératif) + visuel + donnée clé + appel à l'action + signature.
    \item $\square$ Je peux présenter mon affiche à l'oral en deux minutes : décrire l'image, expliquer le slogan, justifier le message.
    \item $\square$ Je connais trois actions concrètes pour le Cameroun : \esp{plantar árboles en el monte Camerún, limpiar las orillas del río Wouri, usar botellas reutilizables}.
\end{itemize}
\end{aretenir}

\findecours

\end{document}
